% concatenated from root.tex, Packages_NewCommands.tex, 0_Abstract.tex, 1_Introduction.tex, 2_Related_Work.tex, 3_Problem_Statement.tex, 4_Preliminaries.tex, 5_Methods.tex, 6_Experiments.tex, 7_Conclusion_Future_Work.tex
\documentclass[letterpaper, 10 pt, conference]{ieeeconf}  

% \documentclass[12pt]{article}
% \usepackage[margin=1in]{geometry}
% \usepackage{amsthm}

% \usepackage{graphicx}
\usepackage{booktabs}
\usepackage{adjustbox}
%\usepackage[margin=1in, top=1in, bottom=1in]{geometry}

% \usepackage[top=20.2mm,bottom=18.0mm,left=16.9mm,right=16.9mm]{geometry}

\usepackage[top=20.2mm,bottom=18mm,left=16.9mm,right=20.2mm]{geometry}

%\usepackage{todonotes}
\usepackage{wrapfig}
% \usepackage[labelfont=bf,figurename=Fig.,font=small]{caption}
% \usepackage{subcaption}
\usepackage{lipsum}
%\usepackage{amsthm}
\usepackage{amssymb}
\usepackage{amsmath}
\usepackage{siunitx}
\usepackage{mathtools}
% \usepackage{xcolor}
% \usepackage{color}
\usepackage{hyperref}
\let\labelindent\relax
\usepackage[inline]{enumitem}
\usepackage[ruled,vlined,linesnumbered]{algorithm2e}
\usepackage[noend]{algpseudocode}
\usepackage{ifthen}
\usepackage{dsfont}
% \usepackage[backend=biber,style=numeric-comp,sorting=none,maxbibnames=99]{biblatex}
\usepackage{balance}
\usepackage{dirtytalk}
% \usepackage{natbib}
\usepackage{afterpage}
\usepackage{changepage}
\usepackage{bm}
\usepackage{cite}

\usepackage{times}

\usepackage{svg}



\allowdisplaybreaks


\newcommand{\A}{\mathcal{A}}
\newcommand{\Avoid}{\mathbb{A}}
\newcommand{\B}{\mathcal{B}}
% \newcommand{\cost}{\text{cost}}
\newcommand{\C}{\mathcal{C}}
\newcommand{\D}{\mathcal{D}}
\newcommand{\diag}{\text{diag}}
\newcommand{\blockDiag}{\text{blkdiag}}
\newcommand{\E}{\mathbb{E}}
\newcommand{\Exp}{\text{Exp}}
\newcommand{\F}{\mathcal{F}}
\newcommand{\G}{\mathcal{G}}
\newcommand{\IMU}{\text{IMU}}
\newcommand{\K}{\mathcal{K}}
\newcommand{\KKT}{\text{KKT}}
\newcommand{\Lin}{\mathcal{L}}
\newcommand{\Log}{\text{Log}}
\newcommand{\M}{\mathcal{M}}
\newcommand{\N}{\mathbb{N}}
\newcommand{\Par}{\mathcal{P}}
\newcommand{\Prob}{\mathbb{P}}
\newcommand{\Q}{\mathbb{Q}}
\newcommand{\ra}{\rightarrow}
\newcommand{\R}{\mathbb{R}}
\newcommand{\Ra}{\Rightarrow}
\newcommand{\Riem}{\mathscr{R}}
\newcommand{\Safe}{\mathbb{S}}
\newcommand{\sym}{\mathbb{S}}
\newcommand{\urk}{{\urcorner{k}}}
\newcommand{\urs}{{\urcorner{s}}}
\newcommand{\U}{\mathcal{U}}
\newcommand{\V}{\mathcal{V}}
\newcommand{\X}{\mathcal{X}}
\newcommand{\Y}{\mathcal{Y}}
\newcommand{\Z}{\mathcal{Z}}
\newcommand{\upRiemannint}[2]{\overline{\int_{#1}^{#2}}}
\newcommand{\loRiemannint}[2]{\underline{\int_{#1}^{#2}}}
\newcommand{\Lim}[1]{\raisebox{0.5ex}{\scalebox{0.8}{$\displaystyle \lim_{#1}\;$}}}
\newcommand{\?}{\stackrel{?}{=}}
\newcommand{\st}{\text{ s.t. }}
\newcommand{\paren}[1]{{({#1})}}


\newcommand{\boldlambda}{\boldsymbol{\lambda}}
\newcommand{\boldnu}{\boldsymbol{\nu}}
\newcommand{\bolddelta}{\boldsymbol{\delta}}
\newcommand{\boldeta}{\boldsymbol{\eta}}

\newcommand{\constraint}{\text{constraint}}


\newcommand{\bolda}{\boldsymbol{a}}
\newcommand{\boldA}{\textbf{A}}
\newcommand{\boldb}{\textbf{b}}
\newcommand{\boldq}{\boldsymbol{q}}
\newcommand{\boldr}{\boldsymbol{r}}
\newcommand{\boldell}{\boldsymbol{\ell}}

\newcommand{\boldf}{\textbf{f}}
\newcommand{\boldg}{\textbf{g}}
\newcommand{\boldh}{\textbf{h}}

\newcommand{\boldx}{\textbf{x}}
\newcommand{\boldu}{\textbf{u}}
\newcommand{\boldw}{\textbf{w}}
\newcommand{\boldy}{\textbf{y}}
\newcommand{\bolde}{\textbf{e}}
\newcommand{\boldv}{\textbf{v}}
\newcommand{\boldz}{\textbf{z}}
\newcommand{\boldzero}{\textbf{0}}

\newcommand{\boldU}{\textbf{U}}
\newcommand{\boldY}{\textbf{Y}}

\newcommand*{\mydprime}{^{\prime\prime}\mkern-1.2mu}
\newcommand*{\mytrprime}{^{\prime\prime\prime}\mkern-1.2mu}
\newcommand{\numpy}{{\tt numpy}}    % tt font for numpy




\newcommand{\obj}{J}
\newcommand{\objSq}{J_{\mathrm{sq}}}
\newcommand{\objSqCons}{J_{\mathrm{sqc}}}
\newcommand{\objLQR}{J_{\mathrm{LQR}}}
\newcommand{\param}{\theta}
\newcommand{\paramDim}{d}
\newcommand{\paramSet}{\Theta}
\newcommand{\paramSetSq}{\Theta_{\mathrm{sq}}}
\newcommand{\paramSetLQR}{\Theta_{\mathrm{LQR}}}
\newcommand{\paramSetSqc}{\Theta_{\mathrm{sqc}}}
\newcommand{\var}{\xi}
\newcommand{\varMap}{\phi}
\newcommand{\varOpt}{\var^\star}
\newcommand{\varDim}{n}
\newcommand{\varConstraintSet}{\C}
\newcommand{\dyn}{f}
\newcommand{\varOptSet}{S_\var^\star}

\newcommand{\varDual}{\lambda}
\newcommand{\varDualOpt}{\varDual^\star}

\newcommand{\outputs}{y}
\newcommand{\outputsStacked}{\boldsymbol{\outputs}}

\newcommand{\varStacked}{\boldsymbol{\var}}
\newcommand{\paramStacked}{\boldsymbol{\param}}
\newcommand{\ratePGDStacked}{\boldsymbol{\ratePGD}}
\newcommand{\outputStacked}{\boldsymbol{y}}

\newcommand{\varNominal}{\hat\var}
\newcommand{\paramNominal}{\hat\param}
\newcommand{\ratePGDNominal}{\hat\ratePGD}

\newcommand{\varNominalStacked}{\hat{\boldsymbol{\var}}}
\newcommand{\paramNominalStacked}{\hat{\boldsymbol{\param}}}
\newcommand{\ratePGDNominalStacked}{\hat{\boldsymbol{\ratePGD}}}

\newcommand{\varError}{\Delta\var}
\newcommand{\paramError}{\Delta\param}
\newcommand{\ratePGDError}{\Delta\ratePGD}
\newcommand{\ratePGDErrorStacked}{\Delta\ratePGDStacked}

\newcommand{\convergenceRadius}{\beta}

\newcommand{\ballUnit}{\mathbb{B}}



\newcommand{\linearizationError}{r}
\newcommand{\disturbance}{d}
\newcommand{\disturbanceStacked}{\boldsymbol{\disturbance}}

\newcommand{\linearizationErrorSmoothed}{\tilde{\linearizationError}}
\newcommand{\disturbanceSmoothed}{\tilde{\disturbance}}
\newcommand{\disturbanceStackedSmoothed}{\tilde{\disturbanceStacked}}

\newcommand{\regularizerSystemResponse}{H}


\newcommand{\smoothnessConstant}{L}
\newcommand{\strongConvexityConstant}{m}

\newcommand{\ratePGD}{\alpha}
\newcommand{\ratePGDMin}{c_\ratePGD}
\newcommand{\ratePGDMax}{C_\ratePGD}
\newcommand{\convergenceRatePGD}{\gamma}
\newcommand{\dynamicsPGD}{f_{\text{PGD}}}
\newcommand{\dynamicsPGDSq}{f_{\mathrm{PGD, sq}}}
\newcommand{\dynamicsPGDLQR}{f_{\mathrm{PGD, LQR}}}

\newcommand{\feedbackMap}{\pi}
\newcommand{\feedbackGain}{K}
\newcommand{\setFeedbackGains}{S_\feedbackGain}
\newcommand{\feedbackGainStacked}{\K}

\newcommand{\downshift}{\Z}

\newcommand{\systemResponse}{\Phi}
\newcommand{\setSystemResponses}{S_\systemResponse}

\newcommand{\jacobianState}{A}
\newcommand{\jacobianControl}{B}
\newcommand{\jacobianStateStacked}{\A}
\newcommand{\jacobianControlStacked}{\B}
\newcommand{\outputMatrixStacked}{\C}

\newcommand{\jacobianStateSmoothed}{\tilde{\jacobianState}}
\newcommand{\jacobianControlSmoothed}{\tilde{\jacobianControl}}


\newcommand{\traj}{\phi}
\newcommand{\proj}{\text{proj}}
\newcommand{\gradScale}{\eta}
\newcommand{\constraintSet}{{S_{\mathcal{C}}}}

% \newcommand{\iterationHorizon}{{\bar K}}
% \newcommand{\iterationHorizon}{{n_k}}
\newcommand{\iterationHorizon}{N}
\newcommand{\LQRHorizon}{T}
\newcommand{\LQRIterate}{t}

\newcommand{\state}{\zeta}
\newcommand{\stateNominal}{\hat\state}
\newcommand{\stateStacked}{\boldsymbol{\state}}
\newcommand{\stateNominalStacked}{\hat{\boldsymbol{\state}}}
\newcommand{\stateError}{\Delta\state}
\newcommand{\stateErrorStacked}{\Delta\stateStacked}

\newcommand{\error}{e}

\newcommand{\dynamicsFullState}{f}

\newcommand{\dynamicsNoise}{w}
\newcommand{\outputNoise}{v}

\newcommand{\tubeSize}{\tau}
\newcommand{\tubeSizeStacked}{\boldsymbol{\tau}}

\newcommand{\varInitialSet}{S_\var^0}

\newcommand{\hessianBound}{\mu}
\newcommand{\hessianBoundStacked}{\boldsymbol{\mu}}
\newcommand{\hessianBoundSmoothed}{\tilde{\hessianBound}}
\newcommand{\hessianBoundSmoothedStacked}{\tilde{\hessianBoundStacked}}


\newcommand{\derivativeFull}{D}

\newcommand{\varForwardReachableSet}[1]{S_{\var, {#1}}}

\newcommand{\varConvexHull}{S_{\var, \text{CH}}}
\newcommand{\conv}{\text{conv}}

\newcommand{\diam}{\text{diam}}
\newcommand{\dist}{\text{dist}}

\newcommand{\outputMatrix}{C}
\newcommand{\zero}{\textbf{0}}

\newcommand{\ratePGDBoundSlope}{a}
\newcommand{\ratePGDBoundOffset}{b}

\newcommand{\inequalityConstraints}{g}

\newcommand{\vol}{\text{Vol}}

\newcommand{\samplingRadius}{\delta}

\newcommand{\samplingRadiusMax}{C_\samplingRadius}
\newcommand{\samplingVariable}{q}

\newcommand{\dynamicsPGDSmoothed}{\tilde{f}_{\text{PGD}}}
\newcommand{\dynamicsFullStateSmoothed}{{\tilde{\dynamicsFullState}}}

\newcommand{\exampleFunction}{h}
\newcommand{\exampleFunctionSmoothed}{\tilde{\exampleFunction}}
\newcommand{\varExample}{z}

\newcommand{\exampleInputDim}{m}
\newcommand{\exampleOutputDim}{{\overline{m}}}

\newcommand{\lipschitz}{\ell}
\newcommand{\exampleConstraintSet}{S_\exampleFunction}

\newcommand{\unif}{\text{Unif}}

\newcommand{\LQRStateCost}{Q}
\newcommand{\LQRControlCost}{R}
\newcommand{\LQRStateDyn}{A}
\newcommand{\LQRControlDyn}{B}

\newcommand{\SqcQuadScale}{c_1}
\newcommand{\SqcLinScale}{c_2}
\newcommand{\SqcConsVal}{c_3}
\newcommand{\SqcLipschitz}{\ell_{\mathrm{sqc}}}




\newtheorem{remark}{Remark}
\newtheorem{assumption}{Assumption}
\newtheorem{proposition}{Proposition}
\newtheorem{definition}{Definition}
\newtheorem{lemma}{Lemma}
\newtheorem{theorem}{Theorem}
\newtheorem{corollary}{Corollary}
\newtheorem{example}{Example}
\newcommand{\runningexample}[1]{\emph{Running example:}~\textit{#1}}


\newcommand{\frank}[1]{{\textcolor[rgb]{0,0,1}{\textsf{[Frank: {#1}]}}}}

\newcommand{\glen}[1]{{\textcolor[rgb]{1,0,0}{\textsf{[Glen: {#1}]}}}}

\newcommand{\antoine}[1]{{\textcolor[rgb]{1,0,0}{\textsf{[Antoine: {#1}]}}}}

\newcommand{\brendan}[1]{{\textcolor[rgb]{1,0,1}{\textsf{[Brendan: {#1}]}}}}


\newcommand{\sam}[1]{{\textcolor[rgb]{1,0,0}{\textsf{[Sam: {#1}]}}}}

\newcommand{\kyriakos}[1]{{\textcolor[rgb]{1,0,0}{\textsf{[Kyriakos: {#1}]}}}}











\newcommand{\revisionRR}[1]{\textcolor[rgb]{0,0,0}{#1}}

\newcommand{\revision}[1]{\textcolor[rgb]{0,0,0}{#1}}

% \newcommand{\revisionA}[1]{\textcolor[rgb]{1,0,0}{#1}}

\newcommand{\revisionA}[1]{\textcolor[rgb]{0,0,0}{#1}}

\newcommand{\revisionB}[1]{\textcolor[rgb]{0,0,0}{#1}}

\newcommand{\INCOMPLETE}{\textcolor[rgb]{1,0,0}{\textbf{INCOMPLETE}}}
    



% \renewcommand*{\bibfont}{\small}

% \bibliography{references}
% \bibliography{references_abbrv}


\IEEEoverridecommandlockouts
\overrideIEEEmargins


\begin{document}


\title{\LARGE \bf
% Constructing Forward Invariant Sets to Over-Approximate Minimizer Sets of Parameterized Convex Functions
% Over-Approximating Minimizer Sets of Parameterized
% % Convex Functions 
% Constrained Convex Programs
% via 
% Forward Reachability of Projected Gradient Descent Dynamics
Over-Approximating Minimizer Sets of Constrained Convex Programs with Parametric Uncertainty via Reachability Analysis
% Constrained Convex Programs via Reachability Analysis
% System Level Synthesis
% \INCOMPLETE: \textcolor{red}{To edit}
}

\author{
Brendan Gould$^{1\star}$,
Chih-Yuan Chiu$^{1\star}$, Antoine P. Leeman$^{2}$, \\ Kyriakos G. Vamvoudakis$^1$, Samuel Coogan$^1$, Glen Chou$^1$
\thanks{$^\star$Equal contribution. $^{1}$Georgia Institute of Technology, Atlanta, GA, USA. \texttt{\{bgould6\} at gatech dot edu}. $^{2}$ETH Zurich, IDSC, Switzerland. }
% \thanks{Georgia Institute of Technology, Atlanta, GA, USA. $^*$ denotes equal contribution. \texttt{\{cyc, zzhang3097, chou\} at gatech dot edu}}
}
% \author{Chih-Yuan Chiu$^{\star}$, Zhouyu Zhang$^{\star}$, and Glen Chou
% \thanks{Georgia Institute of Technology, Atlanta, GA, USA. $^*$ denotes equal contribution. \texttt{\{cyc, zzhang3097, chou\} at gatech dot edu}}
% }

\maketitle

\thispagestyle{empty}
\pagestyle{empty}

% \thispagestyle{plain}
% \pagestyle{plain}


%%%%%%%%%%%%%%%%%%%%%%%%%%%%%%%%%%%%%%%%%%%%%%%%%%%%%%%%%%%%%%%%%%%%%%%%%%%%%%%%

%%%%%%%%%%%%%%%%%%%%%%%%%%%%%%%%%%%%%%%%%%%%%%%%%%%%%%%%%%%%%%%%%%%%%%%%%%%%%%%%

\begin{abstract}
\looseness-1
% In this paper, we
We
study the set of solutions to a parameterized, strongly convex optimization problem whose cost depends on uncertain, bounded parameters.
We compute
% with uncertain parameters drawn from a bounded set, such as parameters in the cost function, 
% seeking 
a certified outer approximation of the corresponding set of optimizers, 
% In contrast to standard reachable-set analysis, which typically focuses on uncertainty in the system dynamics, our framework addresses uncertainty entering through the optimization problem itself. 
using convergence properties of the projected gradient descent (PGD) algorithm for convex programs. Concretely, by treating the cost parameter as constant but unknown, we interpret the PGD iterates as an uncertain dynamical system and analyze its forward reachable sets. 
Since PGD converges exponentially to the unique optimizer for each fixed parameter, these reachable sets provide 
% anytime 
% \antoine{I think I didn't follow the discussion about this ``anytime''. I think it would be good to explain in a few words what we mean by that somewhere in the paper.}
outer approximations of the optimizer set, with an explicit error bound that decays exponentially with the iteration count. 
We apply system-level synthesis (SLS) on the PGD dynamics to optimize the step-size sequence and obtain reachable-set over-approximations. %Antoine: "tight" would only be for LTV; could be confusing.
Our method outperforms existing baselines in over-approximating, with low conservativeness, the minimizer sets of convex programs with uncertain costs
and high-dimensional decision variables.
% while scaling tractably to high-dimensional problems. 
% \antoine{``high-dimensional regimes'' = problems with large number of decision variables? I am sure if it's clear as is.}
% , and scales tractably to high-dimensional problems. 
% we validate the scalability and accuracy of our approach on constrained 
% % optimal control problems 
% , and illustrate that.
\end{abstract}

% \input{0_Outline}

% \vspace{-3pt}
\section{Introduction}
\label{sec: Introduction}
% \vspace{-2pt}

Across a broad range of controls and robotics contexts, optimization is a critical component of planning and decision-making.
Frequently, the data informing these decisions is not known exactly, but subject to many types of \emph{uncertainty}.
% To ensure these tools are robustly applicable in the real world, it is necessary to understand the effect of \emph{uncertainty} on the solutions to such problems. 
For instance, to operate safely in dense traffic, an autonomous vehicle must anticipate the motion of nearby cars and pedestrians, whose intent
% may be uncertain
can only be partially inferred from observations~\cite{englert2017inverse, Peters2023OnlineandOfflineLearningofPlayerObjectivesfromPartialObservationsinDynamicGames}.
When modeling the motion of such agents as cost minimization~\cite{englert2017inverse, Menner2021ConstrainedInverseOptimalControlWithApplicationtoaHumanManipulationTask}, intent uncertainty induces uncertainty in their planned motion, which could cause unsafe behavior.
% To solve robust control and motion planning problems in multi-agent, interactive contexts, it is often critical to predict the behaviors of strategic agents with unknown intent or objectives. 
% Meanwhile, regulating traffic flows effectively over transportation systems often requires anticipating travelers' selfish routing decisions, even when their routing preferences hinge upon network or traveler attributes are imperfectly estimated \cite{Wu2021ValueofInformationinBayesianRoutingGames, Ferguson2022AvoidingUnintendedConsequencesHowIncentivesAidInformationProvisioninginBayesianCongestionGames}.
Separately, quasistatic \cite{li2026certified} and deformable object \cite{li2025convex} contact dynamics models in robotic manipulation can be expressed as the solution to convex conic programs, where uncertainty in the mass matrix manifests as uncertainty in the cost function.
Towards studying robust decision forecasting,
% in the above contexts,
% To close this gap,
we present a method for over-approximating the set of minimizers of a convex program whose cost is characterized by uncertain parameters.
% \antoine{Now, we have ``set of optimizers'', ``set of minimizers'', and ``set of solutions''. I would try to use a single name throuhgout.}
By robustly bounding the solution set for 
all possible parameter realizations,
% every possible realization of the optimization, 
our work 
% In multi-robot motion planning, our method would enable the over-approximation of agent trajectories under bounded objective uncertainty,
% of agents whose objectives are imprecisely determined with bounded uncertainty,
provides valuable guidance for safe downstream planning.
% actions that may be taken by a strategic agent with uncertain intent. 
% Beyond motion forecasting, our techniques are readily applicable to a wide range of controls and robotics applications that require uncertainty propagation through optimization problems. 
% \antoine{The perception idea looks really cool; but would the reader understand how is that different from the ``uncertain'' intent idea? My guess is that we would want to define a cost function based on some observation, e.g., minimize the distance to a perceived object. I would add a few extra words to strengthen this part of the motivation, as it sounds really cool.}
% \frank{Glen will replace this with an example on contact dynamics.}
% ~\\
% \frank{To Glen: Feel free to add your applications here}
% ~\\

% Concretely, given a constraint set $\constraintSet \subseteq \R^\varDim$ over the variable space $\R^\varDim$, and a set of parameters $\paramSet \subseteq \R^\paramDim$, we consider a parametric family of objective functions $\obj: \R^\varDim \times \paramSet \ra \R$, where $\obj(\cdot, \param)$ is strongly convex for each $\param \in \paramSet$.
% % we encode cost uncertainty 
% % % of the agent’s objective 
% % using a parameter $\param$, which is known to belong in a parameter set $\paramSet \subseteq \R^\paramDim$, but whose true value is unknown.
% % only to the agent.
% % We assume that the agent optimizes over a set $\varConstraintSet \subset \R^\varDim$ of admissible actions to minimize its cost $\obj(\var, \param)$, given by the objective function $\obj: \R^\varDim \times \paramSet \ra \R$.
% % We then associate the agent with a set $\varConstraintSet \subseteq \R^\varDim$ of admissible actions, as well as a parameterized cost function $\obj: \varConstraintSet \times \paramSet \ra \R$ such that, given any agent action $\var \in \varConstraintSet$ and objective parameter $\param \in \paramSet$, the agent incurs the cost $\obj(\var, \param)$. 
% % We assume that the agent will take cost-minimizing actions, and thus ...
% Our aim is to compute or over-approximate the set $\varOptSet$
% % , defined below as the set 
% of all optimizers of $\obj(\cdot, \param)$ as
% % optimal agent actions corresponding to all possible values of 
% the parameter $\param$ varies across $\paramSet$:
% \begin{align} 
% \label{Eqn: Set of All Parameterized Minimizers}
%     \varOptSet := \bigcup_{\param \in \paramSet} \arg\min_{\var \in \constraintSet} \obj(\var, \param).
% \end{align}

\looseness-1The tractability of robustly over-approximating the minimizer set varies significantly depending on the structural properties of the objective function, constraint set, and parameter set. 
% the objective map $\obj$, the constraint set $\constraintSet$, and the parameter set $\paramSet$.
Closed-form expressions for the minimizer set are generally obtainable only under restrictive assumptions. 
% \antoine{This claim needs to come with a citation, or be written more mild.}
% on $\obj$, $\constraintSet$, and $\paramSet$. 
% Even when available, 
% % yield clean, geometric descriptions
% analytic expressions for the minimizer set
% may not admit characterizations in terms of standard primitives, e.g., polytopes, limiting their practical utility.
% for downstream engineering tasks.
% polytopes, circles, or other simple geometric primatives)
Sensitivity analysis techniques
% over the parameter $\param$
can be applied to bound the minimizer's variation
as the parameter changes,
% as $\param$ ranges over $\paramSet$,
but such bounds are often highly conservative (see \ref{subsubsec:baseline}), limiting their utility for downstream tasks.
Related work in set-based reachability and robust optimal control studies how uncertainty propagates through dynamical systems \cite{althoff2008reachability, schafer2025robust}. However, these works seek a single solution minimizing worst-case cost under uncertainty, whereas we aim to over-approximate the full set of minimizers across all parameter realizations.
% Antoine: Feel free to remove althoff2008reachability if we need some space. I like schafer2025robust because it's similar to nonlinear SLS.



% (see our baseline results in Sec. \ref{sec: Experiments}).
% Black-box methods such as constrained Bayesian optimization \cite{Berkenkamp2016BayesianOptimizationWithSafetyConstraints} also tackle uncertain optimization problems, but such approaches typically focus on the sample-efficient identification of a feasible optimum from expensive evaluations, differing from the problem we study.
% % rather than on characterizing the full set of optimizers induced by an admissible parameter set.
%\antoine{I would say those two papers are important to cite (more closely related than BO for instance): }
%\frank{Antoine - Could you add a sentence somewhere in this section to include these references? (Wherever you think would be most appropriate; I know far less about these papers than you, but I agree they are very relevant.)}

% For $\varOptSet$ can be computed in closed form 
% % for specific problem instances,
% % for specific classes of objective functions $\var$ and parameter sets $\param$,
% or over-approximated via sensitivity analysis. 
% the obtained set representations or over-approximations are rarely amenable to a clean geometric description, e.g., a vertex- or halfspace-based polytope description, which may render its direct use in downstream engineering applications more challenging.

% To overcome the limitations of existing methods, we use system-level synthesis (SLS) to compute forward reachable sets of projected gradient descent dynamics (PGD) with low conservatism and apply known convergence properties~\cite{WrightRecht2022OptimizationForDataAnalysis, gokhale2024proximal, shikhman2009constrained} to over-approximate the parameterized minimizer set.
% \antoine{Suggestion so the pipeline is more explicit:
% To overcome the limitations of existing methods, we use system-level synthesis (SLS) to compute finite-time forward reachable-set over-approximations of projected gradient descent (PGD) dynamics with low conservatism. We then leverage known convergence properties~\cite{WrightRecht2022OptimizationForDataAnalysis, gokhale2024proximal, shikhman2009constrained} to convert these finite-time bounds into certified outer approximations of the corresponding minimizer set.}
% \frank{I like your sentences above (a lot!) My only concern would be that the contribution list currently duplicates some of the content of your sentences. Perhaps we can streamline one or the other to reduce repetition?
% }
% $\varOptSet$ 
% which are well-known to converge when applied to a standard class of convex programs \cite[Sec. 7.3]{WrightRecht2022OptimizationForDataAnalysis}.
% Our core insight is to exploit 
% The core insight of our work is to over-approximate
% % the computation of the least conservative over-approximation of 
% $\varOptSet$ using the forward reachable sets of the discrete-time projected gradient descent (PGD) dynamics \cite[Sec. 7.3]{WrightRecht2022OptimizationForDataAnalysis} with respect to $\obj(\cdot, \param)$, with $\param$ treated as a fixed but uncertain value over $\paramSet$.
% It is well-established that, for any fixed parameter $\param \in \paramSet$, the PGD dynamics with respect to a strongly convex and smooth cost $\obj(\cdot, \param)$ over a convex set $\constraintSet$ converges exponentially to the corresponding (unique) global minimizer $\arg\min_{\var \in \constraintSet} \obj(\var, \param)$.
% \frank{Perhaps some brief description here of the method, as is standard .}
\looseness-1To overcome the limitations of existing methods, we use system-level synthesis (SLS) to compute finite-time forward reachable-set over-approximations of projected gradient descent (PGD) dynamics with low conservativeness. We then leverage known convergence properties of the PGD dynamics 
% \cite{WrightRecht2022OptimizationForDataAnalysis, shikhman2009constrained} 
\cite{WrightRecht2022OptimizationForDataAnalysis} 
to convert these reachable sets into certified outer approximations of the corresponding minimizer set.
Concretely, the contribution of this paper is fourfold, as follows:
% Concretely: 
\begin{enumerate}
    \item We prove that, under standard assumptions on the objective, constraint, and parameter sets, the PGD dynamics' (true) forward reachable sets converge exponentially to the minimizer set. 
    We thus recast over-approximating the minimizer set as computing forward reachable sets of the PGD dynamics with a fixed but unknown parameter.
    % 
    \item For \textit{differentiable} PGD dynamics, we over-approximate forward reachable sets of PGD trajectories with the SLS-based reachability framework~\cite{Leeman2025RobustNonlinearOptimalControlviaSLS} for smooth nonlinear dynamics, scaling tractably to high-dimensional systems with low over-approximation error 
    % \cite{Leeman2025RobustNonlinearOptimalControlviaSLS, chen2024robust, leeman2026vision, furieri2019input}.
    \cite{Leeman2025RobustNonlinearOptimalControlviaSLS, leeman2026vision, furieri2019input, fang2026safe}.
    % 
    \item To compute reachable sets for \textit{non-differentiable} PGD dynamics, we apply the sampling-based smoothing method in \cite{Flaxman2005OnlineConvexOptimization} to construct a differentiable approximation of PGD for which the SLS techniques in \cite{Leeman2025RobustNonlinearOptimalControlviaSLS} apply.
    % We demonstrate that, in this setting, the tightness of the constructed outer approximation for 
    % % $\varOptSet$
    % the minimizer set
    % hinges upon the choice of a \textit{smoothing radius} parameter, which governs the curvature of the smooth PGD dynamics proxy and its approximation error with respect to the original PGD dynamics. 
    % 
    \item 
    % Across a range of numerical simulations, 
    In numerical simulations,
    our method outperforms existing sensitivity analysis baselines in constructing over-approximations, with low conservativeness, of minimizer sets of optimization programs with parametric cost uncertainty. 
    Our method scales tractably to optimization programs with up to 64-dimensional variables.
\end{enumerate}



% ~\\
% ~\\
% \frank{INCOMPLETE; list of notation to explain:
% \begin{enumerate}
%     \item Index sets $[n]$ and $[0, n]$;
%     % 
%     \item Identity matrices \say{$I_n$}, with indices omitted when clear from context;
%     % 
%     \item Zero matrices \say{$O_{m \times n}$} and \say{$O_n$}; Zero vectors $\textbf{0}_n$. Indices are omitted when clear from context;
%     % 
%     % \item Diagonal matrix construction of the form $\diag\{x_1, \cdots, x_n\}$;
%     \item Block diagonal matrix $\blockDiag\{M_1, \cdots, M_n\}$
%     % 
%     \item Diameter of a set $S$: $\diam(S)$
%     % 
%     \item Distance from a point $p$ to a set $S$: $\dist(p, S)$ OR distance between two sets $S, S'$: $\dist(S, S') := \min_{x \in S, x' \in S'} \Vert x - x' \Vert_2$.
%     % 
%     \item Partial derivatives, e.g., $\frac{\partial}{\partial \param}$.
%     % 
%     \item Total derivative, e.g., $\derivativeFull \dynamicsFullState$, if applicable. 
%     % 
%     \item Minkowski sum $\oplus$
%     % 
%     \item Unit $p$-norm ball $\ballUnit_p^d$ in $\R^d$, and $M \ballUnit_p^d := \{Mx: x \in \ballUnit_p^d \}$ for any $M \in \R^{d' \times d}$.
%     % 
%     \item $\conv$: Convex Hull.
%     % 
%     \item $\vol$: Volume (if used).
%     % 
%     \item $\unif$: Uniform distribution.
% \end{enumerate}
% }



% ~\\
% ~\\
% \frank{(Mon, March 30, 2025, 5:41 pm ET) In response to Brendan's comment below, I've modified just the sentence above so that it won't be confusing. Of course the entire introduction will be rewritten tonight.}
% \brendan{We are not actually finding a robust forward invariant set anymore. This was our old approach to distinguish between iterates of the PGD system and its solution, but now we use the exponentially decaying bound from convergence properties of PGD, with SLS \emph{reachable sets} of the same system to get our over-approximation.}
% \antoine{Add here or later explicitly that the decision variable is the full trajectory, state, input, over time, i.e., $\in \mathbb{R}^{(N +1)n_x + N \nu }$}
% \frank{If we keep the discussion as is, I'd agree with you that further specificity for the variables we use is necessary. 
% But I'm actually thinking of modifying the introduction so that it is less technical, i.e., uses more natural language rather than math, which I think would be more standard. The technical details and definitions can then come later, say, in the next section. Let me know what you think.
% }
% \antoine{I do like the equation above. It's just that computing a reachable set on full trajectories is not a standard thing to do, so we should emphasize it's not a classic reachability analysis here.}


% \section{Related Work}
% \label{sec: Related Work}

\looseness-1Our work is conceptually related to the field of robust optimization, which studies decision-making under worst-case realizations of uncertain data \cite{bertsimas2004price, ben1998robust, bertsimas2011theory}.
However, our aim differs in that
% from the focus of most robust optimization formulations;
instead of identifying a single robust decision, we seek to characterize the set of decisions that are optimal for some admissible parameter value. 
In this sense, prior work on 
% parametric optimization
multi-parametric programming
and solution-map sensitivity analysis
% via variational inequalities, and multi-parametric programming,
are more closely aligned with our objective, since such works study how minimizers change as problem data vary \cite{bonnans2013perturbation, dontchev2009implicit, shapiro2005sensitivity, jones2007lexicographic}. 
Our contribution departs from this literature by focusing on a global over-approximation of the full minimizer set induced by parameter uncertainty, explicitly leveraging reachability techniques from control theory~\cite{bravo2005computation, immrax}.
% This set-valued viewpoint is especially relevant for
% % in settings such as 
% % games and
% motion planning with uncertain objectives, where one is often interested in recovering the set of plausible optimal actions of another agent rather than a single nominal prediction \cite{zhan2025robustly}.

\textbf{Notation}:
For $m, n \in \mathbb{Z}^+$, we define $[n] := \{1, \cdots, n\}$ and $[0, n] := \{0\} \cup [n]$, and denote the $n \times n$ identity matrix, $n \times n$ zero matrix, $m \times n$ zero matrix, and the zero vector in $\R^n$ by $I_{n \times n}$, $O_n$, $O_{m \times n}$, and $\zero_n$ respectively, with indices omitted when clear from context.
% \antoine{In \eqref{Eqn: Stacking State Jacobians}, I think you use $O_{n}$.}
We denote the block-diagonal matrix constructed from given matrix sub-blocks $M_1, \cdots, M_n$ by $\blockDiag\{M_1, \cdots, M_n\}$; when $M_1, \cdots, M_n$ are all scalar, we write $\diag\{\cdot\}$ instead of $\blockDiag\{\cdot\}$.
Given $p \in [1, \infty]$ and $M \in \R^{n \times d}$, and we denote the unit $p$-norm ball in $\R^d$ by $\ballUnit_p^d$, and we define $M \ballUnit_p^d := \{Mx: x \in \ballUnit_p^d \}$.
Given a set $S \subseteq \R^n$, 
% \antoine{$\subseteq$?} 
we define the \textit{diameter} of $S$ by $\diam(S) := \sup_{s, s' \in S} \Vert s - s'\Vert_2$ and the \textit{convex hull} of $S$ by $\conv(S)$.
Given sets $S, S' \subseteq \R^n$, we define the \textit{Minkowski sum} of $S$ and $S'$ by $S \oplus S' := \{s + s': s \in S, s' \in S' \}$, and the Hausdorff distance between $S$ and $S'$ by $\dist(S, S') := \inf \{ \epsilon > 0: S \subseteq S' \oplus \epsilon I_n \ballUnit_2^n, S' \subseteq S \oplus \epsilon I_n \ballUnit_2^n \}$.
% \antoine{Throughout: consistency of “Assumption 1” and “Assumption 1.” + “Cor. 1” and “Corr. 1.” + “overapproximate” and “over-approximation”}

% \begin{itemize}
%     \item Robust optimization under parameter uncertainty studies how to make decisions that remain effective under worst-case uncertainty \cite{bertsimas2004price, ben1998robust, bertsimas2011theory}. While this is the closest classical framework, its goal is fundamentally different from ours: rather than seeking a single robust decision, we aim to characterize the set of solutions that are optimal for some admissible parameter realization.

% \item Parametric optimization and solution-map sensitivity, including related work on variational inequalities, is more closely aligned with our objective since it studies how optimal solutions vary with perturbations in the problem data \cite{bonnans2013perturbation, dontchev2009implicit, shapiro2005sensitivity}. Closely related, work on multi-parametric programming seeks to characterize solution behavior over a range of parameters through explicit solution maps or parameter-space partitions \cite{jones2007lexicographic}. Our setting builds on these viewpoints, but focuses instead on computing a global over-approximation of the set of all possible minimizers induced by a parameter set.

%     \item To analyze this set, we adopt the lens of optimization dynamics, such as gradient flow and projected gradient methods for constrained problems \cite{WrightRecht2022OptimizationForDataAnalysis, gokhale2024proximal,shikhman2009constrained}. This provides a dynamical characterization of minimizers, allowing the parameter-dependent optimizer set to be studied through the behavior of uncertain optimization trajectories.

%     \item This perspective connects our problem naturally to reachable-set analysis for uncertain nonlinear systems \cite{bravo2005computation}. In our setting, uncertainty enters through the optimization problem rather than the system dynamics directly, and the set of possible optimizers is over-approximated by reachable sets of the uncertain projected optimization dynamics, together with explicit bounds from the convergence of projected gradient descent.

%     \item To compute such sets in practice, we build on System Level Synthesis (SLS) methods for robust reachable-set construction \cite{Leeman2025RobustNonlinearOptimalControlviaSLS, chen2024robust, leeman2026vision, furieri2019input}. These tools provide the computational machinery that makes our set-based characterization tractable for the optimization dynamics considered here. Uncontrolled reachability analysis: \cite{althoff2008reachability}.

%     \item Finally, our work is motivated by problems involving intent uncertainty in games and motion planning under uncertain objectives, where one is often interested in the set of plausible optimal actions of another agent rather than a single nominal prediction \cite{zhan2025robustly}. In this sense, our framework provides a bridge between parameter uncertainty in optimization and set-valued prediction in strategic settings.
% \end{itemize}

% ~\\
% \frank{INCOMPLETE}
% ~\\



% \clearpage
\section{Problem Statement}
\label{sec: Problem Statement}


% \section{Problem Statement and Preliminaries}
% \label{sec: Problem Statement and Preliminaries}

% ~\\
% \frank{INCOMPLETE; Exposition to fill in, if there is time.}
% ~\\

% \subsection{Problem Statement}
% \label{subsec: Problem Statement}

% Let $\obj: \R^\varDim \times \R^\paramDim \ra \R$ denote the non-ego objective function, which is characterized by a parameter 
% $\param$,
% % $\param \in \paramSet \subseteq \R^\paramDim$,
% contained in a parameter set $\paramSet \subseteq \R^\paramDim$ whose true value $\param^\star$ is known to the non-ego agent but unknown to the ego agent.
% The non-ego agent aims to minimize $\obj(\cdot, \param^\star)$ over a constraint set $\constraintSet$ which is known to both the non-ego and ego agents.
% \antoine{Reference Eq.(1) here and the exact math definition of its components.}
Consider an objective $\obj: \R^\varDim \times \R^\paramDim \ra \R$ mapping a variable $\var \in \constraintSet \subseteq \R^\varDim$ and a parameter $\param \in \paramSet \subset \R^\paramDim$ to a real number.
We impose the following regularity assumptions.
% \antoine{I think that if $\paramSet$ is a set, it should be $\theta \in \paramSet \subset \R^\paramDim$. Not to be picky, but the reader really need to understand here that $\Theta$ is a set and not a vector.}
\begin{assumption}
\label{Assump: Objective Properties}
All partial second derivatives of $\obj(\var, \param)$ exist and are Lipschitz continuous. 
There exist scalars
% \antoine{scalars} 
$\strongConvexityConstant, \smoothnessConstant > 0$, with $\strongConvexityConstant \leq \smoothnessConstant$, such that $\obj(\cdot, \param)$ is $\smoothnessConstant$-smooth and $\strongConvexityConstant$-strongly convex, 
% over $\constraintSet$ 
% for each $\param \in \paramSet$, s
i.e., for each $\param \in \paramSet$ 
% \antoine{for each $\param \in \paramSet$, i.e., for any $\param \in \paramSet$; this reads weird. I think both mean the same no?}
and $\var, \var' \in \R^\varDim$:
\begin{align} \nonumber
    \obj(\var', \param) &\geq \obj(\var, \param) + \nabla \obj(\var, \param)^\top (\var' - \var) + \textstyle\frac{1}{2}\strongConvexityConstant \Vert \var' - \var \Vert_2^2, \\ \nonumber
    \obj(\var', \param) &\leq \obj(\var, \param) + \nabla \obj(\var, \param)^\top (\var' - \var) + \textstyle\frac{1}{2}\smoothnessConstant \Vert \var' - \var \Vert_2^2.
\end{align}
\end{assumption}
% \antoine{In the game context, are those assumptions satisfied? What type of convex problem would be a good application for this?}
% \frank{I'm thinking specifically of the following scenario which I intend to describe in the introduction: 
% Consider a 2-agent game in which one agent (called the \say{ego}) observes another (the \say{non-ego}) in its vicinity.
% The ego agent would like to know where the non-ego agent is going, i.e., would like to compute the non-ego agent's forward reachable set, but does not perfectly know the non-ego agent's intent $\param$ (which parameterizes the non-ego agent's cost $\obj(\cdot, \param)$), e.g., its goal position, the extent to which it is risk-averse against or agnostic to incurring collisions, etc. (The non-ego agent is assumed to perfectly know their intent). Rather, the ego agent only knows that the non-ego agent's true intent parameter $\param$ fall within a given set $\paramSet$ of possible intent parameter values.
% Thus, in this paper, we adopt the ego agent's perspective, and try to compute the set of all possible non-ego behaviors, as corresponding to each of these intent parameters.
% For any fixed non-ego intent parameter $\param$, I think it is reasonable to assume that $\obj(\cdot, \param)$ is strongly convex; this could just be a LQR problem whose problem data the non-ego agent perfectly knows and is capable of solving accurately.
% Let me know if this framework makes sense.
% }

\begin{assumption}
\label{Assump: Constraint Set is Convex and Closed}
$\constraintSet$ is convex and closed in $\R^\varDim$.
\end{assumption}

\begin{assumption}
\label{Assump: Parameter Set is Convex and Compact}
$\paramSet$ is convex and compact in $\R^\paramDim$. 
% Let $\diam(\paramSet) := \sup_{\param, \param'} \Vert \param - \param' \Vert_2$ denote the \textit{diameter} of $\paramSet$.
\end{assumption}
% \antoine{In the Notation, it's $\diam$, or is it something different?}

Under Assumption \ref{Assump: Objective Properties}, for each $\param \in \paramSet$, $\obj(\cdot, \param)$ has a unique minimizer over $\constraintSet$, denoted $\varOpt(\param)$ below \cite[Thm. 2.8]{WrightRecht2022OptimizationForDataAnalysis}.

\textbf{Problem Statement}:
% We aim to construct an over-approximation for the set $\varOptSet$, defined by:
We aim to construct an outer approximation of the minimizer set $\varOptSet$, defined by:
\begin{align}
\label{Eqn: Set of All Parameterized Minimizers}
    \varOptSet &:= \textstyle\bigcup_{\param \in \paramSet} \arg\min_\var \obj(\var, \param) = \big\{\varOpt(\param): \param \in \paramSet \big\}.
\end{align}

% \eqref{Eqn: Set of All Parameterized Minimizers}, which, under Assumption \ref{Assump: Objective Properties}, can be written as $\varOptSet 
% % := \bigcup_{\param \in \paramSet} \arg\min_\var \obj(\var, \param)
% = \big\{\varOpt(\param): \param \in \paramSet \big\}$.

% The central objective of the ego agent is to 
% % compute or
% over-approximate the set $\varOptSet$ of non-ego actions that are optimal with respect to \textit{some} parameter value $\param$ in $\paramSet$ over the known constraint set $\constraintSet$, defined by:
% \begin{align}
% % \label{Eqn: Set of All Parameterized Minimizers}
%     \varOptSet &:= \bigcup_{\param \in \paramSet} \arg\min_\var \obj(\var, \param) = \big\{\varOpt(\param): \param \in \paramSet \big\}.
% \end{align}
% We will henceforth refer to $\varOptSet$ as the \say{target set}.

% Below, we establish the boundedness of $\varOptSet$ under Assumptions \ref{Assump: Objective Properties} and \ref{Assump: Parameter Set is Convex and Compact}.

% \begin{proposition}
% \label{Prop: Minimizer Set is Bounded}
% Under Assumptions \ref{Assump: Objective Properties} and \ref{Assump: Parameter Set is Convex and Compact}, $\varOptSet$ is bounded.
% \end{proposition}

% \begin{proof}
% Fix $\Tilde{\var} \in \constraintSet$ arbitrarily. By Assumption \ref{Assump: Objective Properties}, the map $\param \mapsto \Vert \nabla \obj(\Tilde{\var}, \param) \Vert_2$ is continuous, and thus, by Assumption \ref{Assump: Parameter Set is Convex and Compact}, it attains a finite maximum value over $\paramSet$.
% Thus, for any $\var \in \constraintSet$ satisfying $\Vert \var - \Tilde{\var} \Vert_2 > \frac{2}{m} \max_{\param' \in \paramSet} \Vert \nabla \obj(\Tilde{\var}, \param') \Vert_2$, and any $\param \in \paramSet$, we obtain from the strong convexity of the map $\var \mapsto \obj(\var, \param)$ that:
% \begin{align} \nonumber
%     \obj(\var, \param) &\geq \obj(\Tilde{\var}, \param) + \nabla \obj(\Tilde{\var}, \param)^\top (\var - \Tilde{\var}) + \frac{1}{2}m \Vert \var - \Tilde{\var} \Vert_2^2 \\ \nonumber
%     &\geq \obj(\Tilde{\var}, \param) - \Vert \nabla \obj(\Tilde{\var}, \param) \Vert_2 \Vert \var - \Tilde{\var} \Vert_2 + \frac{1}{2}m \Vert \var - \Tilde{\var} \Vert_2^2 \\ \nonumber
%     &> \obj(\Tilde{\var}, \param).
% \end{align}
% Thus, for all $\param \in \paramSet$, the minimizer $\varOpt(\param)$ of $\var \mapsto \obj(\var, \param)$ over $\constraintSet$ must satisfy $\Vert \varOpt(\param) - \Tilde{\var} \Vert_2 \leq \frac{2}{m} \max_{\param' \in \paramSet} \Vert \nabla \obj(\Tilde{\var}, \param') \Vert_2$. It thus follows that $\varOptSet = \{ \varOpt(\param): \param \in \paramSet\}$ is bounded.
% \end{proof}



\section{Preliminaries on PGD and SLS}
\label{sec: Preliminaries on Projected Gradient Descent (PGD) and System Level Synthesis (SLS)}

We introduce fundamental properties of the PGD dynamics (Sec. \ref{subsec: Iterative Gradient-Based Descent Methods}) and SLS output feedback error controller design.

\subsection{Iterative Gradient-Based Descent Methods}
\label{subsec: Iterative Gradient-Based Descent Methods}

% To recover an over-approximation of \eqref{Eqn: Set of All Parameterized Minimizers}, w
We first review the convergence properties of
PGD dynamics on strongly convex problems.
% iterative optimization routines that converge point-wise to $\arg\min_{\var} J(\var, \param)$, for any fixed $\param \in \paramSet$.
Consider
% the 
% To recover an over-approximation of \eqref{Eqn: Set of All Parameterized Minimizers}, we consider the 
% iterative projected gradient descent (PGD) 
% PGD dynamics
$f_{\text{PGD}}: \R^\varDim \times \R^\paramDim \times \R \ra \R^\varDim$ defined below.
% \antoine{Most probably it doesn't matter, but just curious why $J$ is defined in $\Theta$ while $f_{PGD}$ is defined on $\mathbb{R}^d$. Is that possible?}
Let $(\ratePGD_k: k \geq 0)$ denote a steplength sequence, and $\proj_\constraintSet$ be the Euclidean projection onto $\constraintSet$,
i.e., $\proj_\constraintSet(\var) := \arg\min_{\var' \in \constraintSet} \Vert \var' - \var \Vert_2$ for each $\var \in \R^\varDim$\footnote{
As proved in \cite[Ch. 7]{WrightRecht2022OptimizationForDataAnalysis}, under Assumption \ref{Assump: Constraint Set is Convex and Closed}, the projection $\proj_\constraintSet(\var)$ of any $\var \in \R^\varDim$ onto $\constraintSet$ is a unique, well-defined element of $\constraintSet$.}.
Then, for any fixed parameter $\param \in \paramSet$
% , starting from an initial iterate $\var_0$,
the PGD updates are:
{
\setlength{\abovedisplayskip}{3.5pt}
\setlength{\belowdisplayskip}{3.5pt}
\begin{align} 
\label{Eqn: PGD Updates Notation}
    \var_{k+1} &= \dynamicsPGD(\var_k, \param, \ratePGD_k) \\ \label{Eqn: Projected Gradient Descent (PGD)}
    &:= \proj_\constraintSet\left( \var_k - \ratePGD_k \nabla_\var \obj(\var_k, \param) \right).
\end{align}
}
% \antoine{Why do we have ``;'' in the cost definition here, but not elsewhere? I guess it should be ``,'' everywhere.}
% \antoine{I understand, maybe we need to state carefully why the system above do satisfy the typical regularity assumptions for the PGD to converge.}
% \frank{Doesn't Prop. \ref{Prop: Convergence of PGD} provide specific details w.r.t. PGD convergence? I'm happy to add more details here if that is what you are suggesting; I'm just unsure what specific exposition you are recommending to add.}
% \antoine{add the definition of $\proj$ in the notation or here, as it's quite a central operator here? Not really important I guess.}
% We first 
% % present a mild condition which bounds the initial iterate $\var_0$ of the PGD dynamics (Assumption \frank{INCOMPLETE; add reference}), and 
% establish necessary notation.

A key result we will leverage from optimization theory is that, under Assumptions \ref{Assump: Objective Properties} and \ref{Assump: Constraint Set is Convex and Closed}, for any fixed $\param \in \paramSet$, the iterative PGD dynamics \eqref{Eqn: PGD Updates Notation} on $\obj$
% \antoine{\eqref{Eqn: PGD Updates Notation}} 
converge exponentially to the global minimizer of $\obj$ over $\constraintSet$ \cite[Sec. 7.3.3]{WrightRecht2022OptimizationForDataAnalysis}.
Below, in Prop. \ref{Prop: Convergence of PGD}, we strengthen this result under our stated assumptions, 
% Assumption \ref{Assump: Parameter Set is Convex and Compact},
by showing that the PGD updates \eqref{Eqn: Projected Gradient Descent (PGD)} converge at a rate 
agnostic to the specific parameter value
% uniform with respect to candidate parameters 
$\param$ in $\paramSet$.

% We note that $\varOptSet$ is bounded under Assumptions \ref{Assump: Objective Properties} and \ref{Assump: Parameter Set is Convex and Compact}, which 

\begin{proposition}[\textbf{Boundedness of $\varOptSet$}]
\label{Prop: Minimizer Set is Bounded}
Under Assumption \ref{Assump: Objective Properties} and \ref{Assump: Parameter Set is Convex and Compact}, $\varOptSet$ is bounded.
\end{proposition}
Prop~\ref{Prop: Minimizer Set is Bounded} is well-known in optimization theory \cite[Sec. 2]{WrightRecht2022OptimizationForDataAnalysis}.

% \begin{proof}
% Fix $\Tilde{\var} \in \constraintSet$. By Assumption \ref{Assump: Objective Properties}, the map $\param \mapsto \Vert \nabla \obj(\Tilde{\var}, \param) \Vert_2$ is continuous, and thus, by Assumption \ref{Assump: Parameter Set is Convex and Compact}, attains a finite maximum over $\paramSet$.
% Thus, for any $\var \in \constraintSet$ satisfying $\Vert \var - \Tilde{\var} \Vert_2 > \frac{2}{m} \max_{\param' \in \paramSet} \Vert \nabla \obj(\Tilde{\var}, \param') \Vert_2$, and any $\param \in \paramSet$, we obtain from strong convexity of $\var \mapsto \obj(\var, \param)$ that:
% \begin{align} \nonumber
%     &\obj(\var, \param) \geq \obj(\Tilde{\var}, \param) + \nabla \obj(\Tilde{\var}, \param)^\top (\var - \Tilde{\var}) + \textstyle\frac{1}{2}m \Vert \var - \Tilde{\var} \Vert_2^2 \\ \nonumber
%     &\hspace{-3pt}\geq \obj(\Tilde{\var}, \param) - \Vert \nabla \obj(\Tilde{\var}, \param) \Vert_2 \Vert \var - \Tilde{\var} \Vert_2 + \textstyle\frac{1}{2}m \Vert \var - \Tilde{\var} \Vert_2^2 > \obj(\Tilde{\var}, \param).
% \end{align}
% Thus, for all $\param \in \paramSet$, the minimizer $\varOpt(\param)$ of $\var \mapsto \obj(\var, \param)$ over $\constraintSet$ must satisfy $\Vert \varOpt(\param) - \Tilde{\var} \Vert_2 \leq \frac{2}{m} \max_{\param' \in \paramSet} \Vert \nabla \obj(\Tilde{\var}, \param') \Vert_2$. It thus follows that $\varOptSet = \{ \varOpt(\param): \param \in \paramSet\}$ is bounded.
% \end{proof}


\begin{assumption}
% [\textbf{Bound on $\var_0$}]
\label{Assump: PGD Dynamics, Initial Iterate}
There is a known, bounded set $\varInitialSet \subseteq \constraintSet$ in which the initial iterate $\var_0$ of the PGD dynamics \eqref{Eqn: Projected Gradient Descent (PGD)} is fixed.
% \antoine{The verb ``drawn'' could refer to something stochastic, and could raise questions, maybe? Or I am not sure if I understand this Assumption.}
% The initial iterate $\var_0$ of the PGD dynamics \eqref{Eqn: Projected Gradient Descent (PGD)} is a fixed value drawn from a fixed, bounded set $\varInitialSet \subseteq \constraintSet$. 
% There exists
% \antoine{mismatch here: Either The initial iterate $\var_0$ is fixed, or it is an arbitrary element of a bounded set}
% \frank{Let me know if the above correction works.}
% Moreover, for some $\ratePGDMin, \ratePGDMax$ such that $0 < \ratePGDMin \leq \ratePGDMax < \frac{2}{\smoothnessConstant}$, 
% % the steplengths $(\ratePGD_k: k \geq 0)$ satisfy 
% we have 
% $\ratePGD_k \in [\ratePGDMin, \ratePGDMax]$ for each $k \geq 0$.
\end{assumption}

For any $k \geq 0$,
given $\var_0 \in \varInitialSet$, $\param \in \paramSet$, and $\ratePGD_{0:k-1} := (\ratePGD_0, \cdots, \ratePGD_{k-1}) \in \R^k$, let $\varMap_k(\var_0, \param, \ratePGD_{0:k-1})$ be the $k$-th iterate 
% \antoine{remove ``rollout''? Not sure if it's clear what is an ``iterate rollout'', I guess you mean just an ``iterate''? Not really important I think.} 
of the PGD dynamics \eqref{Eqn: Projected Gradient Descent (PGD)}.
% We define
For any $k \geq 0$, given $\var_0 \in \varInitialSet$ and $\ratePGD_{0:k-1} \in \R^k$, we define 
the forward reachable set for the PGD dynamics \eqref{Eqn: Projected Gradient Descent (PGD)} at iteration $k$ by:
\begin{align}
\label{Eqn: Forward Reachable Set, def}
    \varForwardReachableSet{k}(\var_0, \ratePGD_{0:k-1}) &:= \{ \varMap_k(\var_0, \param, \ratePGD_{0:k-1}): \param \in \paramSet \}.
\end{align}



\begin{proposition}
[\textbf{Uniform Convergence of PGD Dynamics}]
\label{Prop: Convergence of PGD}
Suppose Assumptions \ref{Assump: Objective Properties}, \ref{Assump: Constraint Set is Convex and Closed}, and \ref{Assump: Parameter Set is Convex and Compact} hold. Fix $\var_0 \in \varInitialSet$, let $\ratePGDMin, \ratePGDMax \in (0, \frac{2}{\smoothnessConstant})$ be given such that $\ratePGDMin \leq \ratePGDMax$, and define:
\begin{align}
    \label{eq:PGD_convergence_rate}
    \textstyle\convergenceRatePGD := \max_{\ratePGD \in [\ratePGDMin, \ratePGDMax]} \big\{\max\{|1 - \ratePGD \strongConvexityConstant|, |1 - \ratePGD \smoothnessConstant| \} \big\} < 1.
\end{align}
Fix $\ratePGD_k \in [\ratePGDMin, \ratePGDMax]$, $k \geq 0$.
Then, 
% under the PGD updates \eqref{Eqn: Projected Gradient Descent (PGD)}, 
for any $\param \in \paramSet$, $k \in \N$:
\begin{align}
    \label{Eqn: Exponential Bound for PGD Convergence}
    \Vert \varMap_k(\var_0, \param, \ratePGD_{0:k-1}) - \varOpt(\param) \Vert_2 &\leq \convergenceRatePGD^k \cdot \Vert \var_0 - \varOpt(\param) \Vert_2. 
\end{align}
\end{proposition}

\begin{proof}
Below, we fix $\var_0 \in \varInitialSet$, $\param \in \paramSet$, and $\ratePGD_{0:k-1} \in [\ratePGDMin, \ratePGDMax]$, and write $\var_j = \varMap_j(\var_0, \param, \ratePGD_{0:j-1})$ for $j \in [0, k]$.
From \cite[Prop. 7.8]{WrightRecht2022OptimizationForDataAnalysis}, $\forall \ \param \in \paramSet$ and $\ratePGD_k \geq 0$, $\varOpt(\param)$ is the unique fixed point of the PGD updates \eqref{Eqn: Projected Gradient Descent (PGD)}. Thus, for any $\param \in \paramSet$ and $j \geq 0$:
\begin{align} \nonumber
    \Vert \var_{j+1} -  &\varOpt(\param) \Vert_2  = \big\Vert \proj_\constraintSet\big(\var_j - \ratePGD_j \nabla \obj(\var_j, \param) \big) \\ \nonumber
    &\hspace{5mm} - \proj_\constraintSet\big(\varOpt(\param) - \ratePGD_j \nabla \obj\big(\varOpt(\param), \param \big) \big) \big\Vert_2 \\
    \label{Eqn: Apply Recht and Wright, for Exponential Bound for PGD convergence}
    \leq &\max\big\{ |1 - \ratePGD_j \strongConvexityConstant|, |1 - \ratePGD_j \smoothnessConstant| \big\} \cdot \Vert \var_j - \varOpt(\param) \Vert_2 \\ \nonumber
    \leq &\convergenceRatePGD \cdot \Vert \var_j - \varOpt(\param) \Vert_2,
\end{align}
%  \antoine{There is an extra comma no?}
%  \frank{Do you mean at the very end of the above equation block? I meant for it to be there, as in \say{Thus, for any $\param \in \paramSet$ and $j \geq 0$, XYZ with (7), where ...}}
% \antoine{I mean here: $\max\big\{ |1 - \ratePGD_j, \strongConvexityConstant|$. At least, I don't know this notation and I am not sure if it is in the notation section.}
% }
% \antoine{Could you check this: $\ \big\Vert \proj_\constraintSet\big(\var_j - \param_j \nabla \obj(\var_j, \param) \big)$ Intuitively, I would say $\ \big\Vert \proj_\constraintSet\big(\var_j - \eta_j \nabla \obj(\var_j, \param) \big)$ }
% \frank{Yes, this was a typo which has now been fixed; thanks}
where \eqref{Eqn: Apply Recht and Wright, for Exponential Bound for PGD convergence} follows from \cite[Eqn. (7.8)]{WrightRecht2022OptimizationForDataAnalysis}. Recursively applying the above for $j \in [0, k-1]$, we obtain \eqref{Eqn: Exponential Bound for PGD Convergence}, as desired.
\end{proof}


Next, we use Prop. \ref{Prop: Convergence of PGD} to recast the problem of over-approximating $\varOptSet$ as the computation of forward reachable sets for PGD \eqref{Eqn: PGD Updates Notation}.
The following theorem asserts that $\varOptSet$ can be over-approximated by fixing any $\var_0 \in \varInitialSet$ and $\ratePGD_{0:k-1} \in [\ratePGDMin, \ratePGDMax]^k$, computing $\varForwardReachableSet{k}(\var_0, \ratePGD_{0:k-1})$, and inflating $\varForwardReachableSet{k}(\var_0, \ratePGD_{0:k-1})$ (in the 2-norm) by the $\convergenceRatePGD$-exponentially shrinking bound furnished by Prop. \ref{Prop: Convergence of PGD}.

\begin{theorem}
(\textbf{Over-Approximating $\varOptSet$ using the Forward Reachable Sets of PGD Dynamics})
\label{Thm: Over-approximating Minimizer Set using Forward Reachable Sets}
Fix $\var_0 \in \varInitialSet$ and $\ratePGD_{0:k-1} \in [\ratePGDMin, \ratePGDMax]^k$ arbitrarily, and define $\convergenceRadius_k(\var_0) := \convergenceRatePGD^k \cdot \max_{\param \in \paramSet} \Vert \varOpt(\param) - \var_0 \Vert_2$ for each $k \geq 0$. Then:
% , for any $k \geq 0$:
% \begin{align}
%     \label{Eqn: Over-approximating Minimizer Set using Forward Reachable Sets}
%     \varOptSet \subseteq \varForwardReachableSet{k}(\var_0, \ratePGD_{0:k-1}) \oplus \convergenceRadius_k(\var_0) I_\varDim \ballUnit_2^\varDim(\textbf{0}_n).
% \end{align}
\begin{align}
    \label{Eqn: Forward Reachable Sets approach Minimizer Sets}
    \dist\big(\varOptSet, \varForwardReachableSet{k}(\var_0, \ratePGD_{0:k-1}) \big) \leq \convergenceRadius_k(\var_0), \hspace{5mm} \forall \ k \geq 0.
\end{align}
\end{theorem}

\begin{proof}
Fix $\var_0 \in \varInitialSet$ and $\ratePGD_{0:k-1} \in [\ratePGDMin, \ratePGDMax]^k$.
% Note that \eqref{Eqn: Over-approximating Minimizer Set using Forward Reachable Sets} is equivalent to the following assertion---
By maximizing the right-hand side of \eqref{Eqn: Exponential Bound for PGD Convergence} over all $\param \in \paramSet$, we obtain that $\Vert \varMap_k(\var_0, \param, \ratePGD_{0:k-1}) - \varOpt(\param) \Vert_2 \leq \convergenceRadius_k(\var_0)$ for each $\param \in \paramSet$, from which 
% \eqref{Eqn: Over-approximating Minimizer Set using Forward Reachable Sets} 
\eqref{Eqn: Forward Reachable Sets approach Minimizer Sets}
readily follows.  
% \begin{align}
%     \Vert \varMap_k(\param) - \varOpt(\param) \Vert_2
% \end{align}
\end{proof}

\begin{remark}
\label{Remark: Upper bound for convergence radius}
The upper bound $\convergenceRadius_k(\var_0)$, though $\var_0$-dependent, is continuous in $\var_0$ and thus uniformly bounded over the bounded set $\varInitialSet$ under Assumption \ref{Assump: PGD Dynamics, Initial Iterate}.

% ~\\
% \frank{Add remark about how an upper bound for $\convergenceRadius_k$ can be constructed without \textit{a priori} knowledge of $\varInitialSet$.}
% ~\\

\end{remark}


% \antoine{What didn't become clear here: do both $\theta$ and the trajectory initial set $\xi$ are uncertain? or just $\theta$? Maybe I juts missed it.}
% \frank{This is an important point that I think we should discuss, e.g., in person.
% I had intended to introduce the theory in a way that provides some flexibility for the \say{ego agent} to choose which $\var_0$ to use, as long as it satisfies the bound prescribed by Assumption \ref{Assump: PGD Dynamics, Initial Iterate}. 
% But we could potentially save ourselves a lot of trouble by just saying that $\var_0$ is permanently fixed at some value.
% }
% ~\\
% \frank{Consensus: Just present the simplest case and add a remark about how the discussion generalizes readily to a broader setting.
% }
% \brendan{On the code side, nothing would need to change - we already support uncertain $\xi_0$. However, I agree that this is not super important for our main message, so if it would add overhead to include in the written presentation we should probably leave it out.}


\subsection{SLS-based Output Feedback Error Controller Design}
\label{subsec: System-Level Synthesis (SLS)-based Output Feedback Error Controller Design}


% \subsection{System Level Synthesis (SLS) for PGD Dynamics}
% \label{subsec: System Level Synthesis (SLS) for PGD Dynamics}

To over-approximate
% the target set
$\varOptSet$
via Thm. \ref{Thm: Over-approximating Minimizer Set using Forward Reachable Sets}, we leverage the SLS approach in~\cite{Leeman2025RobustNonlinearOptimalControlviaSLS, leeman2026vision} to compute a forward reachable set for the PGD dynamics \eqref{Eqn: Projected Gradient Descent (PGD)}, with uncertainty over 
% the intent parameter
$\param \in \paramSet$ modeled as initial state uncertainty. 
Although many methods exist for computing reachable sets~\cite{Bansal2017HJBReachability, immrax}, we adopt the 
SLS-based methods
in~\cite{Leeman2025RobustNonlinearOptimalControlviaSLS, leeman2026vision} due to their tractability for high-dimensional system dynamics \cite{Leeman2024FastSLS}. 
% \antoine{cite something here? any of the following: fast-sls, seeing-is-believing RSS paper, my guaranteed RNMPC ICRA,  GPU solver RSS paper?}
% framework due to its ability to efficiently generate tight over-approximations of forward reachable sets for high-dimensional nonlinear dynamics.
% We adapt the 
% approach 
% % SLS framework
% in \cite{Leeman2025RobustNonlinearOptimalControlviaSLS, leeman2026vision} to closely over-approximate forward reachable sets for PGD dynamics and, via Thm. \ref{Thm: Over-approximating Minimizer Set using Forward Reachable Sets}, produce outer approximations of $\varOptSet$.

We consider the following dynamics over a horizon $[0, \iterationHorizon]$, with variable-parameter tuples $\state := (\var, \param)$ as states, PGD steplengths $\ratePGD$ as control inputs, and variables $\var$ as outputs:
% and 
% % initial uncertainty over the parameter
% parameter uncertainty
% over
% $\paramSet$:
% $\param_0 \in \paramSet$:
% \antoine{The $\theta_0$ is fixed no? I am not sure if it's clear if we say ``uncertainty over $\theta_0$''.}
\begin{subequations}
\label{Eqn: True Dynamics}
\begin{alignat}{2} \label{Eqn: True Dynamics, Variable}
    \var_{k+1} &= \dynamicsPGD(\var_k, \param_k, \ratePGD_k),
    \hspace{5mm} &&\forall \ k \in [0, \iterationHorizon - 1], \\ \label{Eqn: True Dynamics, Parameter}
    \param_{k+1} &= \param_k, \hspace{5mm} &&\forall \ k \in [0, \iterationHorizon - 1], \\
    \outputs_k &= \var_k, \hspace{5mm} &&\forall \ k \in [0, \iterationHorizon - 1].
\end{alignat}
\end{subequations}
% \antoine{I am not sure if we speak about $\theta^{nom}$. Is it also optimized? It could/should be. Also, it should be at least defined/mentionned.}
% \frank{$\paramNominal_0$ is mentioned immediately before \eqref{Eqn: Nominal Dynamics}. Doesn't the NLP we formulate at the end already account for the optimization over the entire nominal state-control trajectory (along with the error feedback gains through the system response), which includes $\paramNominal_0$?
% }
% \frank{In a comment, emphasize we are optimizing over nominal $\param$ as part of the NLP.}

For notational ease, we define $\dynamicsFullState: \R^{\varDim} \times \R^\paramDim \times \R \ra \R^{\varDim + \paramDim}$ by 
% \antoine{$f$ seems to be 3 arguments, while the definition has only 2?}
$\dynamicsFullState(\var, \param, \ratePGD) := \big(\dynamicsPGD(\var, \param, \ratePGD), \param \big)$, simplifying \eqref{Eqn: True Dynamics, Variable}-\eqref{Eqn: True Dynamics, Parameter} to:
\begin{align}
    \label{Eqn: True Dynamics, Full State}
    \state_{k+1} &= \dynamicsFullState(\state_k, \ratePGD_k), \hspace{5mm} \forall \ k \in [0, \iterationHorizon - 1].
\end{align}
We also write $\stateStacked := (\state_0, \cdots, \state_\iterationHorizon) \in \R^{(\varDim + \paramDim)(\iterationHorizon + 1)}$, and for any $k_1, k_2 \in [0, \iterationHorizon]$ with $k_1 \leq k_2$, we write $\stateStacked_{k_1:k_2} := (\state_{k_1}, \cdots, \state_{k_2}) \in \R^{(\varDim + \paramDim)(k_2 - k_1 + 1)}$.
% We also write $\varStacked := (\var_0, \cdots, \var_\iterationHorizon) \in \R^{\varDim(\iterationHorizon + 1)}$ and $\paramStacked := (\param_0, \cdots, \param_\iterationHorizon) \in \R^{\paramDim(\iterationHorizon + 1)}$. 
% Further, for any $k_1, k_2 \in [\iterationHorizon]$ with $k_1 \leq k_2$, we write $\varStacked_{k_1:k_2} := (\var_{k_1}, \cdots, \var_{k_2}) \in \R^{\varDim(k_2 - k_1 + 1)}$ and $\paramStacked_{k_1:k_2} := (\param_{k_1}, \cdots, \param_{k_2}) \in \R^{\paramDim(k_2 - k_1 + 1)}$.

To satisfy Prop. \ref{Prop: Convergence of PGD} and ensure the convergence of PGD dynamics, we enforce that 
% the steplengths 
$(\ratePGD_k: k \in [0, \iterationHorizon])$ satisfy:
% for some pre-specified $\ratePGDMin$ and $\ratePGDMax$ satisfying $0 < \ratePGDMin \leq \ratePGDMax < 2/\smoothnessConstant$:
\begin{align} 
\label{Eqn: Steplength bound constraints}
    \ratePGDMin \leq \ratePGD_k \leq \ratePGDMax, \hspace{5mm} \forall \ k \in [0, \iterationHorizon-1].
\end{align}
Define $\ratePGDBoundSlope_1 := (\zero_{\varDim + \paramDim}, 1), \ratePGDBoundSlope_2 := (\zero_{\varDim + \paramDim}, -1) \in \R^{\varDim + \paramDim + 1}$ and $\ratePGDBoundOffset_1 := - \ratePGDMax, \ratePGDBoundOffset_2 := \ratePGDMin$. 
% \antoine{I think this is a personal taste, but I find it cleaner to always define vector as columns, and hence use a transponse when they are rows for consistency and ease of read.
% $\ratePGDBoundSlope_1 := (\zero_{\varDim + \paramDim}^\top, 1), \ratePGDBoundSlope_2 := (\zero_{\varDim + \paramDim}^\top, -1) \in \R^{\varDim + \paramDim + 1}$.
% }
% We the define $\inequalityConstraints: \R^{\varDim + \paramDim + 1} \ra \R^2$ component-wise by $\inequalityConstraints_p(\state, \ratePGD) := \ratePGDBoundSlope_p^\top (\state, \ratePGD) + \ratePGDBoundOffset_p \leq 0$ for each $p \in [2]$ and any $(\state, \param) \in \R^{\varDim + \paramDim + 1}$. 
We can then rewrite \eqref{Eqn: Steplength bound constraints} as $\ratePGDBoundSlope_p^\top (\state_k, \ratePGD_k) + \ratePGDBoundOffset_p \leq 0$ for each $p \in [2]$.

\begin{remark}
Additional affine constraints of the form \say{$\ratePGDBoundSlope_p^\top (\state_k, \ratePGD_k) + \ratePGDBoundOffset_p \leq 0$} 
% which encode constraints over states $\state_k$ 
% can 
may
% also 
be enforced
% , e.g., 
to bound the distance between each variable iterate $\var_k$ and the minimizer set $\varOptSet$.
\end{remark}

To construct reachable sets for the PGD dynamics $\dynamicsFullState$ in \eqref{Eqn: True Dynamics} over $\param \in \paramSet$ with low conservativeness,
% \antoine{I think the word "tight" is not good here, as the reachable sets aren't really ``tight'', they have low conservatism. They would be tight to LTV.}
% \frank{Can I say \say{closely overapproximates}?}
% a tight over-approximation of the target set $\varOptSet$
at each iteration $k \in [0, N]$, we aim to select appropriate steplengths $\ratePGD_k \in [\ratePGDMin, \ratePGDMax]$ across the iteration horizon $[0, \iterationHorizon]$ as functions of past output realizations 
$\varStacked_{0:k}$, i.e., design a causal feedback law $\feedbackMap_k: \R^{\varDim(k+1)} \ra \R$, and set 
$\ratePGD_k := \feedbackMap_k(\stateStacked_{0:k})$.
Although $\varOptSet$ can also be outer-approximated via computing forward reachable sets for the PGD dynamics with \textit{fixed} steplengths,
% (per Thm. \ref{Thm: Over-approximating Minimizer Set using Forward Reachable Sets}),
doing so may result in more conservative outer approximations of $\varOptSet$.
% \footnote{Although Thm. \ref{Thm: Over-approximating Minimizer Set using Forward Reachable Sets} implies that forward reachable sets for the PGD dynamics can also be computed using fixed steplengths, doing so may generate more conservative outer approximations of $\varOptSet$, as we demonstrate empirically in Sec. \ref{sec: Experiments}}.
% $(\varStacked_{0:k}, \paramStacked_{0:k})$.
% $\{(\var_j, \param_j): j \in [0, k] \}$. 

% $\ratePGD_k := \feedbackMap_k(\varStacked_{0:k}, \paramStacked_{0:k})$.
\looseness-1Since optimizing over all causal feedback laws is in general intractable, we mimic \cite{leeman2026vision, Leeman2025RobustNonlinearOptimalControlviaSLS},
% \antoine{I think the correct citation is this one \cite{leeman2026vision}. I would have both here.}
which instead compute a \textit{nominal state-control trajectory} $(\stateNominalStacked, \ratePGDNominalStacked)$ and an \textit{affine causal output error feedback law} mitigating the deviation of the trajectory rollouts from $(\stateNominalStacked, \ratePGDNominalStacked)$.
We first reformulate the true nonlinear system \eqref{Eqn: True Dynamics}, encoding initial uncertainty over the parameter $\param_0$, 
% \antoine{Maybe there is still something I miss, but why do we speak about ``initial uncertainty $\theta$''? The set $\Theta$ does not change right?}
as the sum of a nominal nonlinear system \textit{without initial state uncertainty}~\eqref{Eqn: Nominal Dynamics} and linear time-varying (LTV) error dynamics~\eqref{Eqn: State Error LTV Dynamics}.
We then design (a) state and control trajectories $\stateNominalStacked
% \stateNominalStacked 
:= (\varNominal_0, \paramNominal_0, \cdots, \varNominal_\iterationHorizon, \paramNominal_\iterationHorizon) 
\in \R^{(\varDim + \paramDim)(\iterationHorizon + 1)}$ and $\ratePGDNominalStacked := (\ratePGDNominal_0, \cdots, \ratePGDNominal_\iterationHorizon) \in \R^{\iterationHorizon + 1}$,
% and $\paramNominalStacked := (\paramNominal_0, \cdots, \paramNominal_\iterationHorizon) \in \R^{\paramDim(\iterationHorizon + 1)}$,
satisfying the nominal dynamics:
% \textit{without any initial uncertainty over $\paramNominal_0$}:
\begin{align}
\label{Eqn: Nominal Dynamics}
    \stateNominal_{k+1} &= \dynamicsFullState(\stateNominal_k, \ratePGDNominal_k), \hspace{5mm} \forall \ k \in [0, \iterationHorizon - 1].
\end{align}
% \begin{subequations}
% \label{Eqn: Nominal Dynamics}
% \begin{alignat}{2} \label{Eqn: Nominal Dynamics, Variable}
%     \varNominal_{k+1} &= \dynamicsPGD(\stateNominal_k, \ratePGDNominal_k), \hspace{5mm} &&\forall \ k \in [0, \iterationHorizon-1], \\
%     \label{Eqn: Nominal Dynamics, Parameter}
%     \paramNominal_{k+1} &= \paramNominal_k, \hspace{5mm} &&\forall \ k \in [0, \iterationHorizon-1].
% \end{alignat}
% \end{subequations}
and (b) a set of causal affine output error feedback gains $\setFeedbackGains := \{\feedbackGain_{k,j} \in \R^{1 \times \varDim}, \ k,j \in [0,\iterationHorizon], k \geq j \}$
% \begin{align} 
% \label{Eqn: Feedback Gains}
%     \setFeedbackGains &:= \{\feedbackGain_{k,j} \in \R^{1 \times \varDim}, \ k,j \in [0,\iterationHorizon-1], k \geq j \},
% \end{align}
% \antoine{Probably need to justify the choice of controller here? Why SLS and not something else? Is it for scalability?}
% \frank{Justifying this point is good / necessary, but should be done in the introduction, I think.}
% \antoine{Sounds good! A sentence to repeat would also be good here. Could be combined just a justification of why adaptive step size yields tighter tubes.}
% ~\\
% \glen{agreed with antoine -- should re-motivate SLS here. should also discuss explicitly: we are not assuming the other agent is solving this SLS problem to solve their optimization problem. rather, we are using the policy to construct the tightest overapproximating set of true minimizers of the other agent (the other agent just optimizes using standard PGD)}
% \frank{I thought we were moving away from a multi-agent problem formulation to a more general one, per your comments regarding the introduction section. I'll see what we can do though.}
parameterizing the following causal affine output error feedback map:
\begin{align} \nonumber
    \ratePGD_k &= \feedbackMap_k(\outputStacked_{0:k}) \\
    \label{Eqn: Affine Output Feedback Map}
    &= \ratePGDNominal_k + \textstyle\sum_{j=0}^k \feedbackGain_{k,j}(\var_{k-j} - \varNominal_{k-j}),
    \hspace{3mm} \forall \ k \in [0, \iterationHorizon].
    % \\ \nonumber
    % &\hspace{1cm} + \sum_{j=0}^k\feedbackGain_{k,j}(\param_{k-j} - \paramNominal_{k-j}).
\end{align}

% \antoine{There is a lot of text that in emph or italic in the paragraph above, in contrast to the rest of the document. Usually chatGPT uses a lot of italic, so I would suggest removing most of them.}
\begin{remark}
\label{Remark: Choice of Initial Nominal Parameter, paramNominal 0}
Designing $\stateNominalStacked$ requires choosing
an initial nominal state $\varNominal_0$, which we assume to be the true initial state $\var_0$ for simplicity, and an initial nominal parameter $\paramNominal_0$, which is then fixed across iterations per \eqref{Eqn: True Dynamics, Parameter}.
% where the nominal parameter value $\paramNominal_0 \in \paramSet$ is initialized without uncertainty
In practice, 
% the nominal intent parameter 
$\paramNominal_0 \in \paramSet$ can be selected to minimize the worst-case parameter deviation over $\paramSet$, i.e., $\paramNominal_0 \in \arg\min_{\param \in  \paramSet} \max_{\param' \in \paramSet} \Vert \param - \param' \Vert_2$, .
\end{remark}

% \begin{remark}
% \label{Remark: Singleton Parameter Set yields Zero Error}
% % In the absence of uncertainty over $\param \in \paramSet$, i.e., w
% When $\paramSet$ is a singleton set, we have $\param_k = \paramNominal_k$ for each $k \in [0, \iterationHorizon]$, from which we can derive by induction over \eqref{Eqn: Nominal Dynamics} and \eqref{Eqn: Affine Output Feedback Map} that $\ratePGD_k = \ratePGDNominal_k$ and $\var_k = \varNominal_k$ for each $k \in [0, \iterationHorizon]$ as well. 
% Thus, the true and nominal dynamics would coincide, with zero state and control (i.e., steplength) error, and we recover the (singleton) minimizer set.
% % \antoine{
% % I would say that we haven't loss anything, i.e,. something like this``
% % We converge immediately to the singleton minimizer set''}
% \end{remark}


\section{Methods}
\label{sec: Methods}

% Prop. \ref{Prop: Convergence of PGD} and Thm. \ref{Thm: Over-approximating Minimizer Set using Forward Reachable Sets} enable our original problem of outer-approximating $\varOptSet$ to be recast as that of over-approximating the reachable sets of the PGD dynamics. 
In this section, we outer-approximate the minimizer set $\varOptSet$ using Thm. \ref{Thm: Over-approximating Minimizer Set using Forward Reachable Sets} by computing reachable sets for the PGD dynamics with SLS.
% for smooth and smoothed dynamics.
% In Sec. \ref{subsec: LTV Error System for Smooth PGD Dynamics},
For differentiable PGD dynamics, we formulate a linear time-varying (LTV) system that describes the propagation of the error between nominal and true trajectories
% , i.e., $(\stateNominalStacked, \ratePGDNominalStacked)$ and $(\stateStacked, \ratePGDStacked)$, 
across iterations (Sec. \ref{subsec: LTV Error System for Smooth PGD Dynamics}), and construct error bounds across iterations (Sec. \ref{subsec: Linearization Error Bounds}).
We then generalize our methods
% in Sec. \ref{subsec: LTV Error System for Smooth PGD Dynamics}-\ref{subsec: Linearization Error Bounds}
to the setting in which the PGD dynamics are non-differentiable
via a smoothing technique
(Sec. \ref{subsec: Smoothing PGD Dynamics when Non-Differentiable}).


% ~\\
% \antoine{The bridge between the previous section and the current one here is a bit abrupt. I would add a sentence like ``Prop. 2 and Cor. 1 reduce the original problem to reachable-set over-approximation of PGD. Section IV shows how to compute those reachable sets using SLS for smooth and smoothed dynamics.'' }
% ~\\
% \frank{(If time and space allow) Add exposition distinguishing the settings where the PGD dynamics are smooth and when they are not, each of which will be explained in its own subsection below.
% }
% ~\\


\subsection{LTV Error System for Smooth PGD Dynamics}
\label{subsec: LTV Error System for Smooth PGD Dynamics}

% \antoine{You argue that every rollout gets close to the bounded minimizer set, and therefore all reachable sets are uniformly bounded. That is plausible, but the step needs more care because closeness as $k \to \infty$ does not by itself give a uniform bound over all $k$. I think the result is correct, only that some steps are missing maybe?}
% \frank{Thm. \ref{Thm: Over-approximating Minimizer Set using Forward Reachable Sets} gives a bound, written as $\convergenceRadius_k(\var_0)$, which is uniform over any choice of $\ratePGD_{0:k-1} \in [\ratePGDMin, \ratePGDMax]^k$ and $\param \in \paramSet$, though it may depend on the initial iterate $\var_0$. The dependence on $\var_0$ can be removed through Assumption \ref{Assump: PGD Dynamics, Initial Iterate} and I intend to explain this in Remark \ref{Remark: Upper bound for convergence radius}.
% Let me know if this would address your comment or if I'm misunderstanding what you mean by \say{uniformly,} misunderstanding which claim you said is \say{plausible but needs more care}, or misunderstanding something else.}

% To design the error feedback gains, we 
We
derive the LTV error system for the true dynamics \eqref{Eqn: True Dynamics},
as a step towards designing error feedback gains,
for the case in which $\constraintSet$ is $\R^\varDim$ or an affine subspace of $\R^\varDim$.
% \antoine{The last word ``is'' looks wrong, no? Or grammatically, I don't understand the sentence. I would say it should be ``or'', i.e., there are two possible constraint set assumed here.}
% \frank{It's \say{or,} thanks.}
In this case, the projection operator $\proj_\constraintSet(\cdot)$ is affine, 
% \antoine{$\proj_\C$: I think this is an old notation?}
and thus $\dynamicsFullState$ in \eqref{Eqn: True Dynamics}
is differentiable and satisfies
% has Lipschitz continuous partial derivatives under 
Assumption \ref{Assump: Objective Properties}.
% The general case is prolonged to Sec. \ref{subsec: Smoothing PGD Dynamics when Non-Differentiable}. 
In Sec. \ref{subsec: Smoothing PGD Dynamics when Non-Differentiable}, we study the general setting where $\dynamicsFullState$ may be non-differentiable.

\looseness=-1The LTV error system admits, across iterations $k \in [0, \iterationHorizon]$, the state errors 
$\stateError_k := \state_k - \stateNominal_k$
% $\varError_k := \var_k - \varNominal_k$ and $\paramError_k := \param_k - \paramNominal_k$
% across iterations $k \in [\iterationHorizon]$ 
as states and steplength errors $\ratePGDError_k := \ratePGD_k - \ratePGDNominal_k$ 
% across iterations $k \in [\iterationHorizon]$ 
as controls. 
Its gain and actuation matrices for each $k \in [0, \iterationHorizon]$ are computed
via the Jacobian linearization of \eqref{Eqn: True Dynamics, Full State} at the nominal state-control 
pair $(\stateNominal_k, \ratePGDNominal_k)$
% trajectory $\stateNominalStacked$ for each $k \in [0, \iterationHorizon]$
via: 
% \antoine{With this notation ``$k \in [\iterationHorizon]$'', according to the notation section, I guess you want to say that we don't need the Jacobian at the initial time step. Is that correct? Or why don't you defined the Jacobian at $k=0$? I think there are other places where $k=0$ is excluded, and I am not sure if I fully understand why or why not.}
% (we note that \eqref{Eqn: True Dynamics, Parameter} is already linear):
{
\setlength{\abovedisplayskip}{3.5pt}
\setlength{\belowdisplayskip}{3.5pt}
\begin{subequations}
\label{Eqn: Jacobian Matrices}
\begin{align}
    \label{Eqn: Jacobian Matrix, State}
    \jacobianState_k(\stateNominal_k, \ratePGDNominal_k) &:= \textstyle\frac{\partial \dynamicsFullState}{\partial \state}(\stateNominal_k, \ratePGDNominal_k), \\
    \label{Eqn: Jacobian Matrix, Control}
    \jacobianControl_k(\stateNominal_k, \ratePGDNominal_k) &:= \textstyle\frac{\partial \dynamicsFullState}{\partial \ratePGD}(\stateNominal_k, \ratePGDNominal_k),
\end{align}
\end{subequations}
}
% \begin{subequations}
% \label{Eqn: Jacobian Matrices}
% \begin{align}
%     \jacobianState_k(\stateNominal_k, \ratePGDNominal_k) &:= \frac{\partial \dynamicsPGD}{\partial \var}(\stateNominal_k, \ratePGDNominal_k), \\
%     \jacobianState_k(\stateNominal_k, \ratePGDNominal_k) &:= \frac{\partial \dynamicsPGD}{\partial \param}(\stateNominal_k, \ratePGDNominal_k), \\
%     \jacobianControl_k(\stateNominal_k, \ratePGDNominal_k) &:= \frac{\partial \dynamicsPGD}{\partial \ratePGD}(\stateNominal_k, \ratePGDNominal_k),
% \end{align}
% \end{subequations}
\hspace{-1mm}and the \textit{linearization error} $\linearizationError_k(\state_k, \ratePGD_k, \stateNominal_k, \ratePGDNominal_k) \in \R^{\varDim + \paramDim}$ as:
{
\setlength{\abovedisplayskip}{3.5pt}
\setlength{\belowdisplayskip}{3.5pt}
\begin{align}
\label{Eqn: Linearization Error}
    \linearizationError_k(\state_k, \ratePGD_k, \stateNominal_k, \ratePGDNominal_k) &= \dynamicsFullState(\state_k, \ratePGD_k) - \dynamicsFullState(\stateNominal_k, \ratePGDNominal_k) \\ \nonumber 
    &\hspace{5mm} - \jacobianState_k(\stateNominal_k, \ratePGDNominal_k) \stateError_k - \jacobianControl_k(\stateNominal_k, \ratePGDNominal_k) \ratePGDError_k.
\end{align}
}
\hspace{-2mm}We then use
\eqref{Eqn: Jacobian Matrices} and \eqref{Eqn: Linearization Error} to describe the propagation of parameter errors as the following LTV system
% $\varError_k$
% % $(\varError_k, \paramError_k, \ratePGDError_k)$
% using the following set of LTV dynamics, in which
while expressing
each linearization error as a disturbance\footnote{
In Sec. \ref{subsec: Smoothing PGD Dynamics when Non-Differentiable}, which considers the setting of non-differentiable PGD dynamics, the disturbance $\disturbance_k$ will also incorporate the approximation error between the PGD dynamics and a differentiable proxy.
% \frank{
% Mention how in subsequent subsections, the disturbance terms will also include approximation errors through smoothing}
% }
}
$\disturbance_k$:
% $\disturbance_k = \linearizationError_k(\state_k, \ratePGD_k, \stateNominal_k, \ratePGDNominal_k)$:
{
\setlength{\abovedisplayskip}{3.5pt}
\setlength{\belowdisplayskip}{3.5pt}
\begin{align} \label{Eqn: Disturbance}
    \disturbance_k &= \linearizationError_k(\state_k, \ratePGD_k, \stateNominal_k, \ratePGDNominal_k), \\
    \nonumber
    \stateError_{k+1} &:= \state_{k+1} - \stateNominal_{k+1} = \dynamicsFullState(\state_k, \ratePGD_k) - \dynamicsFullState(\stateNominal_k, \ratePGDNominal_k) \\ \label{Eqn: State Error LTV Dynamics}
    &= \jacobianState_k(\stateNominal_k, \ratePGDNominal_k) \stateError_k + \jacobianControl_k(\stateNominal_k, \ratePGDNominal_k) \ratePGDError_k + \disturbance_k,
\end{align}
}
% \begin{align} \nonumber
%     \varError_{k+1} &:= \var_{k+1} - \varNominal_{k+1} \\ \nonumber
%     &= \dynamicsPGD(\state_k, \ratePGD_k) - \dynamicsPGD(\stateNominal_k, \ratePGDNominal_k) \\ \nonumber
%     &= \jacobianState_{\var\var}(\stateNominal_k, \ratePGDNominal_k) \varError + \jacobianState_{\var\param}(\stateNominal_k, \ratePGDNominal_k) \paramError \\ \nonumber
%     &\hspace{1cm} + \jacobianControl_{\var\param}(\stateNominal_k, \ratePGDNominal_k) \ratePGDError + \disturbance_k,
% \end{align}
% \begin{align} \nonumber
%     \varError_{k+1} &:= \var_{k+1} - \varNominal_{k+1} \\ \nonumber
%     &= \dynamicsPGD(\state_k, \ratePGD_k) - \dynamicsPGD(\stateNominal_k, \ratePGDNominal_k) \\ \nonumber
%     &= \jacobianState_{\var\var}(\stateNominal_k, \ratePGDNominal_k) \varError + \jacobianState_{\var\param}(\stateNominal_k, \ratePGDNominal_k) \paramError \\ \nonumber
%     &\hspace{1cm} + \jacobianControl_{\var\param}(\stateNominal_k, \ratePGDNominal_k) \ratePGDError + \disturbance_k,
% \end{align}
% Together with the parameter error dynamics $\paramError_{k+1} = \param_{k+1} - \paramNominal_{k+1} = \param_k - \paramNominal_k = \paramError_k$, we then obtain the following LTV error dynamics:
% % the true and nominal parameter dynamics, i.e., \eqref{Eqn: True Dynamics, Parameter} and \eqref{Eqn: True Dynamics, Parameter}, gives:
% \begin{align} \nonumber
%     \begin{bmatrix}
%         \varError_{k+1} \\
%         \paramError_{k+1}
%     \end{bmatrix} &= 
%     \underbrace{\begin{bmatrix}
%         \jacobianState_{k,\var\var}(\stateNominal_k, \ratePGDNominal_k)  &  \jacobianState_{k,\var\param}(\stateNominal_k, \ratePGDNominal_k) \\
%         O_{\paramDim \times \varDim} & I_{\varDim}
%     \end{bmatrix}}_{:= A_k(\stateNominal_k, \ratePGDNominal_k)}
%     \begin{bmatrix}
%         \varError_k \\
%         \paramError_k
%     \end{bmatrix} \\
%     \label{Eqn: }
%     &\hspace{1cm} + \underbrace{\begin{bmatrix}
%         \jacobianControl_k \\ O_{d \times 1}
%     \end{bmatrix}}_{:= \ B_k(\stateNominal_k, \ratePGDNominal_k)}
%     \ratePGDError_k 
%     + \disturbance_k, \\
%     \ratePGDError_k &= \sum_{j=0}^k \feedbackGain_{k,j} \varError_{k-j}.
%     % \big[\feedbackGain_{k,j} \varError_{k-j} + \feedbackGain_{k,j} \paramError_{k-j} \big].
% \end{align}
\hspace{-1.5mm}Meanwhile, we define the \textit{control error} $\ratePGDError_k := \ratePGD_k - \ratePGDNominal_k$ for each $k \in [0, \iterationHorizon]$ and the \textit{output matrix} $\outputMatrix := \begin{bmatrix}
    I_\varDim & O_{\varDim + \paramDim}
\end{bmatrix} \in \R^{\varDim \times (\varDim + \paramDim)}$, and derive from \eqref{Eqn: Affine Output Feedback Map} that:
{
\setlength{\abovedisplayskip}{3.5pt}
\setlength{\belowdisplayskip}{3.5pt}
\begin{align}
    \label{Eqn: Control Error as affine function of State Errors}
    \ratePGDError_k 
    % &= \sum_{j=0}^k \feedbackGain_{k,j} \varError_{k-j} 
    &= \textstyle\sum_{j=0}^k \feedbackGain_{k,j} \outputMatrix \stateError_{k-j},
    \hspace{2mm} \forall \ k \in [0, \iterationHorizon].
\end{align}
}


% \antoine{Something is missing here. What we want is \textbf{tighter tube}. So why do we calculate things with SLS? SLS is a scalable tube for tube ``volume'' minimization! I would explicitly write the cost function as in the fast-SLS paper, say that optimizing over the $\Phi$ variables allow for tube volume minimization, especially the final one.}
% \frank{What is the \say{final one}?
% % Also, since $\Phi$ is only mentioned later in the paper, does this comment pertain to some content later on?
% Also, wouldn't the optimization problem over $\systemResponse$ only make sense after $\systemResponse$ is concretely, mathematically defined, and its connection to $\feedbackGainStacked$ made clear?
% This is why the NLP \eqref{Eqn: NLP for Tightest Tube Size} appears later on in the manuscript instead of here.
% }
% ~\\
% \frank{Consensus: Add a verbal description here and add pointers to specifics in later sections / theorems.
% }
% ~\\
% \frank{Added a preamble to Sec. 4B explaining this.}
% ~\\


A key insight of the SLS framework is that the design of the output feedback gains $\setFeedbackGains$ can be recast as the design of a \textit{system response}, mapping disturbances $\disturbanceStacked$ directly to state and control errors.
% $(\stateErrorStacked, \ratePGDErrorStacked)$. 
For brevity, we stack the LTV system dynamics \eqref{Eqn: State Error LTV Dynamics} and \eqref{Eqn: Control Error as affine function of State Errors} across iterations $k \in [0, \iterationHorizon]$. 
Let $\disturbanceStacked := (\stateError_0, \disturbance_0, \cdots, \disturbance_{\iterationHorizon-1})$, and define 
% $\jacobianStateStacked \in \R^{(\varDim + \paramDim)(\iterationHorizon+1) \times (\varDim + \paramDim)(\iterationHorizon+1)}$, $\jacobianControlStacked \in \R^{(\varDim + \paramDim)(\iterationHorizon+1) \times (\iterationHorizon+1)}$, $\outputMatrixStacked \in \R^{\varDim (\iterationHorizon+1) \times (\varDim + \paramDim)(\iterationHorizon+1)}$,  
% $\feedbackGainStacked \in \R^{(\iterationHorizon+1) \times \varDim(\iterationHorizon+1)}$,
% and $\downshift \in \R^{(\varDim + \paramDim)(\iterationHorizon+1) \times (\varDim + \paramDim)(\iterationHorizon+1)}$ by:
$\jacobianStateStacked$, $\jacobianControlStacked$, $\outputMatrixStacked$, $\feedbackGainStacked$, and $\downshift$ by:
\footnote{For notational simplicity, the explicit dependence of each $\jacobianState$ and $\jacobianControl$ on the nominal state-control pair $(\state_k, \ratePGD_k)$ is omitted from mention in \eqref{Eqn: Stacking State Jacobians}-\eqref{Eqn: Stacking Control Jacobians}.
}
{
\setlength{\abovedisplayskip}{1pt}
\setlength{\belowdisplayskip}{3.5pt}
\begin{subequations}
\label{Eqn: Stacking Matrices}
\begin{align}
    \label{Eqn: Stacking State Jacobians}
    \jacobianStateStacked &:= \blockDiag\{ \jacobianState_0, \cdots, \jacobianState_{\iterationHorizon-1}, O_{\varDim + \paramDim} \}, \\
    \label{Eqn: Stacking Control Jacobians}
    \jacobianControlStacked &:= \blockDiag\{ \jacobianControl_0, \cdots, \jacobianControl_{\iterationHorizon-1}, \zero_{\varDim + \paramDim} \}, \\
    \label{Eqn: Stacking Output Matrices}
    \outputMatrixStacked &:= \blockDiag\{ \outputMatrix, \cdots, \outputMatrix, \outputMatrix \}, \\
    \label{Eqn: Stacking Feedback Gains}
    \feedbackGainStacked &:= \begin{bmatrix}
        \feedbackGain_{0, 0} & O & \cdots & O \\[-7pt]
        \feedbackGain_{1, 0} & \feedbackGain_{1, 1} & \ddots & O \\[-3pt]
        \vdots & \vdots & \ddots & \vdots \\[-7pt]
        \feedbackGain_{\iterationHorizon, 0} & \feedbackGain_{\iterationHorizon, 1} & \ddots & \feedbackGain_{\iterationHorizon, \iterationHorizon}
    \end{bmatrix}, \\
    \label{Eqn: Downshift Operator}
    \downshift &:= \begin{bmatrix}
        O_{(\varDim + \paramDim) \times (\varDim + \paramDim)\iterationHorizon} & O_{\varDim + \paramDim} \\
        I_{(\varDim + \paramDim)\iterationHorizon} & O_{(\varDim + \paramDim)\iterationHorizon \times (\varDim + \paramDim)} 
    \end{bmatrix}
\end{align}
\end{subequations}
}

We use \eqref{Eqn: Stacking Matrices} to rewrite \eqref{Eqn: State Error LTV Dynamics}-\eqref{Eqn: Control Error as affine function of State Errors} across $k \in [0, \iterationHorizon]$ as:%
{
\setlength{\abovedisplayskip}{3.5pt}
\setlength{\belowdisplayskip}{3.5pt}
\begin{subequations}
\label{Eqn: State and Control Errors, Stacked}
\begin{align}
    \label{Eqn: State Errors, Stacked}
    \stateErrorStacked &= \downshift \jacobianStateStacked  \stateErrorStacked + \downshift 
    \jacobianControlStacked \ratePGDErrorStacked + \disturbanceStacked, \\
    \label{Eqn: Control Errors, Stacked}
    \ratePGDErrorStacked &= \feedbackGainStacked \outputMatrixStacked \stateErrorStacked,
\end{align}
\end{subequations}
}
\hspace{-1mm}which can be further written as:
{
\setlength{\abovedisplayskip}{4pt}
\setlength{\belowdisplayskip}{4pt}
\begin{align}
    \label{Eqn: State and Control Errors, Stacked, written using System Responses}
    \begin{bmatrix}
        \stateErrorStacked \\ \ratePGDErrorStacked
    \end{bmatrix}
    &= \begin{bmatrix}
        \systemResponse_{\state \disturbance} & \systemResponse_{\state \outputNoise} \\
        \systemResponse_{\ratePGD \disturbance} & \systemResponse_{\ratePGD \outputNoise}
    \end{bmatrix}
    \begin{bmatrix}
        \disturbanceStacked \\ \zero_{\varDim}
    \end{bmatrix} := \systemResponse \begin{bmatrix}
        \disturbanceStacked \\ \zero_{\varDim}
    \end{bmatrix},
\end{align}
}
\hspace{-1.5mm}where the constituent blocks of the \textit{system response} $\systemResponse \in \R^{(\varDim + \paramDim + 1)(\iterationHorizon+1) \times (2\varDim + \paramDim)(\iterationHorizon+1)}$ are:
{
\setlength{\abovedisplayskip}{3.5pt}
\setlength{\belowdisplayskip}{3.5pt}
\begin{subequations} \label{Eqn: System Responses from Output Feedback Gains}
\begin{align} 
    \label{Eqn: System Response from Feedback Gains (State, Disturbance)}
    \systemResponse_{\state \disturbance} &:= (I_{(\varDim + \paramDim)\iterationHorizon} - \downshift \jacobianStateStacked - \downshift \jacobianControlStacked \feedbackGainStacked \outputMatrixStacked)^{-1} \\ \label{Eqn: System Response from Feedback Gains (State, Output Noise)}
    \systemResponse_{\state \outputNoise} &:= (I_{(\varDim + \paramDim)\iterationHorizon} - \downshift \jacobianStateStacked - \downshift \jacobianControlStacked \feedbackGainStacked \outputMatrixStacked)^{-1} \downshift \jacobianControlStacked \feedbackGainStacked \\ 
    \label{Eqn: System Response from Feedback Gains (Input, Disturbance)}
    \systemResponse_{\ratePGD \disturbance} &:= \feedbackGainStacked \outputMatrixStacked (I - \downshift \jacobianStateStacked - \downshift \jacobianControlStacked \feedbackGainStacked \outputMatrixStacked)^{-1} \\ 
    \label{Eqn: System Response from Feedback Gains (Input, Output Noise)}
    \systemResponse_{\ratePGD \outputNoise} &:= \feedbackGainStacked \outputMatrixStacked (I - \downshift \jacobianStateStacked - \downshift \jacobianControlStacked \feedbackGainStacked \outputMatrixStacked)^{-1} \downshift \jacobianControlStacked \feedbackGainStacked + \feedbackGainStacked.
\end{align}
\end{subequations}
}
\hspace{-1mm}Conversely, from \eqref{Eqn: System Responses from Output Feedback Gains}, we also obtain that:
{
\setlength{\abovedisplayskip}{3.5pt}
\setlength{\belowdisplayskip}{3.5pt}
\begin{align} \label{Eqn: Feedback Error Gains from System Response}
    \feedbackGainStacked &= \systemResponse_{\ratePGD \outputNoise} - \systemResponse_{\ratePGD \disturbance} \systemResponse_{\state \disturbance}^{-1} \systemResponse_{\state \outputNoise}.
\end{align}
}

As proved in \cite{anderson2019levelsynthesis, ZhouTzoumas2023SafeControlofPartiallyObservedLTVSystems} and described below, the system response $\systemResponse$ parameterizes all causal output feedback gains $\feedbackGainStacked$.
% \antoine{I would personally cite Anderson2019 here, as the contribution of \cite{ZhouTzoumas2023SafeControlofPartiallyObservedLTVSystems} isn't really the output feedback SLP. The LTV form of the SLP takes 3 extra lines of proof compared to the one in Anderson2019. But I guess this is a personal preference.}

\begin{proposition}\cite[Prop. 1]{ZhouTzoumas2023SafeControlofPartiallyObservedLTVSystems}
If there exist output feedback gains $\feedbackGainStacked$ of the form \eqref{Eqn: Stacking Feedback Gains} satisfying \eqref{Eqn: State and Control Errors, Stacked}, then $\systemResponse_{\state \disturbance}$, $ \systemResponse_{\state \outputNoise}$, $ \systemResponse_{\ratePGD \disturbance}$, and $\systemResponse_{\ratePGD \outputNoise}$, as computed by \eqref{Eqn: System Responses from Output Feedback Gains},
% , are block lower triangular and 
satisfy: 
% the affine equalities:
{
\setlength{\abovedisplayskip}{3.5pt}
\setlength{\belowdisplayskip}{3.5pt}
\begin{subequations} \label{Eqn: System Responses, Affine Constraints}
\begin{align} \label{Eqn: System Responses, Affine Constraints, Controllability}
    \begin{bmatrix}
        I - \downshift \A & - \downshift \B
    \end{bmatrix} \systemResponse &= 
    \begin{bmatrix}
        I & O
    \end{bmatrix}, \\ \label{Eqn: System Responses, Affine Constraints, Observability}
    \systemResponse \begin{bmatrix}
        I - \downshift \jacobianStateStacked \\ -\outputMatrixStacked
    \end{bmatrix} &= 
    \begin{bmatrix}
        I \\ O
    \end{bmatrix}, \\ \label{Eqn: System Responses, Affine Constraints, Upper Triangular}
    \systemResponse_{\state \disturbance}, \systemResponse_{\state \outputNoise}, \systemResponse_{\ratePGD \disturbance}, \systemResponse_{\ratePGD \outputNoise} &\text{ are lower block-triangular}.
\end{align}
\end{subequations}
}
% \antoine{indices of the $\Phi$ matrices. greek or latin?}
% \frank{These are typos; thanks for pointing them out. I've fixed them.}
\hspace{-1mm}Conversely, if $\systemResponse_{\state \disturbance}$, $ \systemResponse_{\state \outputNoise}$, $ \systemResponse_{\ratePGD \disturbance}$, and $\systemResponse_{\ratePGD \outputNoise}$ 
% are block lower-triangular and 
satisfy \eqref{Eqn: System Responses, Affine Constraints}, then $\feedbackGainStacked$, as computed by \eqref{Eqn: Feedback Error Gains from System Response}, is lower block triangular.
\end{proposition}

Below, given $\systemResponse$ satisfying \eqref{Eqn: System Responses, Affine Constraints}, for each $k \in [0, \iterationHorizon]$ and $j \in [0, k]$, we respectively define $(\systemResponse_{\state \disturbance})_{k,j}$ and $(\systemResponse_{\ratePGD \disturbance})_{k,j}$ to be the $(k, j)$-th block
% \footnote{in the sense of \eqref{Eqn: Stacking Feedback Gains}} 
of $\systemResponse_{\state \disturbance}$ and $\systemResponse_{\ratePGD \disturbance}$ in the sense of \eqref{Eqn: Stacking Feedback Gains}.
% , and we define $(\systemResponse_{\ratePGD \disturbance})_{k,j}$ 


% ~\\
% \frank{To Antoine: If there is some way to simplify the above discussion, specifically, to avoid using $\systemResponse_{\state \outputNoise}$ and $\systemResponse_{\ratePGD \outputNoise}$, please feel free to edit. Since there is no output noise in our problem setting, my intuition suggests that those two sub-blocks of $\systemResponse$ may be unnecessary, but I might be wrong.
% }
% \antoine{I agree; things should simplify.. I don't see immediately what though}
% ~\\
% \frank{I see; I will keep that part of the content as it is for now then.}
% \antoine{I jusz checked and I think there is no possibility to write it more simple. Since C is not identity, we need the full parametrization to account for the lack of information}
% ~\\
% \frank{I see, then I will definitely leave as is for now.}


% \antoine{Is there any result in the literature that says that using a linear controller in this context is good enough? I guess we could argue that a fixed step is no controller, so we are at least better than that.}
% \frank{I know of no such result. I thought a main argument in the robust SLS paper is that optimizing over nonlinear feedback maps was intractable so we instead decouple that problem into the problems of designing a nominal state-control trajectory and designing a causal error feedback map.
% And yes, the literature usually uses a fixed timestep, because usually $\param$ is a fixed, perfectly known value.
% }
% \frank{Empirically it's worse with a fixed step size.}
% ~\\
% \frank{I've added a footnote}


\subsection{Linearization Error Bounds}
\label{subsec: Linearization Error Bounds}

% \antoine{In this section, the definition of the ``disturbance'' changes, and we need to make that really explicit. It now includs the ``smoothing'' error.}
% \frank{In Sec. \ref{subsec: Linearization Error Bounds}, when the PGD dynamics $\dynamicsPGD$ is assumed to be smooth by itself, disturbance is defined via \eqref{Eqn: Disturbance Bound, Stacked} as only the linearization error.  
% In Sec. \ref{subsec: Smoothing PGD Dynamics when Non-Differentiable}, disturbance includes the linearization error on top of the smoothing error, as emphasized in \eqref{Eqn: Disturbance Bound, Stacked, with Smoothed Dynamics}. I've added a sentence above \eqref{Eqn: Disturbance Bound, Stacked, with Smoothed Dynamics} to emphasize the difference between these two definitions.}


% \subsection{Linearization Error Bounds to Ensure Robust Steplength Constraint Satisfaction}
% \label{subsec: Linearization Error Bounds to Ensure Robust Steplength Constraint Satisfaction}


To compute forward reachable sets for the dynamics~\eqref{Eqn: True Dynamics} while enforcing the steplength constraints \eqref{Eqn: Steplength bound constraints}, we first bound the disturbance \eqref{Eqn: Disturbance} to characterize the error accrued from linearizing $\dynamicsFullState$. 
In Prop. \ref{Prop: Lipschitz Continuity of Gradient of PGD Dynamics},
% We 
we
begin by bounding the curvature of $\dynamicsFullState$, 
in Prop. \ref{Prop: Lipschitz Continuity of Gradient of PGD Dynamics}, 
using the Lipschitz constant of its gradient over all possible iterate rollouts.
% \antoine{For an uncertainty set that contains only a single point, do we recover something really close to the optimal solution? Or there is some uncertainty remaining?}
% \frank{Try to emphasize that when $\paramSet$ is just a singleton, the PGD dynamics just converges, per classical optimization theory.}
% ~\\
% \frank{I've added a new remark (Remark \ref{Remark: Singleton Parameter Set yields Zero Error}) to address this.
% }
% ~\\
To this end, define:
{
\setlength{\abovedisplayskip}{3.5pt}
\setlength{\belowdisplayskip}{3.5pt}
\begin{align} \label{Eqn: Convex Hull of all Trajectories}
    \varConvexHull &:= \conv\big\{ \varForwardReachableSet{k}(\var_0, \ratePGD_{0:k-1}): \\ \nonumber
    &\hspace{2cm} k \geq 0, \var_0 \in \varInitialSet, \ratePGD_{0:k-1} \in [\ratePGDMin, \ratePGDMax]^k \big\}
\end{align}
}
\hspace{-1.7mm}as the convex hull of all possible rollouts of the PGD dynamics \eqref{Eqn: Projected Gradient Descent (PGD)} from any $\var_0 \in \varInitialSet$.
We prove that, for a suitable $\constraintSet$, the gradient of \eqref{Eqn: True Dynamics} is Lipschitz continuous over $\varConvexHull \times \paramSet$.

% \antoine{In this proposition and in the next, the assumption $\constraintSet$ is $\R^\varDim$ or an affine subspace of $\R^\varDim$ is quite restrictive. I would mention that much earlier in the paper, maybe even in the abstract, because the method sounds really general, but actually our analysis works out only in this simplified setup. Otherwise, we may oversell the scope.}
% \frank{My opinion is the opposite. Whereas the numbered assumptions, i.e., Assumptions \ref{Assump: Objective Properties}, \ref{Assump: Constraint Set is Convex and Closed}, etc., are invoked throughout the paper, the condition that $\constraintSet$ is $\R^n$ or an affine subspace of $\R^n$ is only enforced within Sec. \ref{subsec: LTV Error System for Smooth PGD Dynamics} and \ref{subsec: Linearization Error Bounds}; this condition is then lifted in Sec. \ref{subsec: Smoothing PGD Dynamics when Non-Differentiable} (the main purpose of that subsection is, indeed to introduce the sampling-based smoothing technique that allows this condition to be removed). 
% Thus, I think adding this condition as an explicit assumption may mislead a reader into believing that all theoretical contributions throughout our manuscript only applies to the setting where this condition holds.
% }
% ~\\
% \frank{There is a separate but related question about whether or not we should clarify earlier in the paper that we have structured the presentation of our theory such that the theory for the \say{$\constraintSet$ is $\R^\varDim$ or an affine subspace of $\R^\varDim$} setting is developed first, and then generalized. I believe the second and third bullet points of our contribution list, in the introduction, with the italicized text, already makes this structure (i.e., order of presenting ideas in the manuscript) clear. However, we can discuss whether we want to clarify this structure earlier, in the abstract.}
% ~\\

\begin{proposition}
(\textbf{$\dynamicsFullState$ Admits Lipschitz Gradients})
\label{Prop: Lipschitz Continuity of Gradient of PGD Dynamics}
Suppose Assumptions \ref{Assump: Objective Properties}, \ref{Assump: Parameter Set is Convex and Compact}, and \ref{Assump: PGD Dynamics, Initial Iterate} hold, and $\constraintSet$ is $\R^\varDim$ or an affine subspace of $\R^\varDim$.
For each $i \in [\varDim + \paramDim]$, let $\dynamicsFullState_i: \R^{\varDim + \paramDim + 1} \ra \R$ denote the $i$-th component of the dynamics $\dynamicsFullState$ in \eqref{Eqn: True Dynamics}. 
Then, $\forall \ i \in [\varDim + \paramDim]$, there exists some $\hessianBound_i > 0$ such that, $\forall \ (\state, \ratePGD), (\state', \ratePGD') \in \varConvexHull \times \paramSet \times [\ratePGDMin, \ratePGDMax]$:
\begin{align}
    \label{Eqn: Lipschitz Continuity of Gradient of PGD Dynamics}
    &\Vert \nabla \dynamicsFullState_i(\state, \ratePGD) - \nabla \dynamicsFullState_i(\state', \ratePGD') \Vert_1 \hspace{-0.4mm} \leq \hspace{-0.4mm} \hessianBound_i \Vert (\state, \ratePGD) - (\state', \ratePGD') \Vert_\infty.
\end{align}
\end{proposition}

\begin{proof}
% To establish \eqref{Eqn: Lipschitz Continuity of Gradient of PGD Dynamics}, we will prove that 
% % $\derivativeFull f(\state, \ratePGD)$, i.e., the total derivative of $\dynamicsFullState$, 
% each $\nabla \dynamicsFullState_i(\cdot)$
% is Lipschitz continuous in $(\state, \ratePGD)$ over $\varConvexHull \times \paramSet \times [\ratePGDMin, \ratePGDMax]$.
We first verify that $\varConvexHull$ is bounded.
From \eqref{Eqn: Forward Reachable Sets approach Minimizer Sets} in Thm. \ref{Thm: Over-approximating Minimizer Set using Forward Reachable Sets},
% By the convergence properties of the PGD dynamics established in Prop. \ref{Prop: Convergence of PGD},
$\forall \ \var_0 \in \varInitialSet$, $\param \in \paramSet$, and $\ratePGD_{0:k-1} \in [\ratePGDMin, \ratePGDMax]^k$, we have 
$\Vert \varMap_k(\var_0, \param, \ratePGD_{0:k-1}) - \varOpt(\param) \Vert_2 \leq \convergenceRadius_k(\var_0)$ as $k \ra \infty$. Maximizing
% of this convergence result 
over $\param \in \paramSet$, we find $\dist(\varForwardReachableSet{k}(\var_0, \ratePGD_{0:k-1}), \varOptSet) \leq \convergenceRadius_k(\var_0)$ as $k \ra \infty$.
Since $\varForwardReachableSet{0}(\var_0, \param) = \varInitialSet$ and $\varOpt$ are bounded (by Assumption \ref{Assump: PGD Dynamics, Initial Iterate} and Prop. \ref{Prop: Minimizer Set is Bounded}, resp.), and $\convergenceRadius_k(\var_0) \geq 0$ decreases 
% exponentially
to 0 as $k \ra \infty$, we conclude that the sets $\varForwardReachableSet{k}(\var_0, \ratePGD_{0:k-1})$ are uniformly bounded across all $k \geq 0$, $\var_0 \in \varInitialSet$, and $\ratePGD_{0:k-1} \in [\ratePGDMin, \ratePGDMax]^k$, and thus $\varConvexHull$ is bounded.

Next, since $\constraintSet$ is $\R^\varDim$ or an affine subspace of $\R^\varDim$, the projection map $\proj_\constraintSet$ is affine. 
Thus, from \eqref{Eqn: Projected Gradient Descent (PGD)} and \eqref{Eqn: True Dynamics}, 
% the total derivative of the dynamics $\dynamicsFullState$, i.e., $\derivativeFull f$, 
each $\nabla \dynamicsFullState_i$
is bilinear in $\ratePGD$ and the first-order and second-order partial derivatives of $\obj$.
For $\nabla \dynamicsFullState_i$ 
% $\derivativeFull f$ 
to be Lipschitz continuous over $\varConvexHull \times \paramSet \times [\ratePGDMin, \ratePGDMax]$, it thus suffices to verify that the second-order partial derivatives of $\obj$ are Lipschitz continuous and bounded over $\varConvexHull \times \paramSet$, which follows from Assumption \ref{Assump: Objective Properties} and \ref{Assump: Parameter Set is Convex and Compact} and the boundedness of $\varConvexHull$ established above.
% This follows from Assumption \ref{Assump: Objective Properties}, which asserts that all second (hence first as well) partial derivatives of $\obj$ are Lipschitz continuous; boundedness then follows from the boundedness of $\varConvexHull \times \paramSet$.
% for any fixed $\ratePGD$, all partial derivatives of the map $(\var, \param) \mapsto \dynamicsFullState(\var, \param, \ratePGD)$ are Lipschitz continuous.
% Moreover, from Assumption \ref{Assump: Objective Properties}, all second partial derivatives of $\obj$ exist and are Lipschitz continuous.
\end{proof}

Equipped with curvature bounds of $\dynamicsFullState$ provided by Prop. \ref{Prop: Lipschitz Continuity of Gradient of PGD Dynamics}, we now bound the disturbance $\disturbance_k = \linearizationError(\state_k, \ratePGD_k, \stateNominal_k, \ratePGDNominal_k)$ in \eqref{Eqn: Disturbance}, in a manner analogous to \cite[Prop. III.1]{Leeman2025RobustNonlinearOptimalControlviaSLS}.

\begin{proposition}[\textbf{Linearization Error Bound}]
\label{Prop: Quadratic Bound on Linearization Error}
\looseness-1Suppose Assumptions \ref{Assump: Objective Properties}, \ref{Assump: Parameter Set is Convex and Compact}, and \ref{Assump: PGD Dynamics, Initial Iterate} hold, and $\constraintSet$ is $\R^\varDim$ or an affine subspace of $\R^\varDim$.
For each $k \in [0, \iterationHorizon]$ and $i \in [\varDim + \paramDim]$, we define $r_{k,i}: \R^{2(\varDim + \paramDim)} \ra \R$ to be the $i$-th coordinate of the linearization error map $r_k$,
% $r_k: \R^{2(\varDim + \paramDim)} \ra \R^{\varDim + \paramDim}$,
for each $i \in [\varDim + \paramDim]$.
Then for any $\var_0 \in  \varInitialSet$, $\ratePGD_{0:k-1}, \ratePGDNominal_{0:k-1} \in [\ratePGDMin, \ratePGDMax]^k$, any $\state_k = (\var_k, \param_k) \in \varForwardReachableSet{k}(\var_0, \ratePGD_{0:k-1}) \times \paramSet$ and $\stateNominal_k = (\varNominal_k, \paramNominal_k) \in \varForwardReachableSet{k}(\var_0, \ratePGDNominal_{0:k-1}) \times \paramSet$, and any $i \in [\varDim + \paramDim]$, we have:
\begin{align}
\label{Eqn: Quadratic Bound on Linearization Error}
    |\linearizationError_{k,i}(\state_k, \ratePGD_k, \stateNominal_k, \ratePGDNominal_k)| &\leq \hessianBound_i \cdot \Vert (\stateError_k, \ratePGDError_k) \Vert_\infty^2.
\end{align}
\end{proposition}

\begin{proof}
\looseness-1 By the Mean Value Theorem, there exists a $\lambda_i \in [0, 1]$ with:
{
\setlength{\abovedisplayskip}{3.5pt}
\setlength{\belowdisplayskip}{3.5pt}
\begin{align} \nonumber
    &\dynamicsFullState_i(\state_k, \ratePGD_k) - \dynamicsFullState_i(\stateNominal_k, \ratePGDNominal_k) \\ \nonumber
    = \ &\nabla \dynamicsFullState_i\big( \lambda_i(\state_k, \ratePGD_k) + (1 - \lambda_i)(\stateNominal_k, \ratePGDNominal_k) \big)^\top (\stateError_k, \ratePGDError_k).
\end{align}
}
\hspace{-1mm}By construction, $(\state_k, \ratePGD_k), (\stateNominal_k, \ratePGDNominal_k) \in \varConvexHull \times \paramSet \times [\ratePGDMin, \ratePGDMax]$, and thus 
$\lambda_i(\state_k, \ratePGD_k) + (1 - \lambda_i)(\stateNominal_k, \ratePGDNominal_k) \in \varConvexHull \times \paramSet \times [\ratePGDMin, \ratePGDMax]$ (Note that $\varConvexHull$ and $[\ratePGDMin, \ratePGDMax]$ are convex by construction, and $\paramSet$ is convex under Assumption \ref{Assump: Parameter Set is Convex and Compact}).
% any convex combination of $(\state_k, \ratePGD_k)$ and $ (\stateNominal_k, \ratePGDNominal_k)$.
Thus, we have:
{
\setlength{\abovedisplayskip}{3.5pt}
\setlength{\belowdisplayskip}{3.5pt}
\begin{align} \nonumber
    &|\linearizationError_{k,i}(\state_k, \ratePGD_k, \stateNominal_k, \ratePGDNominal_k)| \\ \nonumber
    = \ &|\dynamicsFullState_i(\state_k, \ratePGD_k) - \dynamicsFullState_i(\stateNominal_k, \ratePGDNominal_k) - \nabla \dynamicsFullState_i(\stateNominal_k, \ratePGDNominal_k)^\top (\stateError_k, \ratePGDError_k) | \\ \nonumber
    \leq \ &\Vert \nabla \dynamicsFullState_i\big( \lambda_i(\state_k, \ratePGD_k) + (1 - \lambda_i)(\stateNominal_k, \ratePGDNominal_k) \big) - \nabla \dynamicsFullState_i(\stateNominal_k, \ratePGDNominal_k) \Vert_1 \\ \nonumber
    &\hspace{1cm} \cdot \Vert (\stateError_k, \ratePGDError_k) \Vert_\infty \\ 
    \label{Eqn: Apply Lipschitz bounds of Dynamics' Gradient}
    \leq \ &\mu_i \cdot \lambda_i \Vert (\stateError_k, \ratePGDError_k) \Vert_\infty \cdot \Vert (\stateError_k, \ratePGDError_k) \Vert_\infty,
    \\ \nonumber
    \leq \ &\hessianBound_i \cdot \Vert (\stateError_k, \ratePGDError_k) \Vert_\infty^2.
\end{align}
}
where \eqref{Eqn: Apply Lipschitz bounds of Dynamics' Gradient} follows from Prop. \ref{Prop: Lipschitz Continuity of Gradient of PGD Dynamics}.
\end{proof}

We define $\hessianBoundStacked := \diag\{\hessianBound_1, \cdots, \hessianBound_{\varDim + \paramDim} \} \in \R^{(\varDim + \paramDim) \times (\varDim + \paramDim)}$, which allows us to stack \eqref{Eqn: Quadratic Bound on Linearization Error} across all $i \in [\varDim + \paramDim]$ as:
{
\setlength{\abovedisplayskip}{3.5pt}
\setlength{\belowdisplayskip}{3.5pt}
\begin{align}
\label{Eqn: Disturbance Bound, Stacked}
    \disturbance_k = \linearizationError_k(\state_k, \ratePGD_k, \stateNominal_k, \ratePGDNominal_k) \in \Vert (\stateError_k, \ratePGDError_k) \Vert_\infty^2 \hessianBoundStacked \ballUnit_\infty^{\varDim + \paramDim}
\end{align}
}

Below, in Thm. \ref{Thm: Robust Steplength Constraint Satisfaction}, we leverage Prop. \ref{Prop: Quadratic Bound on Linearization Error} to establish robust guarantees for satisfying the steplength constraint \eqref{Eqn: Steplength bound constraints}.
% We emphasize that, by enforcing \eqref{Eqn: Steplength bound constraints}, we ensure that the Lipschitz bounds over the gradients of $\dynamicsFullState$

\begin{theorem}
[\textbf{Robust Steplength Constraint Satisfaction}]
\label{Thm: Robust Steplength Constraint Satisfaction}
Suppose Assumptions \ref{Assump: Objective Properties}, \ref{Assump: Constraint Set is Convex and Closed}, 
\ref{Assump: Parameter Set is Convex and Compact},
and \ref{Assump: PGD Dynamics, Initial Iterate} hold.
Given a nominal trajectory $(\stateNominalStacked, \ratePGDNominalStacked)$ of the dynamics $\dynamicsFullState$ in \eqref{Eqn: True Dynamics} and a system response $\systemResponse$ satisfying \eqref{Eqn: System Responses, Affine Constraints}, 
if there exist auxiliary error upper bounds $\tubeSizeStacked := (\tubeSize_0, \cdots, \tubeSize_\iterationHorizon)$ satisfying, $\forall k \in [0, \iterationHorizon]$:
{
\setlength{\abovedisplayskip}{3.5pt}
\setlength{\belowdisplayskip}{3.5pt}
\begin{align} 
\label{Eqn: Inequality for Robust Constraint Satisfaction}
    &\ratePGDBoundSlope_p^\top (\stateNominal_k, \ratePGDNominal_k) + \ratePGDBoundOffset_p + \left\Vert \ratePGDBoundSlope_p^\top
    \begin{bmatrix}
        (\systemResponse_{\state \disturbance})_{k, 0} \\
        (\systemResponse_{\ratePGD \disturbance})_{k, 0}
    \end{bmatrix}
    \right\Vert_1 \diam(\paramSet) \\ \nonumber
    &\hspace{5mm} + \sum_{j=1}^k \tubeSize_{j-1}^2 \left\Vert \ratePGDBoundSlope_p^\top
    \begin{bmatrix}
        (\systemResponse_{\state \disturbance})_{k, 0} \\
        (\systemResponse_{\ratePGD \disturbance})_{k, 0}
    \end{bmatrix}
    \hessianBoundStacked
    \right\Vert_1 \leq 0, \quad \forall \ p \in [2], \\ 
\label{Eqn: Inequality for Propagating Recursive Tube Bounds}
    &\left\Vert 
    \ratePGDBoundSlope_p^\top
    \begin{bmatrix}
        (\systemResponse_{\state \disturbance})_{k, 0} \\
        (\systemResponse_{\ratePGD \disturbance})_{k, 0}
    \end{bmatrix}
    \right\Vert_\infty \diam(\paramSet) \\ \nonumber
    &\hspace{5mm} + \sum_{j=1}^k \tubeSize_{j-1}^2 \left\Vert 
    \begin{bmatrix}
        (\systemResponse_{\state \disturbance})_{k, 0} \\
        (\systemResponse_{\ratePGD \disturbance})_{k, 0}
    \end{bmatrix}
    \hessianBoundStacked
    \right\Vert_\infty \leq \tubeSize_k,
\end{align}
}
% ~\\
% \frank{INCOMPLETE}
% ~\\
% then the true trajectory $(\stateStacked, \ratePGDStacked)$ satisfies \eqref{Eqn: Steplength bound constraints} 
\hspace{-1.2mm}then $\Vert (\stateError_k, \ratePGDError_k) \Vert_\infty \leq \tubeSize_k$ and $\ratePGD_k \in [\ratePGDMin, \ratePGDMax]$ for each $k \in [0, \iterationHorizon]$,
regardless of the ground truth value of $\param \in \paramSet$.
\end{theorem}

\begin{proof}
We provide a proof by induction.
At $k = 0$, we prove $\Vert (\stateError_0, \ratePGD_0) \Vert_\infty \leq \tubeSize_0$ using \eqref{Eqn: State and Control Errors, Stacked, written using System Responses}, as follows:
{
\setlength{\abovedisplayskip}{3.5pt}
\setlength{\belowdisplayskip}{3.5pt}
\begin{align} \nonumber
    \left\Vert 
    \begin{bmatrix}
        \stateError_0 \\
        \ratePGDError_0
    \end{bmatrix}
    \right\Vert_\infty &= \left\Vert 
    \begin{bmatrix}
        (\systemResponse_{\state \disturbance})_{0, 0} \\
        (\systemResponse_{\ratePGD \disturbance})_{0, 0}
    \end{bmatrix}
    \stateError_0
    \right\Vert_\infty \\
    \label{Eqn: Apply Inequality for Propagating Recursive Tube Bounds, when k is 0}
    &\leq \left\Vert 
    \begin{bmatrix}
        (\systemResponse_{\state \disturbance})_{0, 0} \\
        (\systemResponse_{\ratePGD \disturbance})_{0, 0}
    \end{bmatrix}
    \right\Vert_\infty 
    \Vert \stateError_0 \Vert_\infty
    % \nonumber
    % &\leq \left\Vert 
    % \begin{bmatrix}
    %     (\systemResponse_{\state \disturbance})_{0, 0} \\
    %     (\systemResponse_{\ratePGD \disturbance})_{0, 0}
    % \end{bmatrix}
    % \right\Vert_\infty \diam(\paramSet) \\ 
    \leq \tubeSize_0,
\end{align}
}
\hspace{-1mm}where \eqref{Eqn: Apply Inequality for Propagating Recursive Tube Bounds, when k is 0} follows by taking $k = 0$ in \eqref{Eqn: Inequality for Propagating Recursive Tube Bounds} and noting that $\Vert \stateError_0 \Vert_\infty = \Vert (0, \paramError_0) \Vert_\infty \leq \diam(\paramSet)$. 
Similarly, we verify robust constraint satisfaction at $k = 0$ as follows:
{
\setlength{\abovedisplayskip}{3.5pt}
\setlength{\belowdisplayskip}{3.5pt}
\begin{align} \nonumber
    &\ratePGDBoundSlope_p^\top (\state_0, \ratePGD_0) + \ratePGDBoundOffset_p \\ \nonumber
    = \ &\ratePGDBoundSlope_p^\top (\stateNominal_0, \ratePGDNominal_0) + \ratePGDBoundOffset_p + \ratePGDBoundSlope_p^\top (\stateError_0, \ratePGDError_0) \\ \label{Eqn: Apply Inequality for Robust Constraint Satisfaction, k = 0}
    \leq \ &\ratePGDBoundSlope_p^\top (\stateNominal_0, \ratePGDNominal_0) + \ratePGDBoundOffset_p + 
    \left\Vert 
    \ratePGDBoundSlope_p^\top
    \begin{bmatrix}
        (\systemResponse_{\state \disturbance})_{0, 0} \\
        (\systemResponse_{\ratePGD \disturbance})_{0, 0}
    \end{bmatrix}
    \stateError_0
    \right\Vert_\infty\\ 
    \leq \ &\ratePGDBoundSlope_p^\top (\stateNominal_0, \ratePGDNominal_0) + \ratePGDBoundOffset_p + 
    \left\Vert 
    \ratePGDBoundSlope_p^\top
    \begin{bmatrix}
        (\systemResponse_{\state \disturbance})_{0, 0} \\
        (\systemResponse_{\ratePGD \disturbance})_{0, 0}
    \end{bmatrix}
    \right\Vert_1
    \Vert \stateError_0 \Vert_\infty \leq \ 0,\nonumber
\end{align}
}
\hspace{-1mm}where \eqref{Eqn: Apply Inequality for Robust Constraint Satisfaction, k = 0} follows because $\Vert \stateError_0 \Vert_\infty \leq \diam(\paramSet)$.
From \eqref{Eqn: Inequality for Robust Constraint Satisfaction}, we have $\ratePGDBoundSlope_p^\top (\stateNominal_0, \ratePGDNominal_0) + \ratePGDBoundOffset_p \leq 0$.
Thus, $\ratePGD_0, \ratePGDNominal_0 \in [\ratePGDMin, \ratePGDMax]$. 

Suppose for some $k \geq 1$, we have $\Vert (\stateError_j, \ratePGDError_j) \Vert_\infty \leq \tubeSize_j$ for each $j \in [0, k-1]$ and $\ratePGD_{0:k-1}, \ratePGDNominal_{0:k-1} \in [\ratePGDMin, \ratePGDMax]^k$.
For the inductive step, we must prove that $\Vert (\stateError_k, \ratePGDError_k) \Vert_\infty \leq \tubeSize_k$ and $\ratePGD_k, \ratePGDNominal_k \in [\ratePGDMin, \ratePGDMax]$. 
First, we verify that the disturbance bounds prescribed by Prop. \ref{Prop: Quadratic Bound on Linearization Error} are valid for $\disturbance_j = \linearizationError_j(\state_j, \ratePGD_j, \stateNominal_j, \ratePGDNominal_j)$ across all $j \in [0, k-1]$.
% First, we verify that we can apply Prop. \ref{Prop: Quadratic Bound on Linearization Error} to bound the disturbance $\disturbance_j = \linearizationError_j(\state_j, \ratePGD_j, \stateNominal_j, \ratePGDNominal_j)$ for each $j \in [0, k-1]$. 
Indeed, for each $j \in [0, k-1]$, since $\param_j = \param_0 \in \paramSet$ and $\paramNominal_j = \paramNominal_0 \in \paramSet$:
{
\setlength{\abovedisplayskip}{3.5pt}
\setlength{\belowdisplayskip}{3.5pt}
\begin{subequations}
\begin{align}
    \state_j &= (\var_j, \param_j) \in \varForwardReachableSet{j}(\var_0, \ratePGD_{0:k-1}) \times \paramSet, \\
    \stateNominal_j &= (\varNominal_j, \paramNominal_j) \in \varForwardReachableSet{j}(\varNominal_0, \ratePGDNominal_{0:k-1}) \times \paramSet,
\end{align}
\end{subequations}
}
\hspace{-1mm}where $\ratePGD_{0:k-1}, \ratePGDNominal_{0:k-1} \in [\ratePGDMin, \ratePGDMax]^k$ by the induction hypothesis, and $\var_0, \varNominal_0 \in \varInitialSet$.
Thus, the conditions of Prop. \ref{Prop: Quadratic Bound on Linearization Error} are met at time $j$ for each $j \in [0, k-1]$, so 
\eqref{Eqn: Quadratic Bound on Linearization Error} and \eqref{Eqn: Disturbance Bound, Stacked} imply:
{
\setlength{\abovedisplayskip}{3.5pt}
\setlength{\belowdisplayskip}{3.5pt}
\begin{align}
\label{Eqn: Applying Disturbance Bound, Stacked}
    \disturbance_j 
    % = \linearizationError_j(\state_j, \ratePGD_j, \stateNominal_j, \ratePGDNominal_j)
    \in \Vert (\stateError_j, \ratePGDError_j) \Vert_\infty^2 \hessianBoundStacked \ballUnit_\infty^{\varDim + \paramDim} 
    \subseteq \tubeSize_j \hessianBoundStacked \ballUnit_\infty^{\varDim + \paramDim}
\end{align}
}
\hspace{-1mm}for each $j \in [0, k-1]$. 
Then, by \eqref{Eqn: Applying Disturbance Bound, Stacked}:
% $\Vert (\stateError_k, \paramError_k) \Vert_\infty \leq \tubeSize_k$: %, as follows:
{
\setlength{\abovedisplayskip}{3.5pt}
\setlength{\belowdisplayskip}{3.5pt}
\begin{align} \nonumber
    \left\Vert 
    \begin{bmatrix}
        \stateError_k \\
        \ratePGDError_k
    \end{bmatrix}
    \right\Vert_\infty &= \left\Vert 
    \begin{bmatrix}
        (\systemResponse_{\state \disturbance})_{k, 0} \\
        (\systemResponse_{\ratePGD \disturbance})_{k, 0}
    \end{bmatrix}
    \stateError_0
    + \sum_{j=1}^{k}
    \begin{bmatrix}
        (\systemResponse_{\state \disturbance})_{k, j} \\
        (\systemResponse_{\ratePGD \disturbance})_{k, j}
    \end{bmatrix}
    \disturbance_{j-1}
    \right\Vert_\infty \\ \label{Eqn: Apply Inequality for Propagating Recursive Tube Bounds, general k}
    &\leq \left\Vert 
    \begin{bmatrix}
        (\systemResponse_{\state \disturbance})_{k, 0} \\
        (\systemResponse_{\ratePGD \disturbance})_{k, 0}
    \end{bmatrix}
    \right\Vert_\infty
    \Vert \stateError_0 \Vert_\infty
    \\ \nonumber
    &\hspace{1cm} + \sum_{j=1}^{k}
    \tubeSize_{j-1}^2
    \left\Vert
    \begin{bmatrix}
        (\systemResponse_{\state \disturbance})_{k, j} \\
        (\systemResponse_{\ratePGD \disturbance})_{k, j}
    \end{bmatrix}
    \hessianBoundStacked
    \right\Vert_\infty \leq \tubeSize_k,
    % \\ \nonumber
    % &\leq \left\Vert 
    % \begin{bmatrix}
    %     (\systemResponse_{\state \disturbance})_{0, 0} \\
    %     (\systemResponse_{\ratePGD \disturbance})_{0, 0}
    % \end{bmatrix}
    % \right\Vert_\infty \diam(\paramSet) 
    % \\
    % &
\end{align}
}
\hspace{-1mm}We also prove robust constraint satisfaction, as follows:
{
\setlength{\abovedisplayskip}{4pt}
\setlength{\belowdisplayskip}{4pt}
\begin{align} \nonumber
    &\ratePGDBoundSlope_p^\top (\state_k, \ratePGD_k) + \ratePGDBoundOffset_p 
    = \ \ratePGDBoundSlope_p^\top (\stateNominal_k, \ratePGDNominal_k) + \ratePGDBoundOffset_p + \ratePGDBoundSlope_p^\top (\stateError_k, \ratePGDError_k) \\ \nonumber
    = \ &\ratePGDBoundSlope_p^\top (\stateNominal_k, \ratePGDNominal_k) + \ratePGDBoundOffset_p \\ \nonumber
    &\hspace{5mm} + 
    \left\Vert 
    \ratePGDBoundSlope_p^\top
    \begin{bmatrix}
        (\systemResponse_{\state \disturbance})_{k, 0} \\
        (\systemResponse_{\ratePGD \disturbance})_{k, 0}
    \end{bmatrix}
    \stateError_0
    + \sum_{j=1}^{k}
    \ratePGDBoundSlope_p^\top
    \begin{bmatrix}
        (\systemResponse_{\state \disturbance})_{k, j} \\
        (\systemResponse_{\ratePGD \disturbance})_{k, j}
    \end{bmatrix}
    \disturbance_{j-1}
    \right\Vert_\infty \\ \nonumber
    = \ &\ratePGDBoundSlope_p^\top (\stateNominal_k, \ratePGDNominal_k) + \ratePGDBoundOffset_p + 
    \left\Vert 
    \ratePGDBoundSlope_p^\top
    \begin{bmatrix}
        (\systemResponse_{\state \disturbance})_{k, 0} \\
        (\systemResponse_{\ratePGD \disturbance})_{k, 0}
    \end{bmatrix}
    \right\Vert_1
    \Vert \stateError_0 \Vert_\infty
    \\
    \label{Eqn: Apply Inequality for Robust Constraint Satisfaction, k geq 1}
    &\hspace{5mm} + \sum_{j=1}^{k} \tubeSize_{j-1}^2
    \left\Vert
    \ratePGDBoundSlope_p^\top
    \begin{bmatrix}
        (\systemResponse_{\state \disturbance})_{k, j} \\
        (\systemResponse_{\ratePGD \disturbance})_{k, j}
    \end{bmatrix}
    \hessianBoundStacked
    \right\Vert_1 
    \leq 0.
\end{align}
}
\hspace{-1mm}Finally, from \eqref{Eqn: Inequality for Robust Constraint Satisfaction}, we again directly obtain that $\ratePGDBoundSlope_p^\top (\stateNominal_k, \ratePGDNominal_k) + \ratePGDBoundOffset_p \leq 0$, thus completing the induction step.
\end{proof}

% ~\\
% ~\\
% \frank{Currently modifying the proof; old proof content below.}
% ~\\

% First, we prove that $\Vert (\stateError_k, \ratePGDError_k) \Vert_\infty \leq \tubeSize_k$ for each $k \in [0, \iterationHorizon]$. 
% When $k = 0$, from \eqref{Eqn: State and Control Errors, Stacked, written using System Responses}, we have:
% \begin{align} \nonumber
%     \left\Vert 
%     \begin{bmatrix}
%         \stateError_0 \\
%         \ratePGDError_0
%     \end{bmatrix}
%     \right\Vert_\infty &= \left\Vert 
%     \begin{bmatrix}
%         (\systemResponse_{\state \disturbance})_{0, 0} \\
%         (\systemResponse_{\ratePGD \disturbance})_{0, 0}
%     \end{bmatrix}
%     \stateError_0
%     \right\Vert_\infty \\ \nonumber
%     &\leq \left\Vert 
%     \begin{bmatrix}
%         (\systemResponse_{\state \disturbance})_{0, 0} \\
%         (\systemResponse_{\ratePGD \disturbance})_{0, 0}
%     \end{bmatrix}
%     \right\Vert_\infty 
%     \Vert \stateError_0 \Vert_\infty \\ 
%     % \nonumber
%     % &\leq \left\Vert 
%     % \begin{bmatrix}
%     %     (\systemResponse_{\state \disturbance})_{0, 0} \\
%     %     (\systemResponse_{\ratePGD \disturbance})_{0, 0}
%     % \end{bmatrix}
%     % \right\Vert_\infty \diam(\paramSet) \\ 
%     % \label{Eqn: Apply Inequality for Propagating Recursive Tube Bounds, when k is 0}
%     &\leq \tubeSize_0,
% \end{align}
% where \eqref{Eqn: Apply Inequality for Propagating Recursive Tube Bounds, when k is 0} by taking $k = 0$ in \eqref{Eqn: Inequality for Propagating Recursive Tube Bounds}, and noting that $\Vert \stateError_0 \Vert_\infty = \Vert (0, \paramError_0) \Vert_\infty \leq \diam(\paramSet)$.
% By induction, suppose that for some $k \in [1, \iterationHorizon]$, we have $\Vert (\stateError_j, \ratePGDError_j) \Vert_\infty \leq \tubeSize_j$ for each $j \in [0, k-1]$. 
% Then, for each $j \in [0, k-1]$:
% \begin{align} \nonumber
%     \disturbance_j = \linearizationError_j(\state_j, \ratePGD_j, \stateNominal_j, \ratePGDNominal_j) &\in \Vert (\stateError_j, \ratePGDError_j) \Vert_\infty^2 \hessianBoundStacked \ballUnit_\infty^{\varDim + \paramDim} \\ \nonumber
%     &\subseteq \tubeSize_j^2 \hessianBoundStacked \ballUnit_\infty^{\varDim + \paramDim},
% \end{align}
% so from \eqref{Eqn: State and Control Errors, Stacked, written using System Responses}, we obtain that:
% \begin{align} \nonumber
%     \left\Vert 
%     \begin{bmatrix}
%         \stateError_k \\
%         \ratePGDError_k
%     \end{bmatrix}
%     \right\Vert_\infty &= \left\Vert 
%     \begin{bmatrix}
%         (\systemResponse_{\state \disturbance})_{k, 0} \\
%         (\systemResponse_{\ratePGD \disturbance})_{k, 0}
%     \end{bmatrix}
%     \stateError_0
%     + \sum_{j=1}^{k}
%     \begin{bmatrix}
%         (\systemResponse_{\state \disturbance})_{k, j} \\
%         (\systemResponse_{\ratePGD \disturbance})_{k, j}
%     \end{bmatrix}
%     \disturbance_{j-1}
%     \right\Vert_\infty \\ \nonumber
%     &\leq \left\Vert 
%     \begin{bmatrix}
%         (\systemResponse_{\state \disturbance})_{k, 0} \\
%         (\systemResponse_{\ratePGD \disturbance})_{k, 0}
%     \end{bmatrix}
%     \right\Vert
%     \Vert \stateError_0 \Vert_2
%     \\ \nonumber
%     &\hspace{1cm} + \sum_{j=1}^{k}
%     \tubeSize_{j-1}^2
%     \left\Vert
%     \begin{bmatrix}
%         (\systemResponse_{\state \disturbance})_{k, j} \\
%         (\systemResponse_{\ratePGD \disturbance})_{k, j}
%     \end{bmatrix}
%     \hessianBoundStacked
%     \right\Vert_\infty 
%     % \\ \nonumber
%     % &\leq \left\Vert 
%     % \begin{bmatrix}
%     %     (\systemResponse_{\state \disturbance})_{0, 0} \\
%     %     (\systemResponse_{\ratePGD \disturbance})_{0, 0}
%     % \end{bmatrix}
%     % \right\Vert_\infty \diam(\paramSet) 
%     \\ 
%     % \label{Eqn: Apply Inequality for Propagating Recursive Tube Bounds, general k}
%     &\leq \tubeSize_k,
% \end{align}
% where \eqref{Eqn: Apply Inequality for Propagating Recursive Tube Bounds, general k} follows from \eqref{Eqn: Inequality for Propagating Recursive Tube Bounds} and the fact that $\Vert \stateError_0 \Vert_\infty \leq \diam(\paramSet)$. 
% We have thus proved that $\Vert (\stateError_k, \ratePGDError_k) \Vert_\infty \leq \tubeSize_k$ for each $k \in [0, \iterationHorizon]$. As noted above, we then have $\disturbance_k \in \tubeSize_k^2 \hessianBoundStacked \ballUnit_\infty^{\varDim + \paramDim}$ for each $k \in [0, \iterationHorizon]$.

% We now prove robust constraint satisfaction, i.e., that $\ratePGDBoundSlope_p^\top (\state_k, \ratePGD_k) + \ratePGDBoundOffset \leq 0$ for each $k \in [\iterationHorizon]$, as follows:
% \begin{align} \nonumber
%     &\ratePGDBoundSlope_p^\top (\state_k, \ratePGD_k) + \ratePGDBoundOffset_p \\ \nonumber
%     = \ &\ratePGDBoundSlope_p^\top (\stateNominal_k, \ratePGDNominal_k) + \ratePGDBoundOffset_p + \ratePGDBoundSlope_p^\top (\stateError_k, \ratePGDError_k) \\ \nonumber
%     = \ &\ratePGDBoundSlope_p^\top (\stateNominal_k, \ratePGDNominal_k) + \ratePGDBoundOffset_p \\ \nonumber
%     &\hspace{5mm} + 
%     \left\Vert 
%     \ratePGDBoundSlope_p^\top
%     \begin{bmatrix}
%         (\systemResponse_{\state \disturbance})_{k, 0} \\
%         (\systemResponse_{\ratePGD \disturbance})_{k, 0}
%     \end{bmatrix}
%     \stateError_0
%     + \sum_{j=1}^{k}
%     \ratePGDBoundSlope_p^\top
%     \begin{bmatrix}
%         (\systemResponse_{\state \disturbance})_{k, j} \\
%         (\systemResponse_{\ratePGD \disturbance})_{k, j}
%     \end{bmatrix}
%     \disturbance_{j-1}
%     \right\Vert_\infty \\ \nonumber
%     = \ &\ratePGDBoundSlope_p^\top (\stateNominal_k, \ratePGDNominal_k) + \ratePGDBoundOffset_p \\ \nonumber
%     &\hspace{5mm} + 
%     \left\Vert 
%     \ratePGDBoundSlope_p^\top
%     \begin{bmatrix}
%         (\systemResponse_{\state \disturbance})_{k, 0} \\
%         (\systemResponse_{\ratePGD \disturbance})_{k, 0}
%     \end{bmatrix}
%     \right\Vert_1
%     \Vert \stateError_0 \Vert_\infty
%     \\ \nonumber
%     &\hspace{5mm} + \sum_{j=1}^{k} \tubeSize_{j-1}^2
%     \left\Vert
%     \ratePGDBoundSlope_p^\top
%     \begin{bmatrix}
%         (\systemResponse_{\state \disturbance})_{k, j} \\
%         (\systemResponse_{\ratePGD \disturbance})_{k, j}
%     \end{bmatrix}
%     \hessianBoundStacked
%     \right\Vert_1 \\ 
%     % \label{Eqn: Apply Inequality for Robust Constraint Satisfaction}
%     \leq \ &0,
% \end{align}
% where \eqref{Eqn: Apply Inequality for Robust Constraint Satisfaction} follows from \eqref{Eqn: Inequality for Robust Constraint Satisfaction} and $\Vert \stateError_0 \Vert_\infty \leq \diam(\paramSet)$.

\looseness=-1
We now present a nonlinear program (NLP) to recover the tightest reachable set for the PGD dynamics \eqref{Eqn: Projected Gradient Descent (PGD)}, as characterized by $\tubeSize_\iterationHorizon$ at the end of the iteration horizon $[0, \iterationHorizon]$, while ensuring that the steplength constraints \eqref{Eqn: Steplength bound constraints} are robustly satisfied\footnote{thus enabling the use of the reachable sets of PGD dynamics for constructing valid outer approximations of $\varOptSet$ via Thm. \ref{Thm: Over-approximating Minimizer Set using Forward Reachable Sets}}. The convex cost $\regularizerSystemResponse(\systemResponse)$ is added as a regularizer, to minimize the uncertainty propagation \cite{Leeman2024FastSLS}: 
% ():
{
\setlength{\abovedisplayskip}{3.5pt}
\setlength{\belowdisplayskip}{3.5pt}
\begin{subequations}
\label{Eqn: NLP for Tightest Tube Size}
\begin{align}
\label{Eqn: NLP for Tightest Tube Size, Objective}
    \min_{\stateNominalStacked, \ratePGDNominalStacked, \systemResponse, \tubeSizeStacked} \hspace{5mm} &\tubeSize_\iterationHorizon
    + \regularizerSystemResponse(\systemResponse)
    \\
\label{Eqn: NLP for Tightest Tube Size, Constraints}
    \text{s.t.} \hspace{5mm} &\eqref{Eqn: Nominal Dynamics}, \eqref{Eqn: System Responses, Affine Constraints}, \eqref{Eqn: Inequality for Robust Constraint Satisfaction}, \eqref{Eqn: Inequality for Propagating Recursive Tube Bounds}.
\end{align}
\end{subequations}
}
\hspace{-1.2mm}We tractably compute a feasible point of~\eqref{Eqn: NLP for Tightest Tube Size} using the output feedback SLS solver~\cite{leeman2026vision}, with 
% $\mathcal{O}(\log\iterationHorizon \log^2(\varDim + \paramDim))^3)$
$\mathcal{O}(\iterationHorizon(\varDim^3 + \paramDim^3))$
runtime. 
% prescribed by \cite{Leeman2024FastSLS}.

% \antoine{
% It may be useful to state explicitly that the tube bounds imply that all true and nominal trajectories remain inside the domain on which the Lipschitz and smoothing constants are valid.
% }

\subsection{Smoothing PGD Dynamics when Non-Differentiable}
\label{subsec: Smoothing PGD Dynamics when Non-Differentiable}

% \antoine{Just to make sure I understand
% Are you smoothing in both the state and control variables intentionally?
% Do you want smoothing over 
% ($\xi$,$\theta$,$\eta$) or only over the argument entering the projection?
% Since $\theta$ is a static hidden parameter, is it meaningful to smooth in the $\theta$ direction?
% }
% \frank{Only the $\var$-updates are smoothed; the $\param$, which is not really ever updated via $\param_{k+1} = \param_k$, is left alone.
% This is specified by \eqref{Eqn: Full State Dynamics, Smoothed}, and I've added a clarifying remark below it.
% }
% \brendan{Does this distinction really matter? The update to $\param$ is deterministic in the sample, and always the same thing. Even if we did ``smooth'' along that component, we'd get the same result as if we chose the realization of any particular sample, or even the center point.}
% \frank{As far as I can tell, it matters only in the sense that smoothing over higher dimensional variables / problems results in worse, i.e., larger Lipschitz constants. Qualitatively it doesn't matter though.}

If $\constraintSet$ is neither $\R^\varDim$ nor an affine subspace of $\R^\varDim$, $\proj_\constraintSet(\cdot)$ is continuous but in general non-differentiable \cite[Ex. 7.4-7.5]{WrightRecht2022OptimizationForDataAnalysis},
in which case
% In such cases,
$\dynamicsPGD$ and $\dynamicsFullState$ (see \eqref{Eqn: PGD Updates Notation}) are also, in general, non-differentiable. 
% \antoine{I don't remember what $\dynamicsFullState$ is at this point. Name it maybe?}
% \frank{I added a eqref link back to the equation.}
To 
% facilitate the construction of 
construct
the LTV error system and linearization error bounds in Sec. \ref{subsec: LTV Error System for Smooth PGD Dynamics}--\ref{subsec: Linearization Error Bounds}, we \textit{smooth} $\dynamicsPGD$ and $\dynamicsFullState$, as follows.
Given a \textit{sampling radius} $\samplingRadius > 0$, for 
any
% each state
$\state = (\var, \param) \in \R^{\varDim + \paramDim}$ and 
% steplength control
$\ratePGD \in \R$, define $\dynamicsPGDSmoothed(\cdot; \samplingRadius): \R^{\varDim + \paramDim + 1} \ra \R^{\varDim} $ and $\dynamicsFullStateSmoothed(\cdot; \samplingRadius): \R^{\varDim + \paramDim + 1} \ra \R^{\varDim + \paramDim} $ by:
{
\setlength{\abovedisplayskip}{3.5pt}
\setlength{\belowdisplayskip}{3.5pt}
\begin{align}
    \label{Eqn: True Dynamics, Smoothed}
    \dynamicsPGDSmoothed(\state, \ratePGD; \samplingRadius) &:= \E_{\samplingVariable \sim \unif(\ballUnit_2^{\varDim + \paramDim})} \big[ \dynamicsPGD\big( (\state, \ratePGD) + \samplingRadius \samplingVariable \big) \big], \\
    \label{Eqn: Full State Dynamics, Smoothed}
    \dynamicsFullStateSmoothed(\state, \ratePGD; \samplingRadius) &:= \big( \dynamicsPGDSmoothed(\state, \ratePGD; \samplingRadius), \param \big).
\end{align}
}
We impose a fixed $\samplingRadiusMax > 0$ such that $\samplingRadius \in (0, \samplingRadiusMax]$. 

We review properties of smoothed functions, as below \cite{Flaxman2005OnlineConvexOptimization}.
% \antoine{Highlight the smoothing tradeoff: bias versus improved regularity.}
% \frank{I've added a remark to this effect after the proposition.}
% \antoine{The proposition below has a name, while the others don't. Iwould try to have it uniform for consistency.}

\begin{proposition}[\textbf{Properties of Smoothed Functions}]
\label{Prop: Properties of Smoothed Functions}
Let $\exampleConstraintSet \subseteq \R^\exampleInputDim$, and suppose $\exampleFunction: \R^\exampleInputDim \ra \R^\exampleOutputDim$ is $\lipschitz_\exampleFunction$-Lipschitz continuous in the $\infty$-norm over $\exampleConstraintSet \oplus \samplingRadiusMax I_\exampleInputDim \ballUnit_2^\exampleInputDim$. Define $\exampleFunctionSmoothed(\cdot; \samplingRadius): \R^\exampleInputDim \ra \R^\exampleOutputDim$ as:
{
\setlength{\abovedisplayskip}{3.5pt}
\setlength{\belowdisplayskip}{3.5pt}
\begin{align}
    \exampleFunctionSmoothed(\cdot; \samplingRadius) &:= \E_{\samplingVariable \sim \unif(\ballUnit_2^\exampleInputDim)} \big[ \exampleFunction(\varExample + \samplingRadius \samplingVariable) \big], \hspace{5mm} \forall \ \varExample \in \R^\exampleInputDim.
\end{align}
}
\hspace{-1mm}Then, for each $i \in [\exampleOutputDim]$ and $\varExample, \varExample' \in \exampleConstraintSet$:
\begin{enumerate}
    \item $\Vert \exampleFunctionSmoothed(\varExample; \samplingRadius) - \exampleFunction(\varExample) \Vert_\infty \leq \lipschitz_\exampleFunction \samplingRadius$;
    % 
    \item $\nabla \exampleFunctionSmoothed_i(\varExample; \samplingRadius) = \frac{\exampleInputDim}{\samplingRadius} \E_{\samplingVariable \sim \unif(\ballUnit_2^\exampleInputDim)} \big[ \exampleFunction_i (\varExample + \samplingRadius \samplingVariable) \samplingVariable \big]$;
    % 
    \item $\Vert \nabla \exampleFunctionSmoothed_i(\varExample; \samplingRadius) - \nabla \exampleFunctionSmoothed_i(\varExample'; \samplingRadius) \Vert_1 \leq \frac{\exampleInputDim^{3/2}}{\samplingRadius} \lipschitz_\exampleFunction \cdot \Vert \varExample - \varExample' \Vert_\infty$.
\end{enumerate}
\end{proposition}

\begin{proof}
The above properties of smoothed functions are direct consequences of \cite[Lemma 1]{Flaxman2005OnlineConvexOptimization} and 
\cite[Sec. 4.2]{Bravo2018BanditLearninginConcaveNPersonGames}.
\end{proof}

% \begin{proof}
% \begin{enumerate}
%     \item We have:%
%     \vspace{-5mm}
%     \begin{adjustwidth}{-\leftmargin}{0pt}
%     \begin{align} \nonumber
%         \Vert \exampleFunctionSmoothed(\varExample; \samplingRadius) - \exampleFunction(\varExample) \Vert_\infty \leq \ &\E_{\samplingVariable \sim \unif(\ballUnit_2^\exampleInputDim)} \big[ \Vert \exampleFunction_i (\varExample + \samplingRadius \samplingVariable) - \exampleFunction_i (\varExample) \Vert \big] \\ \nonumber
%         \leq \ &\lipschitz_\exampleFunction \samplingRadius \cdot \E_{\samplingVariable \sim \unif(\ballUnit_2^\exampleInputDim)} \big[ \Vert \samplingVariable \Vert_\infty \big]  \leq \samplingRadius \lipschitz_\exampleFunction.
%     \end{align}
%     \end{adjustwidth}
%     % 
%     \item This result follows directly from \cite[Lemma 1]{Flaxman2005OnlineConvexOptimization}.
%     % 
%     \item From the above expression for $\nabla \exampleFunctionSmoothed_i(\varExample; \samplingRadius)$:
%     \begin{align} \nonumber
%         &\Vert \nabla \exampleFunctionSmoothed(\varExample; \samplingRadius) - \nabla \exampleFunctionSmoothed(\varExample; \samplingRadius) \Vert_1 \\ \nonumber
%         \leq \ &\frac{\exampleInputDim}{\samplingRadius} \E_{\samplingVariable \sim \unif(\ballUnit_2^\exampleInputDim)} \big[|\exampleFunction_i(\varExample + \samplingRadius \samplingVariable) - \exampleFunction_i(\varExample' + \samplingRadius \samplingVariable)| \cdot \Vert \samplingVariable \Vert_1 \big] \\ \nonumber
%         \leq \ &\frac{\exampleInputDim}{\samplingRadius} \lipschitz_\exampleFunction \Vert \varExample - \varExample' \Vert_\infty \cdot \sqrt{\exampleInputDim}.
%     \end{align}
% \end{enumerate}
% \end{proof}

\begin{remark}
% Prop. \ref{Prop: Properties of Smoothed Functions} indicates that selecting a larger sampling radius $\samplingRadius$ would produce a smooth proxy $\exampleFunctionSmoothed(\cdot; \samplingRadius)$ for $\exampleFunction$ that has less curvature, at the cost of deviating more significantly from $\exampleFunction$.
Prop. \ref{Prop: Properties of Smoothed Functions} indicates that selecting a smaller sampling radius $\samplingRadius$ would produce a smooth proxy $\exampleFunctionSmoothed(\cdot; \samplingRadius)$ that more closely approximates $\exampleFunction$, at the cost of having larger curvature and thus inducing larger linearization errors.
\end{remark}

From Prop. \ref{Prop: Properties of Smoothed Functions}, we obtain the following properties of the smoothed dynamics $\dynamicsFullStateSmoothed$ in \eqref{Eqn: True Dynamics, Smoothed}.

\begin{proposition}
(\textbf{Smoothed PGD Dynamics are Lipschitz and admit Lipschitz Gradients})
\label{Prop: Properties of Smoothed Full Dynamics}
There exists $\lipschitz_\dynamicsFullStateSmoothed, \hessianBoundSmoothed_{i, \samplingRadius}$ such that, $\forall \ (\state, \ratePGD), (\state', \ratePGD') \in \varConvexHull \times \paramSet \times [\ratePGDMin, \ratePGDMax]$ and $i \in [\varDim + \paramDim]$:
{
\setlength{\abovedisplayskip}{3.5pt}
\setlength{\belowdisplayskip}{3.5pt}
\begin{align}
\label{Eqn: Smoothed Dynamics vs. Dynamics, Error Bound}
    \Vert \dynamicsFullStateSmoothed(\state, \ratePGD; \samplingRadius) - \dynamicsFullState(\state, \ratePGD) \Vert_\infty &\leq \lipschitz_\dynamicsFullStateSmoothed \samplingRadius, \\
\label{Eqn: Smoothed Dynamics Gradient, Lipschitz Constant}
    \Vert \nabla \dynamicsFullStateSmoothed_i(\state, \ratePGD; \samplingRadius) - \nabla \dynamicsFullStateSmoothed_i(\state', \ratePGD'; \samplingRadius) \Vert_1 &\leq \hessianBoundSmoothed_{i, \samplingRadius} \Vert (\state, \ratePGD) - (\state', \ratePGD') \Vert_\infty.
\end{align}
}
\end{proposition}

\begin{proof}
\looseness-1Since $\dynamicsFullState$ is continuous, it is $\lipschitz_{\dynamicsFullStateSmoothed}$-Lipschitz continuous over the bounded set $\big((\varConvexHull \times \paramSet) \oplus \samplingRadiusMax I_{\varDim + \paramDim} \ballUnit_2^{\varDim + \paramDim} \big) \times [\ratePGDMin, \ratePGDMax]$. Thus, \eqref{Eqn: Smoothed Dynamics vs. Dynamics, Error Bound}-\eqref{Eqn: Smoothed Dynamics Gradient, Lipschitz Constant} follow directly from Prop. \ref{Prop: Properties of Smoothed Functions}.
\end{proof}

% \begin{remark}
% Note that the $\param$-update in $\dynamicsFullState$ is not smoothed when constructing $\dynamicsFullStateSmoothed$ (see \eqref{Eqn: True Dynamics, Smoothed}). Thus, when computing the Lipschitz constant $\lipschitz_\dynamicsFullStateSmoothed$ and $\hessianBoundSmoothed_\dynamicsFullStateSmoothed$ in the context of Prop. \ref{Prop: Properties of Smoothed Full Dynamics}, 
% we do not consider the contribution of the $\theta$-update in the definitions of $\dynamicsFullState$ and $\dynamicsFullStateSmoothed.$
% % one only needs to consider variations in $\dynamicsFullStateSmoothed$ in response to variations in $\var$ and $\ratePGD$, but not $\param$.
% \end{remark}

Equipped with
% the properties of $\dynamicsFullStateSmoothed$ provided by
Prop. \ref{Prop: Properties of Smoothed Full Dynamics}, we proceed to adapt the methods in Sec. \ref{subsec: LTV Error System for Smooth PGD Dynamics}-\ref{subsec: Linearization Error Bounds} to settings in which $\dynamicsFullState$ is non-differentiable.
First, we use $\dynamicsFullStateSmoothed$, rather than $\dynamicsFullState$, to re-define the gain and actuation matrices of~\eqref{Eqn: Jacobian Matrices} for the LTV error system as:
% Concretely, 
% Instead of \eqref{Eqn: Jacobian Matrices}, we compute:
% \eqref{Eqn: Jacobian Matrices, with Smoothed Dynamics}, as shown below:
{
\setlength{\abovedisplayskip}{3.5pt}
\setlength{\belowdisplayskip}{3.5pt}
\begin{subequations}
\label{Eqn: Jacobian Matrices, with Smoothed Dynamics}
\begin{align}
\label{Eqn: Jacobian Matrix, State, with Smoothed Dynamics}
    \jacobianStateSmoothed_k(\stateNominal_k, \ratePGDNominal_k) &:= \textstyle\frac{\partial \dynamicsFullStateSmoothed}{\partial \state}(\stateNominal_k, \ratePGDNominal_k), \\
    \label{Eqn: Jacobian Matrix, Control, with Smoothed Dynamics}
    \jacobianControlSmoothed_k(\stateNominal_k, \ratePGDNominal_k) &:= \textstyle\frac{\partial \dynamicsFullStateSmoothed}{\partial \ratePGD}(\stateNominal_k, \ratePGDNominal_k).
\end{align}
\end{subequations}
}
\hspace{-1mm}Similarly, we modify the definition of the linearization error $\linearizationErrorSmoothed_k(\state_k, \ratePGD_k, \stateNominal_k, \ratePGDNominal_k) \in \R^{\varDim + \paramDim}$, as given by \eqref{Eqn: Linearization Error}, into:
{
\setlength{\abovedisplayskip}{3.5pt}
\setlength{\belowdisplayskip}{3.5pt}
\begin{align}
\label{Eqn: Linearization Error, with Smoothed Dynamics}
    \linearizationErrorSmoothed_k(\state_k, \ratePGD_k, \stateNominal_k, \ratePGDNominal_k)  & = \dynamicsFullStateSmoothed(\state_k, \ratePGD_k) - \dynamicsFullStateSmoothed(\stateNominal_k, \ratePGDNominal_k) \\ \nonumber 
    &\hspace{5mm} - \jacobianStateSmoothed_k(\stateNominal_k, \ratePGDNominal_k) \stateError_k - \jacobianControlSmoothed_k(\stateNominal_k, \ratePGDNominal_k) \ratePGDError_k,
\end{align}
}
\hspace{-1mm}We also modify the disturbance definition $\disturbance_k$ in \eqref{Eqn: Disturbance}
into $\disturbanceSmoothed_k$,
defined below,
% in \eqref{Eqn: Disturbance, with Smoothed Dynamics},
to additionally account for the difference between the true dynamics $\dynamicsFullState$ and its smoothed proxy $\dynamicsFullStateSmoothed$:
% , as well as the linearization error:
{
\setlength{\abovedisplayskip}{3.5pt}
\setlength{\belowdisplayskip}{3.5pt}
\begin{align} 
\label{Eqn: Disturbance, with Smoothed Dynamics}
    \disturbanceSmoothed_k &= \linearizationErrorSmoothed_k(\state_k, \ratePGD_k, \stateNominal_k, \ratePGDNominal_k) + \big[ \dynamicsFullState(\state_k, \ratePGD_k) - \dynamicsFullStateSmoothed(\state_k, \ratePGD_k) \big] \\ \nonumber
    &\hspace{1cm} + \big[ \dynamicsFullStateSmoothed(\stateNominal_k, \ratePGDNominal_k) - \dynamicsFullState(\stateNominal_k, \ratePGDNominal_k) \big].
\end{align}
}
\hspace{-1mm}Using the refined  linearized dynamics and linearization errors \eqref{Eqn: Jacobian Matrices, with Smoothed Dynamics}-\eqref{Eqn: Disturbance, with Smoothed Dynamics},
% \antoine{I think it's misleading that we overload the definition of $d_k$. Somehow the definition \eqref{Eqn: Disturbance, with Smoothed Dynamics} is not consistent with \eqref{Eqn: Disturbance Bound, Stacked, with Smoothed Dynamics}}
% \frank{I've added \say{tildes} to all quantities as necessary, per our discussion.}
% , \eqref{Eqn: Linearization Error, with Smoothed Dynamics}, and 
we derive an expression analogous to \eqref{Eqn: State Error LTV Dynamics} for the state error $\stateError_k := \state_k - \stateNominal_k$, at each iteration $k$:
% \begin{align} \nonumber
%     \stateError_{k+1} &:= \state_{k+1} - \stateNominal_{k+1} \\ \nonumber
%     &= \dynamicsFullState(\state_k, \ratePGD_k) - \dynamicsFullState(\stateNominal_k, \ratePGDNominal_k) \\ 
%     \nonumber
%     % \label{Eqn: State Error LTV Dynamics, with Smoothed Dynamics}
%     &= \jacobianStateSmoothed_k(\stateNominal_k, \ratePGDNominal_k) \stateError_k + \jacobianControlSmoothed_k(\stateNominal_k, \ratePGDNominal_k) \ratePGDError_k + \disturbanceSmoothed_k.
% \end{align}
{
\setlength{\abovedisplayskip}{3.5pt}
\setlength{\belowdisplayskip}{3.5pt}
\begin{align} \nonumber
    \stateError_{k+1} 
    % &:= \state_{k+1} - \stateNominal_{k+1}
    &= \jacobianStateSmoothed_k(\stateNominal_k, \ratePGDNominal_k) \stateError_k + \jacobianControlSmoothed_k(\stateNominal_k, \ratePGDNominal_k) \ratePGDError_k + \disturbanceSmoothed_k.
\end{align}
}

\looseness-1We now adapt Prop. \ref{Prop: Lipschitz Continuity of Gradient of PGD Dynamics}, Prop. \ref{Prop: Quadratic Bound on Linearization Error}, and Thm. \ref{Thm: Robust Steplength Constraint Satisfaction} to the setting in which $\dynamicsFullState$ is non-differentiable.
For Prop. \ref{Prop: Lipschitz Continuity of Gradient of PGD Dynamics}, we replace \eqref{Eqn: Lipschitz Continuity of Gradient of PGD Dynamics} with \eqref{Eqn: Smoothed Dynamics Gradient, Lipschitz Constant}.
For Prop. \ref{Prop: Quadratic Bound on Linearization Error}, we replace $\hessianBound_i$ with $\hessianBoundSmoothed_i$ in \eqref{Eqn: Quadratic Bound on Linearization Error}, i.e.,:
% \antoine{not sure, but the references aboe seem wrong. It shoudl Prop 4 and 5 no?}
% \frank{Thanks. Fixed.}
{
\setlength{\abovedisplayskip}{3.5pt}
\setlength{\belowdisplayskip}{3.5pt}
\begin{align}
\label{Eqn: Quadratic Bound on Linearization Error, with Smoothed Dynamics}
    |\linearizationErrorSmoothed_{k,i}(\state_k, \ratePGD_k, \stateNominal_k, \ratePGDNominal_k)| &\leq \hessianBoundSmoothed_i \cdot \Vert (\stateError_k, \ratePGDError_k) \Vert_\infty^2,
\end{align}
}
which can likewise be proved via \eqref{Eqn: Smoothed Dynamics Gradient, Lipschitz Constant} from Prop. \ref{Prop: Lipschitz Continuity of Gradient of PGD Dynamics}.

Next, we modify the disturbance bounds \eqref{Eqn: Disturbance Bound, Stacked} to incorporate the \textit{smoothing error}, i.e., the difference between $\dynamicsFullState$ and $\dynamicsFullStateSmoothed$ as evaluated on the true and nominal trajectories, in addition to the linearization error $\linearizationError_k(\cdot)$:
{
\setlength{\abovedisplayskip}{3.5pt}
\setlength{\belowdisplayskip}{3.5pt}
\begin{align} \nonumber
    \disturbanceSmoothed_k &= \linearizationErrorSmoothed_k(\state_k, \ratePGD_k, \stateNominal_k, \ratePGDNominal_k) + \big[ \dynamicsFullState(\state_k, \ratePGD_k) - \dynamicsFullStateSmoothed(\state_k, \ratePGD_k) \big] \\ \nonumber
    &\hspace{1cm} + \big[ \dynamicsFullStateSmoothed(\stateNominal_k, \ratePGDNominal_k) - \dynamicsFullState(\stateNominal_k, \ratePGDNominal_k) \big] \\ \label{Eqn: Disturbance Bound, Stacked, with Smoothed Dynamics}
    &\in \big( 2 \lipschitz_\dynamicsFullStateSmoothed I_{\varDim + \paramDim} + \Vert (\stateError_k, \ratePGDError_k) \Vert_\infty^2 \hessianBoundSmoothedStacked \big) \ballUnit_\infty^{\varDim + \paramDim},
\end{align}
}
$\hspace{-1mm}$where the value $2 \lipschitz_\dynamicsFullStateSmoothed$ in \eqref{Eqn: Disturbance Bound, Stacked, with Smoothed Dynamics} arises by applying \eqref{Eqn: Smoothed Dynamics vs. Dynamics, Error Bound}, 
and $\hessianBoundSmoothedStacked := \diag\{\hessianBoundSmoothed_1, \cdots, \hessianBoundSmoothed_{\varDim + \paramDim}\}$.

We then modify Thm. \ref{Thm: Robust Steplength Constraint Satisfaction} by respectively replacing \eqref{Eqn: Inequality for Robust Constraint Satisfaction} and \eqref{Eqn: Inequality for Propagating Recursive Tube Bounds} with \eqref{Eqn: Inequality for Robust Constraint Satisfaction, with Smoothed Dynamics} and \eqref{Eqn: Inequality for Propagating Recursive Tube Bounds, with Smoothed Dynamics}, as given below:
{
\setlength{\abovedisplayskip}{3.5pt}
\setlength{\belowdisplayskip}{3.5pt}
\begin{align} 
\label{Eqn: Inequality for Robust Constraint Satisfaction, with Smoothed Dynamics}
    &\ratePGDBoundSlope_p^\top (\stateNominal_k, \ratePGDNominal_k) + \ratePGDBoundOffset_p + \left\Vert \ratePGDBoundSlope_p^\top
    \begin{bmatrix}
        (\systemResponse_{\state \disturbance})_{k, 0} \\
        (\systemResponse_{\ratePGD \disturbance})_{k, 0}
    \end{bmatrix}
    \right\Vert_1 \diam(\paramSet) \\ \nonumber
    &\hspace{5mm} + \sum_{j=1}^k  \left\Vert \ratePGDBoundSlope_p^\top
    \begin{bmatrix}
        (\systemResponse_{\state \disturbance})_{k, 0} \\
        (\systemResponse_{\ratePGD \disturbance})_{k, 0}
    \end{bmatrix}
    \begin{bmatrix}
        2 \lipschitz_\dynamicsFullStateSmoothed I_{\varDim + \paramDim}
        & \tubeSize_{j-1}^2 \hessianBoundSmoothedStacked
    \end{bmatrix}
    \right\Vert_1 \leq 0, \\ \nonumber
    &\hspace{2cm} \forall \ k \in [0, \iterationHorizon], p \in [2], \\ 
\label{Eqn: Inequality for Propagating Recursive Tube Bounds, with Smoothed Dynamics}
    &\hspace{-5pt}\left\Vert 
    \ratePGDBoundSlope_p^\top
    \begin{bmatrix}
        (\systemResponse_{\state \disturbance})_{k, 0} \\
        (\systemResponse_{\ratePGD \disturbance})_{k, 0}
    \end{bmatrix}
    \right\Vert_\infty \diam(\paramSet) \\ \nonumber
    &\hspace{-6pt} +\hspace{-3pt} \sum_{j=1}^k \left\Vert 
    \begin{bmatrix}
        (\systemResponse_{\state \disturbance})_{k, 0} \\
        (\systemResponse_{\ratePGD \disturbance})_{k, 0}
    \end{bmatrix}
    \begin{bmatrix}
        2 \lipschitz_\dynamicsFullStateSmoothed I_{\varDim + \paramDim}
        & \tubeSize_{j-1}^2 \hessianBoundSmoothedStacked
    \end{bmatrix}
    \right\Vert_\infty\hspace{-7pt} \leq \tubeSize_k,\ \ \forall k \in [0, \iterationHorizon].%\\
    % &\hspace{3cm} \forall \ k \in [0, \iterationHorizon].
\end{align}
}
$\hspace{-1mm}$The validity of \eqref{Eqn: Inequality for Robust Constraint Satisfaction, with Smoothed Dynamics} and \eqref{Eqn: Inequality for Propagating Recursive Tube Bounds, with Smoothed Dynamics} as replacements for \eqref{Eqn: Inequality for Robust Constraint Satisfaction} and \eqref{Eqn: Inequality for Propagating Recursive Tube Bounds} in Thm. \ref{Thm: Robust Steplength Constraint Satisfaction}, for non-differentiable dynamics $\dynamicsFullState$, can be verified by retracing the proof of Thm. \ref{Thm: Robust Steplength Constraint Satisfaction}, and applying the modified disturbance bound \eqref{Eqn: Disturbance Bound, Stacked, with Smoothed Dynamics} instead of its original form \eqref{Eqn: Disturbance Bound, Stacked} when appropriate. 
Similarly, we refine the NLP \eqref{Eqn: NLP for Tightest Tube Size}, used to recover the tightest end-of-horizon tube size $\tubeSize_\iterationHorizon$, by replacing \eqref{Eqn: Inequality for Robust Constraint Satisfaction} and \eqref{Eqn: Inequality for Propagating Recursive Tube Bounds} with \eqref{Eqn: Inequality for Robust Constraint Satisfaction, with Smoothed Dynamics} and \eqref{Eqn: Inequality for Propagating Recursive Tube Bounds, with Smoothed Dynamics}, respectively.
% ~\\
% \antoine{Maybe we add the usual time complexity here. As we work on augmented state, we would need $\mathcal{O}(N^3 n^3)$}
% \frank{Just checking: Are \say{$N$} and \say{$n$} the time horizon and state dimension, respectively, in this context?
% }
% \glen{should make some discussion on how one solves the resulting SLS problem (refer to the Fast OF SLS paper). can also refer to the GPU paper (which has the computational complexities also written out, for the GPU case)}


\section{Simulations}
\label{sec: Experiments}
% \antoine{Something I missed: Shuyu's solver can only deal with delayed measurement. How did we use it here?}
Here, we apply Thms. \ref{Thm: Over-approximating Minimizer Set using Forward Reachable Sets}-\ref{Thm: Robust Steplength Constraint Satisfaction} compute over-approximations of the parameterized minimizer sets for specific objectives.
% apply
% % demonstrate an application of 
% Thms. \ref{Thm: Over-approximating Minimizer Set using Forward Reachable Sets}-\ref{Thm: Robust Steplength Constraint Satisfaction} to specific objectives, computing over-approximations of their parameterized minimizer sets $\varOptSet$. 
% We 
% apply
% % demonstrate an application of 
% Thms. \ref{Thm: Over-approximating Minimizer Set using Forward Reachable Sets}-\ref{Thm: Robust Steplength Constraint Satisfaction} to specific objectives, computing over-approximations of their parameterized minimizer sets $\varOptSet$. 
% \antoine{Use either ``over-approximations'' or ``over-approximations''}
% \frank{Fixed (I'm uniformly using the latter, dashed version).}
% In~\ref{subsec: ex_quadratic_optim}, we consider a simple scalar, quadratic, unconstrained optimization problem for ease of exposition. 
% Next, in~\ref{subsec: ex_lqr}, we present our main example: bounding the planned trajectories that result as solutions to a parameterized linear-quadratic regulator (LQR) problem. 

% \antoine{Here or earlier: The algo performance depends on the initial guess $\xi_0$. How is it picked (I would say ``CE''), and does it influence the performance a lot?}

% \antoine{Another potential limitation, is that we can't deal with large set $\Theta$, because it could make the over-approximation too loose. Also, the horizon length influences the conservatism.}

% \antoine{I would make one plot where we can the step size fixed and see how much tighter the tube become.}

% \antoine{Note to self: read this document with KKT conditions mental model in mind}

\subsection{Quadratic Optimization}
\label{subsec: ex_quadratic_optim}

% \brendan{TODO. Should this example go here, or earlier so that we can use it as a tool in discussions? }

% \looseness=-1
Consider the unconstrained minimization of a scalar cost:
{
\setlength{\abovedisplayskip}{3.5pt}
\setlength{\belowdisplayskip}{3.5pt}
\begin{equation}
    \label{eq:scalar_quadratic_cost}
    \min_{\var \in \R} \objSq(\var, \param) := \var^2 + \param \var
\end{equation}
}
\hspace{-1.5mm}over $\var \in \R$,
% with no constraints (i.e., $\constraintSet = \R$) 
% and
with parameter uncertainty over $\param \in \paramSetSq := [-0.1, 0.1]$.
$\objSq$ is $m$-strongly convex with $L$-Lipschitz gradients, where
% constant 2 in each case, 
$m = L = 2$.
% , thus
% satisfying Assumption~\ref{Assump: Objective Properties}.
% Since minimization is unconstrained, the PGD dynamics are given simply by 
The PGD dynamics are:
\begin{equation}
    \label{eq:scalar_quadratic_pgd}
    \dynamicsPGDSq(\var, \param, \ratePGD) := \var - \ratePGD (2 \var + \param).
\end{equation}
% It is simple to 
We 
compute $\nabla^2 \dynamicsFullState_i$ for $i \in \{1, 2\}$ to find that the curvature of $\dynamicsFullState$ can be bounded by $\hessianBound_1=2$, $\hessianBound_2=0$, as in Prop. \ref{Prop: Lipschitz Continuity of Gradient of PGD Dynamics}.

\looseness-1Choosing $\var_0 = 1.0$, $\ratePGDMin=0.4$, $\ratePGDMax=0.6$ and applying Thm. \ref{Thm: Robust Steplength Constraint Satisfaction}, we bound in 
% $\lVert \cdot \rVert_\infty$ 
$\infty$-norm
% the location of the 
\emph{iterates} of PGD dynamics with horizon $\iterationHorizon=20$. 
By~\eqref{eq:PGD_convergence_rate} and Prop. \ref{Prop: Convergence of PGD}, these iterates converge exponentially to $\varOptSet$ with rate $\convergenceRatePGD = 0.2$.
As our SLS bounds ensure that the iterates are bounded between $\overline{\var} = 10$ and $\underline{\var} = - \overline{\var}$, it can be shown that $\lVert \var_0 - \varOpt(\param) \rVert \le \frac{\lVert \overline{\var} - \underline{\var} \rVert}{1 - \convergenceRatePGD^\iterationHorizon}$ for all $\param \in \paramSet$, overapproximating the distance between the iterates and solution ($\convergenceRadius_k(\var_0)$ in Thm. \ref{Thm: Over-approximating Minimizer Set using Forward Reachable Sets}).
These bounds are combined in Fig.~\ref{fig:unconstrained_quadratic} to overapproximate the parameterized minimizer set of~\eqref{eq:scalar_quadratic_cost}.

\begin{figure}
    \centering
    \includegraphics[width=\linewidth]{Figures/bilinear_1d.pdf}
    \caption{%
        Robust over-approximation of the parameterized solution set of~\eqref{eq:scalar_quadratic_cost}.
        Individual rollouts of PGD corresponding to realizations of $\param$ Monte Carlo sampled from $\paramSetSq$ are shown in black, and the ``nominal'' rollout with $\param = 0$ is shown in red. 
        The blue bounds show the reachable-set bounds computed via Thm. \ref{Thm: Robust Steplength Constraint Satisfaction} that must contain the iterate trajectory for any parameter in $\paramSetSq$. The outer green bounds are guaranteed to contain $\varOptSet$ through Thm. \ref{Thm: Over-approximating Minimizer Set using Forward Reachable Sets}.%
    }
    \label{fig:unconstrained_quadratic}
\end{figure}

\looseness-1
% In this problem, the 
% set of parameterized minimizers 
The
minimizer set
for \eqref{eq:scalar_quadratic_cost}
can be explicitly expressed as $\varOptSet = \{ - \frac{\param}{2} \mid \param \in \paramSet \}$.
While our SLS-computed bounds are conservative relative to
$\varOptSet$,
% this exact solution,
our method can provide tighter bounds than classic robust optimization methods when
% a closed-form expression for 
% $\varOptSet$ is intractable to obtain 
no closed-form expression for 
$\varOptSet$ is available
(see Sec. \ref{subsec: ex_lqr}).
% an explicit characterization of $\varOptSet$ is intractable to obtain.
% However, as we show below, our method is applicable in cases where explicit characterization of $\varOptSet$ is not tractable, and can provide tighter bounds than classic robust optimization techniques. 

\subsection{LQR with Intent Uncertainty}
\label{subsec: ex_lqr}

Consider the following parameterized LQR problem, with variable $\var := (x_0, u_0, \ldots, x_{\LQRHorizon-1}, u_{\LQRHorizon-1}, x_\LQRHorizon) \in \R^{4(T+1) + 2T}$, objective $\objLQR(\var, \param)$, and parameter $\param \in \R$:
% , as formulated below:
{
\setlength{\abovedisplayskip}{3.5pt}
\setlength{\belowdisplayskip}{3.5pt}
\begin{subequations}
\label{eq:lqr}
\begin{align}
\label{eq:lqr_cost}
    \min_{\var} \hspace{3mm} &\frac{1}{2} \textstyle x_\LQRHorizon^\top \LQRStateCost_\LQRHorizon x_\LQRHorizon + \frac{1}{2} \sum_{\LQRIterate=0}^{\LQRHorizon} x_\LQRIterate^\top \LQRStateCost_\LQRIterate x_\LQRIterate + \param u_\LQRIterate^\top \LQRControlCost_\LQRIterate u_\LQRIterate, \\
\label{eq:lqr_dyn}
    &x_{\LQRIterate+1} = \LQRStateDyn_{\LQRIterate} x_{\LQRIterate} + \LQRControlDyn_{\LQRIterate} u_{\LQRIterate}, \hspace{5mm} x_0 = \bar{x}_0,
\end{align}
\end{subequations}
}
% , which seeks to minimize a quadratic cost 
% \begin{equation}
%     \label{eq:lqr_cost}
%     \textstyle\objLQR(\var, \param) = \frac{1}{2} x_\LQRHorizon^\top \LQRStateCost_\LQRHorizon x_\LQRHorizon + \frac{1}{2} \sum_{\LQRIterate=0}^{\LQRHorizon} x_\LQRIterate^\top \LQRStateCost_\LQRIterate x_\LQRIterate + u_\LQRIterate^\top \param \LQRControlCost_\LQRIterate u_\LQRIterate
% \end{equation}
% subject to the dynamics constraints
% \begin{equation}
%     \label{eq:lqr_dyn}
%     x_{\LQRIterate+1} = \LQRStateDyn_{\LQRIterate} x_{\LQRIterate} + \LQRControlDyn_{\LQRIterate} u_{\LQRIterate}, \quad x_0 = \bar{x}_0
% \end{equation}
% \antoine{I guess that $\theta \in \paramSetLQR$?}
% We stack the \emph{entire state-control trajectory} into a decision variable $\var = (x_0, u_0, \ldots, x_{\LQRHorizon-1}, u_{\LQRHorizon-1}, x_\LQRHorizon)$ and minimize~\eqref{eq:lqr_cost} using PGD.
% Specifically, consider 2-dimensional double integrator dynamics with equally weighted costs in all components: 
\hspace{-1mm}where:
\begin{equation}
    \LQRStateDyn_\LQRIterate = \begin{bsmallmatrix}
        1 & 1 & 0 & 0 \\
        0 & 1 & 0 & 0 \\
        0 & 0 & 1 & 1 \\
        0 & 0 & 0 & 1
    \end{bsmallmatrix}, 
    \LQRControlDyn_\LQRIterate = \begin{bsmallmatrix}
        0 & 0 \\
        1 & 0 \\
        0 & 0 \\
        0 & 1 
    \end{bsmallmatrix}, 
    \LQRStateCost_\LQRIterate = 0.1 I,
    \LQRControlCost_\LQRIterate = 0.1 I.
\end{equation}
In words, $\LQRStateDyn_\LQRIterate$ and $\LQRControlDyn_\LQRIterate$ describe time-invariant double integrator dynamics.
Note that the combined state $\var$ has dimension $4(\LQRHorizon + 1) + 2\LQRHorizon = 64$ for $T=10$.
Above, the parameter uncertainty over $\param \in \paramSetLQR := [0.9, 1.1]$ influences the relative weight of state cost and control cost to 
% the ego
a path planning agent.
% \antoine{Suggestion:  The ``-- 64'' is a bit confusing, especially right after a mathematical expression, because it looks like a minus.}
% Defining the block matrices 
% \antoine{I understand, but I think it's a lot of symbol introduction here. I would suggest something like this ``By stacking the full trajectory the finite-horizon LQR becomes a strongly convex quadratic objective subject to linear equality constraints, equivalently optimization over an affine subspace.''}
% \frank{I agree with (what I think is) Antoine's main point here, i.e., that the description of $\objLQR$ and the associated constraint set can be substantially streamlined once $A, B, Q, R$, $H_Q$, $H_R$ have been defined. I.e., it may not be necessary to write $M$ and $b$ explicitly.}
We compactly write the cost~\eqref{eq:lqr_cost} as $\objLQR(\var, \param) = \frac{1}{2} \var^\top H \var = \frac{1}{2} \var^\top (H_\LQRStateCost + \param H_\LQRControlCost) \var$ and dynamics~\eqref{eq:lqr_dyn} as $M \var = b$ for appropriate $H$, $M$, and $b$. 
Concretely,
{
\setlength{\abovedisplayskip}{3.5pt}
\setlength{\belowdisplayskip}{3.5pt}
\begin{align*}
    H_\LQRStateCost &= \blockDiag(\LQRStateCost_0, O, \ldots, \LQRStateCost_{\LQRHorizon-1}, O, \LQRStateCost_{\LQRHorizon}), \\
    H_\LQRControlCost &= \blockDiag(O, \LQRControlCost_0, \ldots, O, \LQRControlCost_{\LQRHorizon-1}, \LQRStateCost_{\LQRHorizon}), % \\
    % \begin{bsmallmatrix}
    %     \LQRStateCost_{0} & 0 & 0 & 0 & \cdots & 0 \\
    %     0 & \param \LQRControlCost_{0} & 0 & 0 & \cdots & 0 \\
    %     0 & 0 & \LQRStateCost_{1} & 0 & \cdots & 0 \\
    %     0 & 0 & 0 & \param \LQRControlCost_{1} & \cdots & 0 \\
    %     \vdots & \vdots & \vdots & \vdots & \ddots & \vdots \\ 
    %     0 & 0 & 0 & 0 & \cdots & \LQRStateCost_{\LQRHorizon}
    % \end{bsmallmatrix}, 
    % M &= \begin{bsmallmatrix}
    %     I & O & O & O & \cdots & O & O & O \\
    %     - \LQRStateDyn_{0} & - \LQRControlDyn_{O} & I & O & \cdots & O & O & O \\
    %     O & O & - \LQRStateDyn & - \LQRControlDyn & \cdots & O & O & O \\
    %     \vdots & \vdots & \vdots & \vdots & \ddots & \vdots & \vdots & \vdots \\
    %     O & O & O & O & \cdots & - \LQRStateDyn & - \LQRControlDyn & I 
    % \end{bsmallmatrix}, 
    % b = \begin{bsmallmatrix} \bar{x}_0 \\ \zero \\ \zero \\ \vdots \\ \zero \end{bsmallmatrix},
\end{align*}
}
% \antoine{I am not sure if I understand the construction of the 2 matrices above. $Q_{T-1}$ come up two times inside $H_Q$ and $H_R$ somehow also has a $Q_T-1$ ? Is it intended? Also, according to the notation section, a matrix of zeros should be $O$, i.e., the letter, and not the digit $\zero$.}
% \brendan{Good catches. The duplicated $T-1$ was a typo, and I fixed the notation to be consistent with the rest of the paper}
The constraint manifold is the affine space given by $\{\var = Pz + d: z \in \R^\varDim \}$, with $P := I - M^\top (M M^\top)^{-1}M$ and $d = M^\top (MM^\top)^{-1}b$, 
%\antoine{What's the affine space? I assume something like $Pz + d, z \in \mathbb{R}^n$ ?}
%\brendan{That is exactly correct. I'm not sure how to clarify this, is it easier in the context of the below equation? I could add an explicit reference}
% \frank{I think the affine space is just $\{\var: M \var = b\}$.}
yielding PGD dynamics 
{
\setlength{\abovedisplayskip}{3.5pt}
\setlength{\belowdisplayskip}{3.5pt}
\begin{equation}
    \label{eq:lqr_pgd}
    \dynamicsPGDLQR(\var, \param, \ratePGD) = P(I - \ratePGD H) \var + d.
\end{equation}
}
\hspace{-3mm}$\objLQR$ is $\min_{\param \in \paramSetLQR} \lambda_{\rm min}(H) = 0.09$-strongly convex and has gradients that are Lipschitz with constant $\max_{\param \in \paramSetLQR} \lambda_{\rm max}(H) = 0.11$.
% , satisfying Assumption~\ref{Assump: Objective Properties}.
% Due to the high-dimensional nature of the
% % problem of the
% LQR problem, 
Since the LQR problem is high-dimensional,
% \antoine{Suggestion ``the LQR problem'' instead of ``the problem of the LQR problem'' ?}
the 
% exact 
computation of the curvature bounds $\hessianBoundStacked$ is 
% significantly 
much
more challenging than was the case for the
% quadratic optimization
example in Sec. \ref{subsec: ex_quadratic_optim}.
% Due to the size of this problem, it is not simple to exactly compute the curvature bounds $\hessianBoundStacked$ as it was in the previous example. 
Fixing the region $\overline{\var} = (20, 20, \ldots, 20)$, $\underline{\var} = - \overline{\var}$, $\ratePGDMin=9.9$, $\ratePGDMax=10.1$ (guaranteed by SLS to contain the true iterates), we 
% automatically 
compute a guaranteed local over-approximation using the interval arithmetic toolbox {\tt immrax}~\cite{immrax}.
Then,
as detailed
% using the pipeline detailed 
in Sec. \ref{subsec: ex_quadratic_optim}, 
% with $\iterationHorizon=10$ PGD iterations, 
we compute bounds on the iterates $\xi_k$ of PGD minimizing~\eqref{eq:lqr_cost}
across $\iterationHorizon = 10$ iterations, and bloat these bounds using the convergence rate $\convergenceRatePGD = 0.111$.
In Fig. \ref{fig:lqr_slices}, we visualize components of the resulting anytime over-approximations of
$\varOptSet$
% parameterized optimal $\varOpt$,
% specifically those 
corresponding to position states at LQR timesteps 0, 3, and 6.
% In Fig. \ref{fig:lqr_slices}, we visualize components of the resulting anytime over-approximations of parameterized optimal $\varOpt$, specifically those corresponding to the position states at LQR timesteps 0, 3, and 6.
Our bounds are only slightly larger than 
those indicated by
% the states obtained via
Monte Carlo sampled rollouts. 
% Our bounds are only slightly larger than the empirically achievable states shown in Monte Carlo sampled rollouts. 

\begin{figure}
    \centering
    \includegraphics[width=\linewidth]{Figures/pgd_lqr_slices.pdf}
    \caption{%
        Robust over-approximation of the planned positions of a parameterized LQR system.
        All visual elements are as in Fig. \ref{fig:unconstrained_quadratic}.
        For space, we only show the components of $\var$ corresponding to position on timesteps $\LQRIterate = 0, 3, 6$.%
        }
    \label{fig:lqr_slices}
\end{figure}

\subsubsection{Baseline Comparison}
\label{subsubsec:baseline}
% By contrast \antoine{In contrast to what?},
Compared to our method,
classical Implicit Function Theorem methods~\cite[Section 11.5.2]{WrightRecht2022OptimizationForDataAnalysis}
% robust optimization techniques based on the Implicit Function Theorem 
perform poorly.
% on this problem.
For
% \antoine{new sentence? I don't think we should capitalize ``Letting'' here}
dual variables
$\varDual$,
% the dual variables to $\var$ 
differentiating the KKT conditions of 
% the constrained optimization problem~
\eqref{eq:lqr} at an optimal $(\varOpt, \varDualOpt)$, yields:
% the differential system
{
\setlength{\abovedisplayskip}{3.5pt}
\setlength{\belowdisplayskip}{3.5pt}
\begin{align*}
    \begin{bmatrix}
        H & M^\top \\
        M & O
    \end{bmatrix}
    \begin{bmatrix}
        \nabla_\param \varOpt(\param) \\ 
        \nabla_\param \varDualOpt(\param)
    \end{bmatrix} = 
    \begin{bmatrix}
        - H_\LQRControlCost \var \\ 
        \zero
    \end{bmatrix},
\end{align*}
}
\hspace{-1mm}with solution $\nabla_\param \varOpt = [ H^{-1} M^\top (M H^{-1} M^\top)^{-1} M H^{-1} - H^{-1} ] H_\LQRControlCost \var := \mathcal{P} H_\LQRControlCost \var$.
Over the parameter region $[\underline{\var}, \overline{\var}]$ guaranteed to be robustly forward invariant for PGD by SLS, $L_{\varOpt} := \sup_{\var \in [\underline{\var}, \overline{\var}], \param \in \paramSetLQR} \lVert \nabla_\param \varOpt \rVert_2$ is a local Lipschitz constant for $\varOpt(\param)$. 
By sub-multiplicativity of norms and standard linear algebraic techniques, for any $\param_1, \param_2 \in \paramSetLQR$
{
\setlength{\abovedisplayskip}{3.5pt}
\setlength{\belowdisplayskip}{3.5pt}
\begin{align*}
    \lVert \varOpt(\param_1) - \varOpt(\param_2) &\rVert_2 \le L_\varOpt \lVert \param_1 - \param_2 \rVert_2 \\
    & \le \lVert \mathcal{P} \rVert_2 \lVert H_\LQRControlCost \rVert_2 \diam(\paramSetLQR) \textstyle\sup_{\var \in [\underline{\var}, \overline{\var}]} \lVert \var \rVert_2 \\
    % & \le \lVert H^{-1} \rVert_2 \lVert H_\LQRControlCost \rVert_2 \diam(\paramSetLQR) \sqrt{\sum_{i=0}^{T n_{LQR} + (T-1) m_{LQR}} 20^2}  \\
    & = \lVert H^{-1} \rVert_2 \lVert H_\LQRControlCost \rVert_2 \diam(\paramSetLQR) \lVert \overline{\var} \rVert_2  \approx 35.
    % \\
    % & = \textstyle\frac{1}{0.09} \cdot 0.1 \cdot 0.2 \cdot 20 \sqrt{64}  \approx 35,
\end{align*}
}
% which is $\approx$35
% \antoine{it looks odd to have an equation without left-hand side, no? I would have this $\approx 35$ included in the previous equation},
% \hspace{-1mm}a much looser bound than those in Fig. \ref{fig:lqr_slices}.
\hspace{-1mm}The maximum bound width computed by our method for \emph{any} component of $\var$ is $0.2$, which is much less conservative. 
% Further,
We note that
an analytical expression for $\varOptSet$ is unavailable 
since the LQR dynamics impose multiple equality constraints.
% due to the 
% % presence of
% multiple equality constraints imposed by LQR dynamics.
% These factors demonstrate 
% % the unique ability of
% that
% our method can derive tight bounds on the planned actions of a strategic agent minimizing an uncertain cost.  

\subsection{Constrained Quadratic Optimization}
\label{subsec:ex_constrained_quadratic}

Consider the following constrained variant of~\eqref{eq:scalar_quadratic_cost}:
{
\setlength{\abovedisplayskip}{3pt}
\setlength{\belowdisplayskip}{3pt}
\begin{subequations}
\label{eq:scalar_quadratic_constrained_cost}
\begin{align}
    \min_{\var \in \R} \hspace{5mm} & \objSqCons(\var, \param) := \SqcQuadScale \var^2 + \SqcLinScale \param \var, \\
    \text{s.t.} \hspace{5mm} &\var \ge \SqcConsVal.
\end{align}
\end{subequations}
}

% \begin{equation}
%     \label{eq:scalar_quadratic_constrained_cost}
%     \objSqCons(\var, \param) := \SqcQuadScale \var^2 + \SqcLinScale \param \var, \quad \var \ge \SqcConsVal.
% \end{equation}
% \looseness=-1
Since the constraint 
set above
% defines $\constraintSet$ as a halfspace, which 
is \emph{not} an affine subspace, the PGD dynamics are non-smooth, so we apply the smoothing 
process in Sec.~\ref{subsec: Smoothing PGD Dynamics when Non-Differentiable} to obtain 
% SLS-based 
reachable set bounds.
% , and we must use the smoothing procedures of Sec. \ref{subsec: Smoothing PGD Dynamics when Non-Differentiable} for SLS reachable set bounds to apply. 

The objective satisfies Assumption~\ref{Assump: Objective Properties} with $m = L = 2 \SqcQuadScale$.
The gradient descent dynamics without projection are similar to~\eqref{eq:scalar_quadratic_pgd}, and have Lipschitz constant $\SqcLipschitz = \lvert 1 - 2 \SqcQuadScale \ratePGD \rvert + \lvert \SqcLinScale \ratePGD \rvert + \lvert 2 \SqcQuadScale \var + \SqcLinScale \param \rvert$ 
% \begin{equation}
%     \SqcLipschitz = \lvert 1 - 2 \SqcQuadScale \ratePGD \rvert + \lvert \SqcLinScale \ratePGD \rvert + \lvert 2 \SqcQuadScale \var + \SqcLinScale \param \rvert
% \end{equation}
as in Prop~\ref{Prop: Properties of Smoothed Functions}. 
Since this problem is scalar, projection is non-expansive in every norm, and $\SqcLipschitz$ is also a Lipschitz constant for PGD dynamics.
Applying Prop~\ref{Prop: Properties of Smoothed Functions}, points 3) and 1) with sampling radius $\samplingRadius = 0.1$ bound the curvature of the smoothed dynamics and their distance from the true PGD update. 
Then, \eqref{Eqn: Disturbance Bound, Stacked, with Smoothed Dynamics}--\eqref{Eqn: Inequality for Propagating Recursive Tube Bounds, with Smoothed Dynamics} 
over-approximate
% produce robust over-approximations of the locations of 
PGD iterates for any parameterized realization of~\eqref{eq:scalar_quadratic_constrained_cost}. 

Using $\var_0 = 0.5$, 
% uncertainty set
$\paramSet = [-0.1, 0.1]$, 
% $\paramSet = \paramSetSqc := [-0.1, 0.1]$, 
$\SqcQuadScale = 0.0521$, $\SqcLinScale = 0.0054$, $\ratePGDMin = 9.54$, and $\ratePGDMax = 9.56$,
we outer-approximate
% compute over-approximations of
$\varOptSet$ for~\eqref{eq:scalar_quadratic_constrained_cost} 
% are computed and plotted in 
(see Fig. \ref{fig:constrained_quadratic}). 
Due to error between the true nonsmooth PGD update and the smoothed dynamics used for SLS, each PGD iteration incurs additional conservativeness. 
% Individual rollouts 
Rollouts
quickly approach the constraint boundary, causing the over-approximation to contain non-differentiable points for 
% the majority 
most
of the iteration horizon;
% , and demonstrating the necessity of
we thus require
Prop.~\ref{Prop: Properties of Smoothed Functions} to bound the curvature of the smoothed dynamics. 
Both factors increase conservativeness compared to previous examples. 
% Both these factors contribute to the increased conservatism, compared to previous examples. 

\begin{figure}
    \centering
    \includegraphics[width=\linewidth]{Figures/bilinear_nd.pdf}
    \caption{%
        Robust over-approximation of the parameterized minimizer set of~\eqref{eq:scalar_quadratic_constrained_cost}.
        All visual elements are as in Fig.~\ref{fig:unconstrained_quadratic}. 
        The boundary of the constraint set $\constraintSet$ is also shown in purple, for convenience. 
    }
    \label{fig:constrained_quadratic}
\end{figure}


% \begin{figure}
%     \centering
%     \includesvg[width=\linewidth]{Figures/pgd_lqr_2d_static}
%     \caption{Caption}
%     \label{fig:lqr_pos_space}
% \end{figure}



% \subsection{Constrained Quadratic Optimization}
% \label{subsec: ex_constrained_quadratic_optim}

% \brendan{TODO}

% \INCOMPLETE

\section{Conclusion}
% \vspace{-0.3em}
\label{sec: Conclusion and Future Work}

We present a novel system-level synthesis (SLS)-based method for over-approximating the minimizer set of strongly convex objective costs characterized by uncertain parameters.
Our method leverages the convergence properties of projected gradient descent (PGD) dynamics over convex optimization landscapes, to recast the problem of over-approximating the minimizer set as that of computing the PGD dynamics' forward reachable sets.
% Our method leverages the convergence properties of projected gradient descent (PGD) dynamics over convex optimization landscapes.
% Leveraging the convergence properties of projected gradient descent (PGD) dynamics when applied to convex programs, we recast the problem of over-approximating the minimizer set as that of computing the PGD dynamics' forward reachable sets.
By considering the cost parameter as constant but unknown, we model PGD dynamics as a nonlinear system with initial state uncertainty, and refine existing nonlinear SLS techniques to compute its forward reachable sets.
% to over-approximate the set of parameterized minimizers.
Our numerical results illustrate that our method outperforms sensitivity analysis and fixed-steplength PGD baselines in over-approximating the minimizer set with low conservativeness.
Future work will fuse our method for minimizer set over-approximation with adaptive control techniques for shrinking $\Theta$ online from observations. 

% \antoine{Future work: Apply adaptive control theory to reduce the set $\Theta$ online. Set membership estimation, etc should apply directly.}
% ~\\
% \frank{INCOMPLETE}
% ~\\


% \INCOMPLETE



% \clearpage
% \input{temp_figures}

\bibliographystyle{IEEEtran}
% \bibliography{references}
\bibliography{references_trunc}
% \printbibliography

% \appendix


\section*{Acknowledgement}
The authors thank Shuyu Zhan for his assistance with the simulation experiments.

% \addtolength{\textheight}{-12cm}   
% This command serves to balance the column lengths on the last page of the document manually. It shortens the textheight of the last page by a suitable amount. This command does not take effect until the next page so it should come on the page before the last. Make sure that you do not shorten the textheight too much.






%%%%%%%%%%%%%%%%%%%%%%%%%%%%%%%%%%%%%%%%%%%%%%%%%%%%%%%%%%%%%%%%%%%%%%%%%%%%%%%%







% \begin{thebibliography}{99}

% \bibitem{c1} G. O. Young, ``Synthetic structure of industrial plastics (Book style with paper title and editor),'' 	in Plastics, 2nd ed. vol. 3, J. Peters, Ed.  New York: McGraw-Hill, 1964, pp. 15--64.
% \bibitem{c2} W.-K. Chen, Linear Networks and Systems (Book style).	Belmont, CA: Wadsworth, 1993, pp. 123--135.
% \bibitem{c3} H. Poor, An Introduction to Signal Detection and Estimation.   New York: Springer-Verlag, 1985, ch. 4.

% \end{thebibliography}

\end{document}
